% Generated mechanically from frozen input p12/ko/build/ch115/obsolete_full.tex.
% Do not edit this generated file; edit the bound locale source instead.

% OK, start here.
%

\setcounter{chapter}{114}
\chapter{폐기된 결과}


\phantomsection
\label{obsolete-section-phantom}


\section{서론}
\label{obsolete-section-introduction}

\noindent
이 장에는 이제 ``폐기된'' 몇 가지 보조정리를 모아 둔다
(\cite{Miller}를 보라).


\section{예비 사항}
\label{obsolete-section-preliminaries}

\begin{remark}
\label{obsolete-remark-composition-of-adjoints-isomorphic-to-identity}
이 주석에 들어 있던 내용은 이제
『범주』의 보조정리 \ref{categories-lemma-adjoint-fully-faithful}와
\ref{categories-lemma-left-adjoint-composed-fully-faithful}를 함께 적용하면
모두 얻어진다.
\end{remark}


\section{호몰로지 대수학}
\label{obsolete-section-homological-algebra}

\begin{remark}
\label{obsolete-remark-weak-serre-subcategory}
다음 논의는 『호몰로지』의 보조정리
\ref{homology-lemma-biregular-ss-converges}와
\ref{homology-lemma-first-quadrant-ss}에 포함되므로 폐기되었다.
$\mathcal{A}$를 아벨 범주라 하고,
$\mathcal{C} \subset \mathcal{A}$를 약한 Serre 부분범주라 하자
(『호몰로지』의 정의
\ref{homology-definition-serre-subcategory}을 보라).
『호몰로지』의 보조정리 \ref{homology-lemma-first-quadrant-ss}가 적용되는
이중 복합체 $K^{\bullet, \bullet}$가 주어졌고, 어떤 $r \geq 0$에 대해
모든 대상 ${}'E_r^{p, q}$가 $\mathcal{C}$에 속한다고 하자. 그러면 모든
코호몰로지 군 $H^n(sK^\bullet)$가 $\mathcal{C}$에 속한다. 실제로 가정에
따르면 ${}'d_r^{p, q}$의 핵과 상은 $\mathcal{C}$에 속한다. 따라서 각
${}'E_{r + 1}^{p, q}$가 $\mathcal{C}$에 속한다. 귀납법으로 각
${}'E_\infty^{p, q}$가 $\mathcal{C}$에 속함을 얻는다. 그러므로 각
$H^n(sK^\bullet)$에는 모든 부분몫이 $\mathcal{C}$에 속하는 유한 여과가
있다. $\mathcal{C}$가 확대에 대해 닫혀 있음을 이용하면, 주장한 대로
$H^n(sK^\bullet)$가 $\mathcal{C}$에 속한다. 같은 결과가
$K^{\bullet, \bullet}$에 딸린 두 번째 스펙트럼 열에도 성립한다. 마찬가지로,
여과 복합체 $(K^\bullet, F)$에 『호몰로지』의 보조정리
\ref{homology-lemma-biregular-ss-converges}가 적용되고 어떤 $r \geq 0$에 대해
모든 대상 $E_r^{p, q}$가 $\mathcal{C}$에 속하면, 각 $H^n(K^\bullet)$은
$\mathcal{C}$의 대상이다.
\end{remark}

\section{폐기된 대수학 보조정리}
\label{obsolete-section-algebra}

\begin{lemma}
\label{obsolete-lemma-finite-presentation-module-independent}
$M$을 유한 표시 $R$-가군이라 하자.
임의의 전사 $\alpha : R^{\oplus n} \to M$에 대해
$\alpha$의 핵은 유한 $R$-가군이다.
\end{lemma}

\begin{proof}
이는 『대수학』의 보조정리 \ref{algebra-lemma-extension}의 특수한 경우이다.
\end{proof}

\begin{lemma}
\label{obsolete-lemma-p-ring-map}
$\varphi : R \to S$를 환 준동형이라 하자. 다음이 성립한다고 하자.
\begin{enumerate}
\item 모든 $x \in S$에 대해 $x^n$이 $\varphi$의 상에 속하는
$n > 0$이 존재하고,
\item 모든 $x \in \Ker(\varphi)$에 대해 $x^n = 0$인
$n > 0$이 존재한다.
\end{enumerate}
그러면 $\varphi$는 스펙트럼 사이의 위상동형을 유도한다. 소수 $p$가
주어져서 다음이 성립한다고 하자.
\begin{enumerate}
\item[(a)] $S$는 $R$-대수로서, 어떤 $n > 0$에 대해
$x^{p^n} \in \varphi(R)$ 및 $p^nx \in \varphi(R)$를 만족하는
원소 $x$들로 생성되고,
\item[(b)] $\varphi$의 핵은 멱영원들로 생성된다.
\end{enumerate}
그러면 (1)과 (2)가 성립하며, 임의의 환 준동형 $R \to R'$에 대해
환 준동형 $R' \to R' \otimes_R S$도 (a), (b), (1), (2)를 만족한다.
특히 이 준동형은 스펙트럼 사이의 위상동형을 유도한다.
\end{lemma}

\begin{proof}
이는 『대수학』의 보조정리 \ref{algebra-lemma-powers}와
\ref{algebra-lemma-p-ring-map}을 결합한 것이다.
\end{proof}

\noindent
다음 기술적 보조정리는 적절한 의미에서 본질적으로 유한형인 환 준동형의
한 올에서 주어진 임의의 관계열을 전체 공간으로 올릴 수 있음을 말한다.

\begin{lemma}
\label{obsolete-lemma-lift-elements-ideal}
$R \to S$를 환 준동형이라 하자.
$\mathfrak p \subset R$를 소 아이디얼이라 하자.
$\mathfrak q \subset S$를 $\mathfrak p$ 위에 놓인 소 아이디얼이라 하자.
$S_{\mathfrak q}$가 $R_\mathfrak p$ 위에서 본질적으로 유한형이라고 가정하자.
다음이 주어졌다고 하자.
\begin{enumerate}
\item 정수 $n \geq 0$,
\item 소 아이디얼 $\mathfrak a \subset \kappa(\mathfrak p)[x_1, \ldots, x_n]$,
\item 전사인 $\kappa(\mathfrak p)$-준동형
$$
\psi : (\kappa(\mathfrak p)[x_1, \ldots, x_n])_{\mathfrak a}
\longrightarrow
S_{\mathfrak q}/\mathfrak p S_{\mathfrak q},
$$
\item $\Ker(\psi)$의 원소 $\overline{f}_1, \ldots, \overline{f}_e$.
\end{enumerate}
그러면 다음이 존재한다.
\begin{enumerate}
\item 정수 $m \geq 0$,
\item 원소 $g \in S$, $g \not\in \mathfrak q$,
\item 사상
$$
\Psi :
R[x_1, \ldots, x_n, x_{n + 1}, \ldots, x_{n + m}]
\longrightarrow
S_g,
$$
\item $\Ker(\Psi)$의 원소
$f_1, \ldots, f_e, f_{e + 1}, \ldots, f_{e + m}$.
\end{enumerate}
이들은 다음을 만족한다.
\begin{enumerate}
\item 다음 도표가 가환이고
$$
\xymatrix{
R[x_1, \ldots, x_{n + m}] \ar[d]_\Psi
\ar[rr]_-{x_{n + j} \mapsto 0} & &
(\kappa(\mathfrak p)[x_1, \ldots, x_n])_{\mathfrak a} \ar[d]^\psi \\
S_g \ar[rr] & &
S_{\mathfrak q}/\mathfrak p S_{\mathfrak q}
},
$$
\item 원소 $f_i$는 $i \leq n$일 때 국소환
$$
(\kappa(\mathfrak p)[x_1, \ldots, x_{n + m}])_{
(\mathfrak a, x_{n + 1}, \ldots, x_{n + m})},
$$
에서 $\overline{f}_i$에 단원을 곱한 원소로 사상되고,
\item 원소 $f_{e + j}$는 같은 국소환에서 $x_{n + j}$에 단원을 곱한
원소로 사상되며,
\item 유도된 사상 $R[x_1, \ldots, x_{n + m}]_{\mathfrak b}
\to S_{\mathfrak q}$는 전사이다. 여기서
$\mathfrak b = \Psi^{-1}(\mathfrak qS_g)$이다.
\end{enumerate}
\end{lemma}

\begin{proof}
$R$과 $S$가 각각 극대 아이디얼 $\mathfrak p$와 $\mathfrak q$를 갖는
국소환인 경우에 보조정리를 증명하는 것으로 충분하다고 주장한다.
실제로 필요한 성질을 모두 갖는
$$
\Psi' : R_{\mathfrak p}[x_1, \ldots, x_{n + m}]
\longrightarrow
S_{\mathfrak q}
$$
와 $f_1', \ldots, f_{e + m}' \in R_{\mathfrak p}[x_1, \ldots, x_{n + m}]$
를 구성했다고 하자. 그러면 각 $ff_k'$가 어떤 원소
$f_k \in R[x_1, \ldots, x_{n + m}]$에서 오는 원소가 되게 하는
$f \in R$, $f \not \in \mathfrak p$가 존재한다. 또한 적절한
$g \in S$, $g \not \in \mathfrak q$에 대해 원소 $\Psi'(x_i)$들은
원소 $y_i \in S_g$의 상이다. 규칙 $\Psi(x_i) = y_i$로 정의되는
$R$-대수 준동형을 $\Psi$라 하자. $\Psi(f_i)$는 국소화
$S_{\mathfrak q}$에서 영이므로, 필요하면 $g$를 바꾸어
$\Psi(f_i) = 0$이라고 가정할 수 있다. 이로써 주장이 증명된다.

\medskip\noindent
따라서 $R$과 $S$가 각각 극대 아이디얼 $\mathfrak p$와
$\mathfrak q$를 갖는 국소환이라고 가정해도 된다.
$y_i \bmod \mathfrak pS = \psi(x_i)$가 되도록
$y_1, \ldots, y_n \in S$를 택한다.
$S$가 그 국소화가 되는 어떤 $R$-부분대수를 생성하는 원소
$y_{n + 1}, \ldots, y_{n + m} \in S$를 택한다.
$S$가 $R$ 위에서 본질적으로 유한형이라는 가정에 의해 이런 원소들이
존재한다. $\psi$가 전사이므로 어떤
$h_j \in \kappa(\mathfrak p)[x_1, \ldots, x_n]_{\mathfrak a}$에 대해
$y_{n + j} \bmod \mathfrak pS = \psi(h_j)$라고 쓸 수 있다.
어떤 공통 분모
$d \in \kappa(\mathfrak p)[x_1, \ldots, x_n]$,
$d \not \in \mathfrak a$와
$g_j \in \kappa(\mathfrak p)[x_1, \ldots, x_n]$에 대해
$h_j = g_j/d$라고 쓰자. $g_j$와 $d$의 올림
$G_j, D \in R[x_1, \ldots, x_n]$을 택한다. 다음과 같이 놓는다.
$y_{n + j}' = D(y_1, \ldots, y_n) y_{n + j} - G_j(y_1, \ldots, y_n)$.
구성에 의해 $y_{n + j}' \in \mathfrak p S$이다.
$y_1, \ldots, y_n, y_n', \ldots, y_{n + m}'$가 $S$의 어떤
$R$-부분대수를 생성하고 그 부분대수의 국소화가 $S$임은 명백하다.
다음 사상을 정의한다.
$$
\Psi : R[x_1, \ldots, x_{n + m}] \to S
$$
이는 $i = 1, \ldots, n$에 대해 $x_i$를 $y_i$로 보내고,
$j = 1, \ldots, m$에 대해 $x_{n + j}$를 $y'_{n + j}$로 보낸다.
성질 (1)과 (4)는 구성에서 명백하다. 또한 아이디얼 $\mathfrak b$는
다항식환 $\kappa(\mathfrak p)[x_1, \ldots, x_{n + m}]$의 아이디얼
$(\mathfrak a, x_{n + 1}, \ldots, x_{n + m})$ 위로 사상된다.

\medskip\noindent
$J = \Ker(\Psi)$라 쓰자. 짧은 완전열
$$
0 \to J_{\mathfrak b}
\to R[x_1, \ldots, x_{n + m}]_{\mathfrak b}
\to S_{\mathfrak q}
\to 0.
$$
을 얻는다. 전사성은 위에서 택한
$y_1, \ldots, y_n, y_n', \ldots, y_{n + m}'$에서 따른다.
따라서 다음 열이 완전하다.
$$
J_{\mathfrak b}/ \mathfrak pJ_{\mathfrak b}
\to \kappa(\mathfrak p)[x_1, \ldots, x_{n + m}]_{
(\mathfrak a, x_{n + 1}, \ldots, x_{n + m})}
\to S_{\mathfrak q}/\mathfrak pS_{\mathfrak q}
\to 0
$$
구성에 의해 마지막 사상은 $x_i$를 $\psi(x_i)$로 보내고
$x_{n + j}$를 영으로 보낸다. 따라서 보조정리의 (2)와 (3)에 나오는
$f_i$들을 쉽게 택할 수 있다.
\end{proof}

\begin{lemma}
\label{obsolete-lemma-smooth-independent-presentation}
$R \to S$를 유한 표시 환 준동형이라 하자. 어떤 $R$ 위의 $S$의 표시
$\alpha$에 대해 순진한 여접 복합체 $\NL(\alpha)$가 차수 $0$에 놓인
유한 사영 $S$-가군과 준동형이면, 모든 표시에 대해서도 그렇다.
\end{lemma}

\begin{proof}
『대수학』의 보조정리 \ref{algebra-lemma-NL-homotopy}에서 즉시 따른다.
\end{proof}

\begin{remark}[사영 분해]
\label{obsolete-remark-projective-resolution}
$R$을 환이라 하자. 임의의 집합 $S$에 대해 $F(S)$로 $S$ 위의 자유
$R$-가군을 나타내자. 그러면 모든 왼쪽 $R$-가군에는 다음과 같은
두 단계 분해가 있다.
$$
F(M \times M) \oplus F(R \times M) \to F(M) \to M \to 0.
$$
첫 번째 사상은 다음 규칙으로 주어진다.
$$
[m_1, m_2] \oplus [r, m] \mapsto [m_1 + m_2] - [m_1] - [m_2] + [rm] - r[m].
$$
\end{remark}

\begin{lemma}
\label{obsolete-lemma-spec-localization-first}
$S$를 $A$의 곱셈적 집합이라 하자. 표준 환 준동형
$A \to S^{-1}A$가 유도하는 사상
$$
f: \Spec(S^{-1}A)\longrightarrow \Spec(A)
$$
은 그 상 위의 위상동형이고,
$\Im(f) = \{ \mathfrak p \in \Spec(A) : \mathfrak p\cap S = \emptyset \}$이다.
\end{lemma}

\begin{proof}
이는 『대수학』의 보조정리 \ref{algebra-lemma-spec-localization}과 중복된다.
\end{proof}

\begin{lemma}
\label{obsolete-lemma-finite-type-flat-over-integral-algebra}
$A \to B$를 유한형인 평탄 환 준동형이라 하고 $A$가 정역이라고 하자.
그러면 $B$는 유한 표시 $A$-대수이다.
\end{lemma}

\begin{proof}
『평탄성 심화』의 명제
\ref{flat-proposition-flat-finite-type-finite-presentation-domain}의
특수한 경우이다.
\end{proof}

\begin{lemma}
\label{obsolete-lemma-helper-finite-type-flat-finite-presentation}
$R$을 분수체 $K$를 갖는 정역이라 하자.
$S = R[x_1, \ldots, x_n]$를 $R$ 위의 다항식환이라 하자.
$M$을 유한 $S$-가군이라 하자. $M$이 $R$ 위에서 평탄하다고 가정하자.
모든 부분환 $R \subset R' \subset K$, $R \not = R'$에 대해
$M \otimes_R R'$가 $S \otimes_R R'$ 위의 유한 표시 가군이면,
$M$은 $S$ 위의 유한 표시 가군이다.
\end{lemma}

\begin{proof}
이 보조정리는 참이다. 사실 명제에 나오는 모든 $R'$에 대해
$M \otimes_R R'$가 유한 표시라는 가정 없이도 $M$은 유한 표시이다.
이는 『평탄성 심화』의 명제
\ref{flat-proposition-flat-finite-type-finite-presentation-domain}에서 따른다.
원래 이 보조정리에는 잘못된 증명이 있었고(그 틈을 찾아 준 Ofer Gabber에게
감사한다), 인용한 명제의 다른 증명에 사용되었다. 이 보조정리를 복원하려면
$R$이 국소 정규 정역인 경우에 가환대수학 장들의 결과만을 사용하는 올바른
논증이 필요하다. 그런 논증을 알고 있다면
\href{mailto:stacks.project@gmail.com}{stacks.project@gmail.com}으로
이메일을 보내 주기 바란다.
\end{proof}

\begin{lemma}
\label{obsolete-lemma-relative-effective-cartier-algebra}
$A \to B$를 환 준동형이라 하고 $f \in B$라 하자. 다음을 가정하자.
\begin{enumerate}
\item $A \to B$는 평탄하고,
\item $f$는 영인자가 아니며,
\item $A \to B/fB$는 평탄하다.
\end{enumerate}
그러면 모든 아이디얼 $I \subset A$에 대해 사상
$f : B/IB \to B/IB$는 단사이다.
\end{lemma}

\begin{proof}
$B$와 $B/fB$가 $A$ 위에서 평탄하므로
$IB = I \otimes_A B$이고 $I(B/fB) = I \otimes_A B/fB$이다.
특히 $IB/fIB \cong I \otimes_A B/fB$는 $B/fB$로 단사적으로 사상된다.
따라서 완전한 행을 갖는 도표
$$
\xymatrix{
0 \ar[r] &
I \otimes_A B \ar[r] \ar[d]^f &
B \ar[r] \ar[d]^f &
B/IB \ar[r] \ar[d]^f &
0 \\
0 \ar[r] &
I \otimes_A B \ar[r] &
B \ar[r] &
B/IB \ar[r] &
0
}
$$
에 뱀 보조정리를 적용하면 결과가 따른다.
\end{proof}

\begin{lemma}
\label{obsolete-lemma-faithfully-flat-injective}
$R \to S$가 충실히 평탄한 환 준동형이면, 모든 $R$-가군 $M$에 대해
사상 $M \to S \otimes_R M$, $x \mapsto 1 \otimes x$는 단사이다.
\end{lemma}

\begin{proof}
이 보조정리는 『대수학』의 보조정리
\ref{algebra-lemma-faithfully-flat-universally-injective}와 중복된다.
\end{proof}

\begin{remark}
\label{obsolete-remark-section-colimits}
이 참조/태그는 전에는 『환 준동형의 평활화』 장의 한 절을 가리켰으나,
그 내용은 이후 『대수학』의 절 \ref{algebra-section-colimits-flat}에
포함되었다.
\end{remark}

\begin{lemma}
\label{obsolete-lemma-not-domain}
$(R, \mathfrak m)$을 차원이 $1$인 축소 뇌터 국소환이라 하고
$x \in \mathfrak m$을 영인자가 아닌 원소라 하자.
$\mathfrak q_1, \ldots, \mathfrak q_r$을 $R$의 극소 소 아이디얼들이라 하자.
그러면
$$
\text{length}_R(R/(x)) = \sum\nolimits_i \text{ord}_{R/\mathfrak q_i}(x)
$$
이다.
\end{lemma}

\begin{proof}
『Chow 호몰로지』의 보조정리
\ref{chow-lemma-additivity-divisors-restricted}의 특수한(매우 쉬운) 경우이다.
\end{proof}

\begin{lemma}
\label{obsolete-lemma-bound-primes}
$A$를 차원이 $2$인 뇌터 국소 정규 정역이라 하자.
영이 아닌 $f \in \mathfrak m$에 대해 $A$의 뚫린 스펙트럼 위에서
$f$에 딸린 인자를
$\text{div}(f) = \sum n_i (\mathfrak p_i)$로 나타내자.
$|f| = \sum n_i$로 놓는다. 모든 $g \in \mathfrak m^N$에 대해
$|f + g| \leq M$이 되게 하는 정수 $N$과 $M$이 존재한다.
\end{lemma}

\begin{proof}
$f, h$가 $A$의 정칙열이 되도록 $h \in \mathfrak m$을 택한다
(이는 『대수학』의 보조정리 \ref{algebra-lemma-criterion-normal}과
\ref{algebra-lemma-depth-drops-by-one}에서 따른다).
$M = \text{length}_A(A/(f, h))$로 놓고,
$\mathfrak m^N \subset (f, h)$인 임의의 정수 $N$을 택하면 보조정리가
성립함을 증명하겠다. $\sqrt{(f, h)} = \mathfrak m$이므로 그런 정수
$N$이 존재한다. $(f, h) = (f + g, h)$이므로 모든
$g \in \mathfrak m^N$에 대해
$M = \text{length}_A(A/(f + g, h))$임에 유의하자. 또한 이로부터
$f + g, h$도 $A$의 정칙열임이 따른다. 『대수학』의 보조정리
\ref{algebra-lemma-reformulate-CM}를 보라.
이제 $\text{div}(f + g ) = \sum m_j (\mathfrak q_j)$라고 하자.
다음 사상을 생각한다.
$$
c : A/(f + g) \longrightarrow \prod A/\mathfrak q_j^{(m_j)}
$$
여기서 $\mathfrak q_j^{(m_j)}$는 기호적 거듭제곱이다. 『대수학』의 절
\ref{algebra-section-symbolic-power}를 보라.
$A$가 정규이므로 $A_{\mathfrak q_i}$는 이산 값매김환이고, 따라서
$$
A_{\mathfrak q_i}/(f + g) =
A_{\mathfrak q_i}/\mathfrak q_i^{m_i} A_{\mathfrak q_i} =
(A/\mathfrak q_i^{(m_i)})_{\mathfrak q_i}
$$
이다. $V(f + g, h) = \{\mathfrak m\}$이므로 $h$를 가역화하면 $c$가
동형이 된다(작은 세부 사항은 생략한다). $h$는 $A/(f + g)$ 위의
영인자가 아니므로 $A/(f + g, h)$의 길이는 『Chow 호몰로지』의 절
\ref{chow-section-periodic-complexes}에서 정의한 Herbrand 몫
$e_A(A/(f + g), 0, h)$와 같다.
마찬가지로 $A/(h, \mathfrak q_j^{(m_j)})$의 길이는
$e_A(A/\mathfrak q_j^{(m_j)}, 0, h)$와 같다. 이제
\begin{align*}
M & = \text{length}_A(A/(f + g, h) \\
& =
e_A(A/(f + g), 0, h) \\
& =
\sum\nolimits_i e_A(A/\mathfrak q_j^{(m_j)}, 0, h) \\
& =
\sum\nolimits_i \sum\nolimits_{m = 0, \ldots, m_j - 1}
e_A(\mathfrak q_j^{(m)}/\mathfrak q_j^{(m + 1)}, 0, h)
\end{align*}
를 얻는다. 위에서 논의했듯이 $c$의 여핵이 유한 길이를 갖는다는 사실을
특히 사용하여 『Chow 호몰로지』의 보조정리
\ref{chow-lemma-additivity-periodic-length}와
\ref{chow-lemma-finite-periodic-length}를 적용하면 이 등식들이 따른다.
Nakayama 보조정리를 써서
$e_A(\mathfrak q^{(m)}/\mathfrak q^{(m + 1)}, 0, h)$가 적어도 $1$임을
곧바로 증명할 수 있다. 이로써 보조정리의 증명이 끝난다.
\end{proof}

\begin{lemma}
\label{obsolete-lemma-flat-over-gorenstein-gorenstein-fibre}
$A \to B$를 뇌터 국소환들의 평탄한 국소 준동형이라 하자.
$A$와 $B/\mathfrak m_A B$가 Gorenstein이면 $B$도 Gorenstein이다.
\end{lemma}

\begin{proof}
『쌍대화 복합체』의 보조정리
\ref{dualizing-lemma-flat-under-gorenstein}에서 즉시 따른다.
\end{proof}

\begin{lemma}
\label{obsolete-lemma-kill-local}
$(A, \mathfrak m)$을 뇌터 국소환이라 하자.
$I \subset A$를 아이디얼, $M$을 유한 $A$-가군, $s$를 정수라 하자.
다음을 가정하자.
\begin{enumerate}
\item $A$에는 쌍대화 복합체가 있고,
\item $\mathfrak p \not \in V(I)$이고
$V(\mathfrak p) \cap V(I) \not = \{\mathfrak m\}$이면
$\text{depth}_{A_\mathfrak p}(M_\mathfrak p) + \dim(A/\mathfrak p) > s$이다.
\end{enumerate}
그러면 $V(J) \cap V(I) = \{\mathfrak m\}$이고 $JI^n$이
$i \leq s$에 대해 $H^i_\mathfrak m(M)$을 소멸시키게 하는
$n > 0$과 아이디얼 $J \subset A$가 존재한다.
\end{lemma}

\begin{proof}
『국소 코호몰로지』의 보조정리
\ref{local-cohomology-lemma-sitting-in-degrees}에 따르면,
$i \leq s$인 유한 $A$-가군
$E^i = \text{Ext}^{-i}_A(M, \omega_A^\bullet)$에 대해 이를 보이면 된다.
$E^0 \oplus \ldots \oplus E^s$의 지지 $Z$는 $\Spec(A)$에서 닫혀 있고
(2)에 나오는 어떤 소 아이디얼도 포함하지 않는다. 따라서 보조정리의
명제에 나오는 것과 같은 어떤 $J$에 대해 이는
$V(JI^n)$에 포함된다.
\end{proof}

\begin{lemma}
\label{obsolete-lemma-algebraize-local-cohomology-bis-bis}
$(A, \mathfrak m)$을 뇌터 국소환이라 하자.
$I \subset A$를 아이디얼, $M$을 유한 $A$-가군이라 하자.
$s$와 $d$를 정수라 하자. 다음을 가정하자.
\begin{enumerate}
\item[(a)] $A$에는 쌍대화 복합체가 있고,
\item[(b)] $\text{cd}(A, I) \leq d$이며,
\item[(c)] $\mathfrak p \not \in V(I)$이면
$\text{depth}_{A_\mathfrak p}(M_\mathfrak p) > s$이거나
$\text{depth}_{A_\mathfrak p}(M_\mathfrak p) + \dim(A/\mathfrak p) > d + s$이다.
\end{enumerate}
그러면 $A, I, \mathfrak m, M$에 대해 『대수 및 형식 기하학』의 보조정리
\ref{algebraization-lemma-algebraize-local-cohomology-bis}의 가정들이
성립하고, $i \leq s$에 대해
$H^i_\mathfrak m(M) \to \lim H^i_\mathfrak m(M/I^nM)$은 동형이며
이 가군들은 $I$의 어떤 거듭제곱에 의해 소멸한다.
\end{lemma}

\begin{proof}
『대수 및 형식 기하학』의 보조정리
\ref{algebraization-lemma-algebraize-local-cohomology-bis}의 가정들은
더 일반적인 『대수 및 형식 기하학』의 보조정리
\ref{algebraization-lemma-bootstrap-bis-bis}에 의해 성립한다.
그러면 『대수 및 형식 기하학』의 보조정리
\ref{algebraization-lemma-algebraize-local-cohomology-bis}의 결론이
두 번째 명제를 준다.
\end{proof}

\begin{lemma}
\label{obsolete-lemma-combine-one}
『대수 및 형식 기하학』의 상황 \ref{algebraization-situation-bootstrap}에서
$H^s_\mathfrak a(M) = \lim H^s_\mathfrak a(M/I^nM)$이다.
\end{lemma}

\begin{proof}
이는 『대수 및 형식 기하학』의 정리
\ref{algebraization-theorem-final-bootstrap}에서 즉시 따른다.
추가 가정들이 있던 이 보조정리의 원래 판은 이 정리로 대체되었다.
\end{proof}

\begin{lemma}
\label{obsolete-lemma-fully-faithful-simple-one}
$A$를 뇌터 환이라 하자. $f \in \mathfrak a$를 $A$의 아이디얼의
한 원소라 하자. $U = \Spec(A) \setminus V(\mathfrak a)$로 놓는다.
다음을 가정하자.
\begin{enumerate}
\item $A$에는 쌍대화 복합체가 있고 이는 $f$에 관해 완비이며,
\item $A_f$는 $(S_2)$이고, 모든 극소 소 아이디얼
$\mathfrak p \subset A$에 대해 $f \not \in \mathfrak p$이며,
$\mathfrak q \in V(\mathfrak p) \cap V(\mathfrak a)$이면
$\dim((A/\mathfrak p)_\mathfrak q) \geq 3$이다.
\end{enumerate}
그러면 완비화 함자
$$
\textit{Coh}(\mathcal{O}_U)
\longrightarrow
\textit{Coh}(U, I\mathcal{O}_U),
\quad
\mathcal{F} \longmapsto \mathcal{F}^\wedge
$$
는 유한 국소 자유 대상들의 충만한 부분범주 위에서 완전충실하다.
\end{lemma}

\begin{proof}
이 보조정리는 『대수 및 형식 기하학』의 보조정리
\ref{algebraization-lemma-fully-faithful-simple-two}의 특수한 경우이다.
\end{proof}

\section{ZMT와 관련된 보조정리}
\label{obsolete-section-ZMT}

\noindent
이 절의 보조정리들은 원래 Zariski 주정리의 (대수적 판인)
『대수학』의 정리 \ref{algebra-theorem-main-theorem}의 증명에 사용되었다.

\begin{lemma}
\label{obsolete-lemma-change-equation-multiply}
$R$을 환이라 하고 $\varphi : R[x] \to S$를 환 준동형이라 하자.
$t \in S$라 하자. $t$가 $R[x]$ 위에서 정수적이면, 모든 $a \in R$에 대해
$\varphi(a)^\ell t$가 $\varphi_a : R[y] \to S$ 위에서 정수적이 되게 하는
$\ell \geq 0$이 존재한다. 여기서 이 사상은
$y \mapsto \varphi(ax)$와 $r \mapsto \varphi(r)$ ($r\in R$)로 정의된다.
\end{lemma}

\begin{proof}
$f_i \in R[x]$에 대해
$t^d + \sum_{i < d} \varphi(f_i)t^i = 0$이라 하자.
$\ell$을 모든 $f_i$의 $x$에 관한 차수 가운데 최댓값이라 하자.
등식에 $\varphi(a)^\ell$을 곱하면
$\varphi(a)^\ell t^d + \sum_{i < d} \varphi(a^\ell f_i)t^i = 0$을 얻는다.
각 $\varphi(a^\ell f_i)$가 $\varphi_a$의 상에 속함에 유의하자.
『대수학』의 보조정리 \ref{algebra-lemma-make-integral-trivial}에서 결과가
따른다.
\end{proof}

\begin{lemma}
\label{obsolete-lemma-make-integral-less-trivial}
$\varphi : R \to S$를 환 준동형이라 하자.
$t \in S$가 관계식
$\varphi(a_0) + \varphi(a_1)t + \ldots + \varphi(a_n) t^n = 0$을
만족한다고 하자. $u_n = \varphi(a_n)$,
$u_{n-1} = u_n t + \varphi(a_{n-1})$로 놓고,
$u_1 = u_2 t + \varphi(a_1)$에 이를 때까지 같은 방식으로 정의한다.
그러면 $u_n, u_{n-1}, \ldots, u_1$과
$u_nt, u_{n-1}t, \ldots, u_1t$는 모두 $R$ 위에서 정수적이고,
$S$의 아이디얼 $(\varphi(a_0), \ldots, \varphi(a_n))$과
$(u_n, \ldots, u_1)$은 같다.
\end{lemma}

\begin{proof}
$n$에 대한 귀납법으로 증명한다. $u_n = \varphi(a_n)$이므로
『대수학』의 보조정리 \ref{algebra-lemma-make-integral-trivial}에서
$u_nt$가 $R$ 위에서 정수적임을 얻는다. 물론
$u_n = \varphi(a_n)$도 $R$ 위에서 정수적이다. 따라서
$u_{n - 1} = u_n t  + \varphi(a_{n - 1})$는 $R$ 위에서 정수적이다
(『대수학』의 보조정리 \ref{algebra-lemma-integral-closure-is-ring}을 보라).
또한
$$
\varphi(a_0) + \varphi(a_1)t + \ldots + \varphi(a_{n - 1})t^{n - 1} +
u_{n - 1}t^{n - 1} = 0.
$$
이다. $S'$을 $S$ 안에서 $R$의 정수적 폐포라 하자. 사상 $S' \to S$와
위 등식에 귀납 가정을 적용하면
$u_{n-1}, \ldots, u_1$과 $u_{n-1}t, \ldots, u_1t$도 모두 $S'$에
속함을 알 수 있다. 아이디얼에 관한 명제는 원소들의 형태와
$u_1t + \varphi(a_0) = 0$이라는 사실에서 즉시 따른다.
\end{proof}

\begin{lemma}
\label{obsolete-lemma-make-integral-not-in-ideal}
$\varphi : R \to S$를 환 준동형이라 하자.
$t \in S$가 관계식
$\varphi(a_0) + \varphi(a_1)t + \ldots + \varphi(a_n) t^n = 0$을
만족한다고 하자. $J \subset S$를 어떤 하나 이상의 $i$에 대해
$\varphi(a_i) \not \in J$인 아이디얼이라 하자.
그러면 $u \not\in J$이고 $u$와 $ut$가 모두 $R$ 위에서 정수적인
$u \in S$가 존재한다.
\end{lemma}

\begin{proof}
원소 $u_i$ 가운데 하나는 $J$에 속하지 않으므로, 보조정리
\ref{obsolete-lemma-make-integral-less-trivial}에서 즉시 따른다.
\end{proof}

\noindent
다음 두 보조정리는 하나의 (퇴화하지 않은) 방정식으로 잘라 낸
$\mathbf{P}^1_R$의 닫힌 부분스킴을 기술하는 방법이다.

\begin{lemma}
\label{obsolete-lemma-P1}
$R$을 환이라 하자. $F(X, Y) \in R[X, Y]$를 차수가 $d$인 동차다항식이라
하자. $R$의 모든 소 아이디얼 $\mathfrak p$에 대해 $F$의 계수 가운데
적어도 하나는 $\mathfrak p$에 속하지 않는다고 가정하자.
$S = R[X, Y]/(F)$를 등급환으로 보자. 그러면 모든 $n \geq d$에 대해
$R$-가군 $S_n$은 계수 $d$인 유한 국소 자유 가군이다.
\end{lemma}

\begin{proof}
$R$-가군 $S_n$은 다음 표시를 갖는다.
$$
R[X, Y]_{n-d} \to R[X, Y]_n \to S_n \to 0.
$$
따라서 『대수학』의 보조정리 \ref{algebra-lemma-cokernel-flat}에 의해,
모든 소 아이디얼 $\mathfrak p$에 대해 $F$를 곱하는 사상
$\kappa(\mathfrak p)[X, Y] \to \kappa(\mathfrak p)[X, Y]$가 단사임을
보이면 충분하다. $F$가 $\mathfrak p$로 나눈 나머지에서 영다항식으로
사상되지 않는다는 가정으로부터 이는 명백하다. 계수에 관한 주장도
이로부터 명백하다.
\end{proof}

\begin{lemma}
\label{obsolete-lemma-rel-prime-pols}
$k$를 체라 하자. $F, G \in k[X, Y]$를 각각 차수가 $d, e$인
동차다항식이라 하자. $F, G$가 서로소라고 가정하자.
그러면 $G$를 곱하는 사상은 $S = k[X, Y]/(F)$ 위에서 단사이다.
\end{lemma}

\begin{proof}
이는 ``서로소''를 정의하는 한 방법이다. 다른 정의를 사용한다면,
그 정의가 이 정의와 동치임을 보일 수 있다.
\end{proof}

\begin{lemma}
\label{obsolete-lemma-P1-localize}
$R$을 환이라 하자. $F(X, Y) \in R[X, Y]$를 차수가 $d$인 동차다항식이라
하고 $S = R[X, Y]/(F)$를 등급환으로 보자.
$\mathfrak p \subset R$를 소 아이디얼이라 하고, $F$의 어떤 계수가
$\mathfrak p$에 속하지 않는다고 하자. 모든 $n \geq d$에 대해 $G$를 곱하는 사상이 동형
$(S_n)_f \to (S_{n + e})_f$을 유도하도록 하는
$f \in R$, $f \not\in \mathfrak p$, 정수 $e$, 그리고
$G \in R[X, Y]_e$가 존재한다.
\end{lemma}

\begin{proof}
증명하는 동안 $f\in R$, $f\not\in \mathfrak p$에 대해 $R$을 $R_f$로
(유한 번) 바꾸어도 된다. 첫 단계로, $F$의 어떤 계수가 $R$에서 가역이
되도록 이렇게 바꾼다. 특히 보조정리 \ref{obsolete-lemma-P1}에 의해 이제
$n \geq d$에 대해 가군 $S_n$은 계수 $d$인 국소 자유 가군이다.
$\kappa(\mathfrak p)[X, Y]$에서의 $G$의 상이 $F(X, Y)$의 상과 서로소가
되도록 하는 임의의 $G \in R[X, Y]_e$를 택한다(어떤 $e$에 대해서는
가능하다). $G$를 곱하여 유도되는 사상 $S_d \to S_{d + e}$에
『대수학』의 보조정리 \ref{algebra-lemma-cokernel-flat}을 적용한다.
$G$의 선택과 보조정리 \ref{obsolete-lemma-rel-prime-pols}에 의해
$S_d \otimes \kappa(\mathfrak p) \to S_{d + e} \otimes \kappa(\mathfrak p)$
는 전단사이다. 따라서 적절한 $f$에 대해 $R$을 $R_f$로 바꾸면
$G : S_d \to S_{d + e}$가 전단사라고 가정할 수 있다.
그러면 $R$의 모든 소 아이디얼 $\mathfrak p'$에 대해
$\kappa(\mathfrak p')[X, Y]$에서의 $G$의 상은 $F$의 상과 서로소이다.
이제 『대수학』의 보조정리 \ref{algebra-lemma-cokernel-flat}을 다시
적용하면 모든 사상
$G : S_d \to S_{d + e}$, $n \geq d$가 동형임을 알 수 있다.
\end{proof}

\begin{remark}
\label{obsolete-remark-algebra}
$R$을 환이라 하자. $F \in R[X, Y]_d$와 $G \in R[X, Y]_e$가 주어졌고,
$S = R[X, Y]/(F)$로 놓았을 때 다음이 성립한다고 하자.
(1) 모든 $n \geq d$에 대해 $S_n$은 계수 $d$인 유한 국소 자유 가군이고,
(2) 모든 $n \geq d$에 대해 $G$를 곱하는 사상은 동형
$S_n \to S_{n + e}$을 정의한다. 이 경우 다음과 같이 유한 국소 자유
$R$-대수 $A$를 정의할 수 있다.
\begin{enumerate}
\item $R$-가군으로서 $A = S_{ed}$이고,
\item 곱셈 $A \times A \to A$는 $S_{2ed}$에서
$G^d H_3 = H_1 H_2$일 때 그리고 그때에만 $H_1 H_2 = H_3$이 되도록
정의한다.
\end{enumerate}
$G^d$를 곱하는 사상이 전단사 $S_{de} \to S_{2de}$를 유도하므로 이는
잘 정의된다. 이것이 환 구조를 정의함은 쉽게 알 수 있다. 혼동하기 쉬운
사실은 원소 $G^d$가 환 $A$의 단위원을 정의한다는 점이다.
\end{remark}

\begin{lemma}
\label{obsolete-lemma-finite-after-localization}
$R$을 환이라 하고 $f \in R$라 하자. $S$, $S'$과 다음 가환 도표의
실선 화살표들이 주어졌다고 하자.
$$
\xymatrix{
& S'' \ar@{-->}[rd] \ar@{-->}[dd] &
\\
R \ar[rr] \ar@{-->}[ru] \ar[d] &  & S \ar[d]
\\
R_f \ar[r] & S' \ar[r] & S_f
}
$$
$R_f \to S'$가 유한이라고 가정하자. 그러면 $S' = (S'')_f$가 되도록
도표에 표시된 점선 화살표들과 유한 환 준동형 $R \to S''$을 찾을 수 있다.
\end{lemma}

\begin{proof}
$S'$이 $R_f$ 위에서 $x_i$, $i = 1, \ldots, w$로 생성된다고 하자.
$P_i(x_i) = 0$인 일계수 다항식 $P_i(t) \in R_f[t]$를 택한다.
$P_i$의 차수가 $d_i > 0$이라고 하자. 하나의 공통된 $n$에 대해
$P_i(t) = t^{d_i} + \sum_{j < d_i} (a_{ij}/f^n) t^j$라고 쓴다.
또한 적절한 $g_i \in S$에 대해 $S_f$에서 $x_i$의 상을
$g_i / f^n$으로 쓴다. 그러면 $S$의 원소
$\xi_i = f^{nd_i} g_i^{d_i} + \sum_{j < d_i} f^{n(d_i - j)} a_{ij} g_i^j$는
$f$의 어떤 거듭제곱에 의해 소멸한다. 따라서 $n$을 $n'$으로 늘리면
$g_i$가 $f^{n' - n}g_i$로 바뀌고, $\xi_i = 0$이라고 가정할 수 있다.
이제 $S'$은 원소 $f^n x_i$들로 생성되며, 각 원소는 $R$에 계수를 갖는
일계수 다항식 $Q_i(t) = t^{d_i} +
\sum_{j < d_i} f^{n(d_i - j)} a_{ij} t^j$의 영점이다.
또한 구성에 의해 $S$에서 $Q_i(f^ng_i) = 0$이다.
따라서 유한 $R$-대수
$S'' = R[z_1, \ldots, z_w]/(Q_1(z_1), \ldots, Q_w(z_w))$를 얻으며,
이는 위와 같은 가환 도표에 들어간다. 사상 $\alpha : S'' \to S$는
$z_i$를 $f^ng_i$로 보내고, 사상 $\beta : S'' \to S'$는
$z_i$를 $f^nx_i$로 보낸다. 아직은 $\beta$가 동형
$(S'')_f \cong S'$를 유도하지 않을 수도 있다. 지금까지는 이 사상이
전사라는 것만 안다. 문제는 $S'$에서 영으로 사상되지만 영은 아닌 원소
$h/f^n \in (S'')_f$가 있을 수 있다는 것이다. 이 경우 $\beta(h)$는
어떤 $N$에 대해 $f^N \beta(h) = 0$을 만족하는 $S$의 원소이다.
따라서 $f^N h$는 $S''$의 아이디얼
$J = \{h \in S'' \mid \alpha(h) = 0 \text{ 이고 }
\beta(h) = 0\}$의 원소이다. 좋다. 이제 $S''/J$가 원하는 대수임을
쉽게 알 수 있다.
\end{proof}

\section{형식적으로 매끄러운 환 준동형}
\label{obsolete-section-formally-smooth}

\begin{lemma}
\label{obsolete-lemma-formally-smooth-smooth}
$R$을 환이라 하고 $S$를 $R$-대수라 하자.
$S$가 $R$ 위에서 유한 표시이고 형식적으로 매끄러우면,
$S$는 $R$ 위에서 매끄럽다.
\end{lemma}

\begin{proof}
『대수학』의 명제 \ref{algebra-proposition-smooth-formally-smooth}를 보라.
\end{proof}

\begin{remark}
\label{obsolete-remark-equation-derivatives}
이 태그는 『형식 공간의 대수화』의 명제
\ref{restricted-proposition-approximate}의 증명에 있던 한 등식을
가리켰으나, 내용을 재배열하면서 더 이상 쓰이지 않게 되었다.
\end{remark}

\begin{remark}
\label{obsolete-remark-equation-ci}
이 태그는 『형식 공간의 대수화』의 명제
\ref{restricted-proposition-approximate}의 증명에 있던 한 등식을
가리켰으나, 내용을 재배열하면서 더 이상 쓰이지 않게 되었다.
\end{remark}

\begin{remark}
\label{obsolete-remark-equation-in-ideal}
이 태그는 『형식 공간의 대수화』의 명제
\ref{restricted-proposition-approximate}의 증명에 있던 한 등식을
가리켰으나, 내용을 재배열하면서 더 이상 쓰이지 않게 되었다.
\end{remark}

\begin{remark}
\label{obsolete-remark-equation-derivatives-analogue}
이 태그는 『형식 공간의 대수화』의 명제
\ref{restricted-proposition-approximate}의 증명에 있던 한 등식을
가리켰으나, 내용을 재배열하면서 더 이상 쓰이지 않게 되었다.
\end{remark}

\begin{remark}
\label{obsolete-remark-equation-go-down}
이 태그는 『형식 공간의 대수화』의 보조정리
\ref{restricted-lemma-lift-approximation}의 증명에 있던 한 등식을
가리켰으나, 내용을 재배열하면서 더 이상 쓰이지 않게 되었다.
\end{remark}

\begin{lemma}
\label{obsolete-lemma-get-morphism-general}
$A$를 뇌터 환이라 하고 $I \subset A$를 아이디얼이라 하자.
$t$를 $I$의 최소 생성원 수라 하자.
$C$를 $I$-진 완비 뇌터 $A$-대수라 하자.
$I \subset A \to C$에만 의존하고 다음 성질을 갖는 정수
$d \geq 0$이 존재한다. 다음이 주어졌다고 하자.
\begin{enumerate}
\item $c \geq 0$과 『형식 공간의 대수화』의 등식
(\ref{restricted-equation-C-prime})에 나오는 $B$가 주어지고,
$a \in I^c$에 대해 $\NL_{B/A}^\wedge$ 위에서 $a$를 곱하는 사상이
$D(B)$에서 영이며,
\item 정수 $n > 2t\max(c, d)$와,
\item $A/I^n$-대수 준동형 $\psi_n : B/I^nB \to C/I^nC$가 주어졌다.
\end{enumerate}
그러면 $m = \lfloor \frac{n}{t} \rfloor$일 때
$\psi_n \bmod I^{m - c} = \varphi \bmod I^{m - c}$를 만족하는
$A$-대수 준동형 $\varphi : B \to C$가 존재한다.
\end{lemma}

\begin{proof}
이 보조정리는 더 강한 『형식 공간의 대수화』의 보조정리
\ref{restricted-lemma-get-morphism-general-better}로 대체되었다.
실제로 그 보조정리에서 이 보조정리를 도출하겠다.

\medskip\noindent
위 명제와 같은 $I \subset A \to C$가 주어졌다고 하자.
『국소 코호몰로지』의 절 \ref{local-cohomology-section-uniform}에서 얻은
정수들을 $d(\text{Gr}_I(C))$와 $q(\text{Gr}_I(C))$로 나타내자.
$t$는 $IC$의 최소 생성원 수의 상계이고, 따라서 『국소 코호몰로지』의
절 \ref{local-cohomology-section-uniform}의 논의에 의해
$d(\text{Gr}_I(C)) + 1 \leq t$이다.
보조정리에서 실제 내용이 있는 경우만 생각하면 되므로 $t \geq 1$이라고
가정한다. 다음 값에 대해 보조정리가 참이라고 주장한다.
$$
d = q(\text{Gr}_I(C))
$$
실제로 $c$, $B$, $n$, $\psi_n$이 위 명제와 같이 주어졌다고 하자.
그러면
$$
n > 2t\max(c, d) \Rightarrow n \geq 2tc + 1 \Rightarrow
n \geq 2(d(\text{Gr}_I(C)) + 1)c + 1
$$
이고, 한편
$$
n > 2t\max(c, d) \Rightarrow n > t(c + d) \Rightarrow
n \geq q(C) + tc \geq q(\text{Gr}_I(C)) + (d(\text{Gr}_I(C)) + 1)c
$$
이다. 따라서 『형식 공간의 대수화』의 보조정리
\ref{restricted-lemma-get-morphism-general-better}의 가정들이 성립하고,
$I^{n - (d(\text{Gr}_I(C)) + 1)c}C$를 법으로 $\psi_n$과 합동인
$A$-대수 준동형 $\varphi : B \to C$를 얻는다.
다음 계산에 의해
\begin{align*}
n - (d(\text{Gr}_I(C)) + 1)c
& = \frac{n}{t} + \frac{(t - 1)n}{t} - (d(\text{Gr}_I(C)) + 1)c \\
& \geq \frac{n}{t} + \frac{(d(\text{Gr}_I(C))n}{t} - (d(\text{Gr}_I(C)) + 1)c \\
& > \frac{n}{t} + \frac{d(\text{Gr}_I(C))2tc}{t} - (d(\text{Gr}_I(C)) + 1)c \\
& = \frac{n}{t} + 2d(\text{Gr}_I(C))c - (d(\text{Gr}_I(C)) + 1)c \\
& = \frac{n}{t} + d(\text{Gr}_I(C))c - c \\
& \geq m - c
\end{align*}
원하는 대로 $I^{m - c}C$를 법으로 $\varphi$와 $\psi_n$이 합동이다.
\end{proof}

\section{사이트와 층}
\label{obsolete-section-sites}

\begin{remark}[아래 느낌표 함자에서 밂앞으로 가는 사상은 없다]
\label{obsolete-remark-from-shriek-to-star}
$U$를 사이트 $\mathcal{C}$의 대상이라 하자. $\mathcal{C}/U$ 위의 임의의
아벨 층 $\mathcal{G}$에 대해 표준 사상
$$
c : j_{U!}\mathcal{G} \longrightarrow j_{U*}\mathcal{G}
$$
이 있는지 물을 수 있다. 그런 사상을 구성하는 것은
$j_U^{-1}j_{U!}\mathcal{G} \to \mathcal{G}$를 구성하는 것과 같다.
제한은 층화와 가환함에 유의하자. 따라서 『사이트 위의 가군』의 보조정리
\ref{sites-modules-lemma-extension-by-zero}에 나오는 준층을 사용할 수 있다.
그러므로 제한 사상들과 호환되는 사상
$$
\bigoplus\nolimits_{\varphi \in \Mor_\mathcal{C}(V, U)}
\mathcal{G}(V \xrightarrow{\varphi} U)
\longrightarrow
\mathcal{G}(V/U)
$$
을 각 $V/U$에 대해 정의하면 충분하다. $\varphi$가 구조 사상
$\varphi_0 : V \to U$인 성분에서는 $1$을 취하고 다른 모든 성분에서는
영인 사상을 취할 수 있을 것처럼 보인다. 그러나 이는 제한 사상과 호환되지
않는다. 실제로 $\alpha : V' \to V$가 $\mathcal{C}$의 사상이면,
구조 사상이 $\varphi'_0 = \varphi_0 \circ \alpha$인
$\mathcal{C}/U$의 대상을 $V'/U$라 나타내자. 다음 도표가 가환인지
확인해야 한다.
$$
\xymatrix{
\bigoplus\nolimits_{\varphi \in \Mor_\mathcal{C}(V, U)}
\mathcal{G}(V \xrightarrow{\varphi} U)
\ar[d] \ar[r] &
\mathcal{G}(V/U) \ar[d] \\
\bigoplus\nolimits_{\varphi' \in \Mor_\mathcal{C}(V', U)}
\mathcal{G}(V' \xrightarrow{\varphi'} U)
\ar[r] &
\mathcal{G}(V'/U)
}
$$
문제는 $\varphi \circ \alpha = \varphi'_0$을 만족하면서
$\varphi_0$와는 다른 사상 $\varphi : V \to U$가 있을 수 있다는 점이다.
그러면 왼쪽 수직 화살표는 $\varphi$에 대응하는 성분을 아래 수평 화살표가
$1$인 성분으로 보내므로, 도표는 거의 확실히 가환하지 않는다.
\end{remark}

\section{코호몰로지}
\label{obsolete-section-cohomology}

\noindent
코호몰로지에 관한 폐기된 보조정리들이다.

\begin{lemma}
\label{obsolete-lemma-ML-general}
$I$를 환 $A$의 아이디얼이라 하고 $X$를 $\Spec(A)$ 위의 스킴이라 하자.
$$
\ldots \to \mathcal{F}_3 \to \mathcal{F}_2 \to \mathcal{F}_1
$$
을 $\mathcal{F}_n = \mathcal{F}_{n + 1}/I^n\mathcal{F}_{n + 1}$을
만족하는 $\mathcal{O}_X$-가군의 역계라 하자.
$$
\bigoplus\nolimits_{n \geq 0} H^1(X, I^n\mathcal{F}_{n + 1})
$$
가 등급 $\bigoplus_{n \geq 0} I^n/I^{n + 1}$-가군으로서 오름 사슬
조건을 만족한다고 가정하자. 그러면 역계
$M_n = \Gamma(X, \mathcal{F}_n)$은 Mittag--Leffler 조건을 만족한다.
\end{lemma}

\begin{proof}
더 일반적인 『코호몰로지』의 보조정리
\ref{cohomology-lemma-ML-general}의 특수한 경우이다.
\end{proof}

\begin{lemma}
\label{obsolete-lemma-ML-general-better}
$I$를 환 $A$의 아이디얼이라 하고 $X$를 $\Spec(A)$ 위의 스킴이라 하자.
$$
\ldots \to \mathcal{F}_3 \to \mathcal{F}_2 \to \mathcal{F}_1
$$
을 $\mathcal{F}_n = \mathcal{F}_{n + 1}/I^n\mathcal{F}_{n + 1}$을
만족하는 $\mathcal{O}_X$-가군의 역계라 하자. $n$이 주어졌을 때
$$
H^1_n =
\bigcap\nolimits_{m \geq n}
\Im\left(
H^1(X, I^n\mathcal{F}_{m + 1}) \to H^1(X, I^n\mathcal{F}_{n + 1})
\right)
$$
로 정의한다. $\bigoplus H^1_n$이 등급
$\bigoplus_{n \geq 0} I^n/I^{n + 1}$-가군으로서 오름 사슬 조건을
만족하면, 역계 $M_n = \Gamma(X, \mathcal{F}_n)$은 Mittag--Leffler
조건을 만족한다.
\end{lemma}

\begin{proof}
더 일반적인 『코호몰로지』의 보조정리
\ref{cohomology-lemma-ML-general-better}의 특수한 경우이다.
\end{proof}

\begin{lemma}
\label{obsolete-lemma-topology-I-adic-general}
$I$를 환 $A$의 유한 생성 아이디얼이라 하고 $X$를 $\Spec(A)$ 위의
스킴이라 하자.
$$
\ldots \to \mathcal{F}_3 \to \mathcal{F}_2 \to \mathcal{F}_1
$$
을 $\mathcal{F}_n = \mathcal{F}_{n + 1}/I^n\mathcal{F}_{n + 1}$을
만족하는 $\mathcal{O}_X$-가군의 역계라 하자.
$$
\bigoplus\nolimits_{n \geq 0} H^0(X, I^n\mathcal{F}_{n + 1})
$$
가 등급 $\bigoplus_{n \geq 0} I^n/I^{n + 1}$-가군으로서 오름 사슬
조건을 만족한다고 가정하자. 그러면
$M = \lim \Gamma(X, \mathcal{F}_n)$ 위의 극한 위상은 $I$-진 위상이다.
\end{lemma}

\begin{proof}
더 일반적인 『코호몰로지』의 보조정리
\ref{cohomology-lemma-topology-I-adic-general}의 특수한 경우이다.
\end{proof}

\begin{lemma}
\label{obsolete-lemma-pullback-K-flat}
$(\Sh(\mathcal{C}), \mathcal{O}_\mathcal{C})$를 환 달린 토포스라 하자.
$\mathcal{O}_\mathcal{C}$-가군의 임의의 복합체 $\mathcal{G}^\bullet$에
대해 준동형 $\mathcal{K}^\bullet \to \mathcal{G}^\bullet$이 존재하여,
환 달린 토포스의 임의의 사상
$f : (\Sh(\mathcal{D}), \mathcal{O}_\mathcal{D}) \to
(\Sh(\mathcal{C}), \mathcal{O}_\mathcal{C})$에 대해
$f^*\mathcal{K}^\bullet$은 $\mathcal{O}_\mathcal{D}$-가군의 K-평탄
복합체가 된다.
\end{lemma}

\begin{proof}
『사이트 위의 코호몰로지』의 보조정리
\ref{sites-cohomology-lemma-K-flat-resolution}과
\ref{sites-cohomology-lemma-pullback-K-flat}에서 따른다.
\end{proof}

\begin{remark}
\label{obsolete-remark-pullback-K-flat}
이 주석은 전에는 K-평탄 복합체의 당김뒤가 K-평탄인지 여부에 관해
알려진 내용을 논의했으나, 이제 『사이트 위의 코호몰로지』의 보조정리
\ref{sites-cohomology-lemma-pullback-K-flat}로 대체되었다.
\end{remark}

\noindent
다음 보조정리는 특수한 경우에 유도 극한의 코호몰로지 층을 계산한다.

\begin{lemma}
\label{obsolete-lemma-Rlim-of-system}
$(\mathcal{C}, \mathcal{O})$를 환 달린 사이트라 하자.
$(K_n)$을 $D(\mathcal{O})$의 대상들의 역계라 하자.
$\mathcal{B} \subset \Ob(\mathcal{C})$를 부분집합이라 하고
$d \in \mathbf{N}$이라 하자. 다음을 가정하자.
\begin{enumerate}
\item 모든 $n$에 대해 $K_n$은 $D^+(\mathcal{O})$의 대상이고,
\item $q \in \mathbf{Z}$마다 $n \geq n(q)$이면
$H^q(K_{n + 1}) \to H^q(K_n)$이 동형이 되게 하는 $n(q)$가 존재하며,
\item $\mathcal{C}$의 모든 대상에는 그 원소들이 $\mathcal{B}$에 속하는
피복이 있고,
\item 모든 $U \in \mathcal{B}$에 대해 $p > d$와 모든 $q$에 대하여
$H^p(U, H^q(K_n)) = 0$이다.
\end{enumerate}
그러면 모든 $m \in \mathbf{Z}$에 대해
$H^m(R\lim K_n) = \lim H^m(K_n)$이다.
\end{lemma}

\begin{proof}
$K = R\lim K_n$으로 놓고 $U \in \mathcal{B}$라 하자.
각 $n$에 대해 스펙트럼 열
$$
H^p(U, H^q(K_n)) \Rightarrow H^{p + q}(U, K_n)
$$
이 있다. $K_n$이 아래로 유계이므로 이 열은 수렴한다. 『유도 범주』의
보조정리 \ref{derived-lemma-two-ss-complex-functor}를 보라.
$m \in \mathbf{Z}$를 고정하면 가정 (4)에 의해 $H^m(U, K_n)$에
기여하는 항은 $0 \leq p \leq d$ 및 $m - d \leq q \leq m$인
$H^p(U, H^q(K_n))$뿐이다. 가정 (2)에 의해
$n \geq \max{n(m), \ldots, n(m - d)}$이면
$H^m(U, K_{n + 1}) \to H^m(U, K_n)$은 동형이다.
『단사 대상』의 보조정리
\ref{injectives-lemma-RF-commutes-with-Rlim}에 의해 함자
$R\Gamma(U, -)$는 유도 극한과 가환한다. 따라서
$$
H^m(U, K) = H^m(R\lim R\Gamma(U, K_n))
$$
이다. 한편 방금 복합체 $R\Gamma(U, K_n)$의 코호몰로지 군이 결국
상수가 됨을 보였다. 그러므로 『대수학 심화』의 주석
\ref{more-algebra-remark-compare-derived-limit}에 의해,
$U \in \mathcal{B}$와 무관한 어떤 경계에 대해 모든 $n \gg 0$에서
$H^m(U, K)$는 $H^m(U, K_n)$과 같다. 그런 $n$을 택한다.
마지막으로 $H^m(K)$는 준층 $U \mapsto H^m(U, K)$의 층화이고
$H^m(K_n)$은 준층 $U \mapsto H^m(U, K_n)$의 층화임을 상기하자.
이 준층들은 $\mathcal{B}$의 원소들 위에서 같은 값을 갖는다. 따라서
가정 (3)은 두 층화도 같음을 보장한다. 보조정리가 따른다.
\end{proof}

\begin{lemma}
\label{obsolete-lemma-trivialities-cohomological-descent-abelian}
『단체적 공간』의 상황
\ref{spaces-simplicial-situation-simplicial-site}에서,
『단체적 공간』의 주석
\ref{spaces-simplicial-remark-augmentation-site}에 나오는 것처럼
$a_0$를 사이트 $\mathcal{D}$를 향하는 증대라 하자.
엄밀히 충만한 약한 Serre 부분범주
$$
\mathcal{A} \subset \textit{Ab}(\mathcal{D}),\quad
\mathcal{A}_n \subset \textit{Ab}(\mathcal{C}_n)
$$
이 주어졌다고 하자. 그러면
\begin{enumerate}
\item[(1)] 모든 $n$에 대해 $\mathcal{C}_n$으로 제한한 것이
$\mathcal{A}_n$에 속하는 $\mathcal{C}_{total}$ 위의 아벨 층
$\mathcal{F}$들의 모임은 엄밀히 충만한 약한 Serre 부분범주
$\mathcal{A}_{total} \subset \textit{Ab}(\mathcal{C}_{total})$이다.
\end{enumerate}
모든 $n$에 대해 $a_n^{-1}$이 $\mathcal{A}$를 $\mathcal{A}_n$으로
보내면
\begin{enumerate}
\item[(2)] $a^{-1}$은 $\mathcal{A}$를 $\mathcal{A}_{total}$로 보내고,
\item[(3)] $a^{-1}$은 $D_\mathcal{A}(\mathcal{D})$를
$D_{\mathcal{A}_{total}}(\mathcal{C}_{total})$로 보낸다.
\end{enumerate}
모든 $n, q$에 대해 $R^qa_{n, *}$가 $\mathcal{A}_n$을
$\mathcal{A}$로 보내면
\begin{enumerate}
\item[(4)] 모든 $q$에 대해 $R^qa_*$는 $\mathcal{A}_{total}$을
$\mathcal{A}$로 보내고,
\item[(5)] $Ra_*$는 $D_{\mathcal{A}_{total}}^+(\mathcal{C}_{total})$을
$D_\mathcal{A}^+(\mathcal{D})$로 보낸다.
\end{enumerate}
\end{lemma}

\begin{proof}
내용이 있는 주장은 (4)와 (5)뿐이다. (4)는 『단체적 공간』의 보조정리
\ref{spaces-simplicial-lemma-augmentation-spectral-sequence}의 스펙트럼
열과 『호몰로지』의 보조정리
\ref{homology-lemma-biregular-ss-converges}에서 따른다.
이어서 (5)는 절단으로 주어지는
$D_{\mathcal{A}_{total}}^+(\mathcal{C}_{total})$의 대상 $K$ 위의 표준
여과에 딸린 스펙트럼 열을 생각하면 따른다. 세부 사항은 생략한다.
\end{proof}

\begin{remark}
\label{obsolete-remark-cohomology-topics}
이 태그는 전에는 다룰 주제들을 나열한 코호몰로지 장의 한 절을 가리켰다.
\end{remark}

\begin{remark}
\label{obsolete-remark-sites-cohomology-topics}
이 태그는 전에는 다룰 주제들을 나열한 코호몰로지 장의 한 절을 가리켰다.
\end{remark}

\begin{remark}
\label{obsolete-remark-V-implies-C}
이 태그는 『사이트 위의 코호몰로지』의 보조정리
\ref{sites-cohomology-lemma-V-implies-C-general}의 특수한 경우를 가리켰다.
이는 『사이트 위의 코호몰로지』의 보조정리
\ref{sites-cohomology-lemma-compare-qc-zar}에 기술된 상황에 관한 것이다.
\end{remark}

\begin{remark}
\label{obsolete-remark-V-implies-cohomology}
이 태그는 『사이트 위의 코호몰로지』의 보조정리
\ref{sites-cohomology-lemma-V-implies-cohomology-general}의 특수한 경우를
가리켰다. 이는 『사이트 위의 코호몰로지』의 보조정리
\ref{sites-cohomology-lemma-compare-qc-zar}에 기술된 상황에 관한 것이다.
\end{remark}

\begin{remark}
\label{obsolete-remark-induction-step-V-C}
이 태그는 『사이트 위의 코호몰로지』의 보조정리
\ref{sites-cohomology-lemma-induction-step-V-C-general}의 특수한 경우를
가리켰다. 이는 『사이트 위의 코호몰로지』의 보조정리
\ref{sites-cohomology-lemma-compare-qc-zar}에 기술된 상황에 관한 것이다.
\end{remark}

\begin{remark}
\label{obsolete-remark-V-implies-C-etale-fppf}
이 태그는 『사이트 위의 코호몰로지』의 보조정리
\ref{sites-cohomology-lemma-V-implies-C-general}의 특수한 경우를 가리켰다.
이는 『에탈 코호몰로지』의 보조정리
\ref{etale-cohomology-lemma-compare-fppf-etale}에 기술된 상황에 관한 것이다.
\end{remark}

\begin{remark}
\label{obsolete-remark-V-implies-cohomology-etale-fppf}
이 태그는 『사이트 위의 코호몰로지』의 보조정리
\ref{sites-cohomology-lemma-V-implies-cohomology-general}의 특수한 경우를
가리켰다. 이는 『에탈 코호몰로지』의 보조정리
\ref{etale-cohomology-lemma-compare-fppf-etale}에 기술된 상황에 관한 것이다.
\end{remark}

\begin{remark}
\label{obsolete-remark-induction-step-V-C-etale-fppf}
이 태그는 『사이트 위의 코호몰로지』의 보조정리
\ref{sites-cohomology-lemma-induction-step-V-C-general}의 특수한 경우를
가리켰다. 이는 『에탈 코호몰로지』의 보조정리
\ref{etale-cohomology-lemma-compare-fppf-etale}에 기술된 상황에 관한 것이다.
\end{remark}

\begin{remark}
\label{obsolete-remark-V-implies-C-etale-ph}
이 태그는 『사이트 위의 코호몰로지』의 보조정리
\ref{sites-cohomology-lemma-V-implies-C-general}의 특수한 경우를 가리켰다.
이는 『에탈 코호몰로지』의 보조정리
\ref{etale-cohomology-lemma-compare-ph-etale}에 기술된 상황에 관한 것이다.
\end{remark}

\begin{remark}
\label{obsolete-remark-V-implies-cohomology-etale-ph}
이 태그는 『사이트 위의 코호몰로지』의 보조정리
\ref{sites-cohomology-lemma-V-implies-cohomology-general}의 특수한 경우를
가리켰다. 이는 『에탈 코호몰로지』의 보조정리
\ref{etale-cohomology-lemma-compare-ph-etale}에 기술된 상황에 관한 것이다.
\end{remark}

\begin{remark}
\label{obsolete-remark-V-implies-cohomology-etale-ph-extra}
이 태그는 『사이트 위의 코호몰로지』의 보조정리
\ref{sites-cohomology-lemma-V-implies-cohomology-extra-general}의 특수한
경우를 가리켰다. 이는 『에탈 코호몰로지』의 보조정리
\ref{etale-cohomology-lemma-compare-ph-etale}에 기술된 상황에 관한 것이다.
\end{remark}

\begin{remark}
\label{obsolete-remark-make-class-zero}
이 태그는 『사이트 위의 코호몰로지』의 보조정리
\ref{sites-cohomology-lemma-make-class-zero-general}의 특수한 경우를
가리켰다. 이는 『에탈 코호몰로지』의 보조정리
\ref{etale-cohomology-lemma-compare-ph-etale}에 기술된 상황에 관한 것이다.
\end{remark}

\begin{remark}
\label{obsolete-remark-induction-step-V-C-etale-ph}
이 태그는 『사이트 위의 코호몰로지』의 보조정리
\ref{sites-cohomology-lemma-induction-step-V-C-general}의 특수한 경우를
가리켰다. 이는 『에탈 코호몰로지』의 보조정리
\ref{etale-cohomology-lemma-compare-ph-etale}에 기술된 상황에 관한 것이다.
\end{remark}

\begin{remark}
\label{obsolete-remark-V-implies-C-etale-h}
이 태그는 『사이트 위의 코호몰로지』의 보조정리
\ref{sites-cohomology-lemma-V-implies-C-general}의 특수한 경우를 가리켰다.
이는 『에탈 코호몰로지』의 보조정리
\ref{etale-cohomology-lemma-compare-h-etale}에 기술된 상황에 관한 것이다.
\end{remark}

\begin{remark}
\label{obsolete-remark-V-implies-cohomology-etale-h}
이 태그는 『사이트 위의 코호몰로지』의 보조정리
\ref{sites-cohomology-lemma-V-implies-cohomology-general}의 특수한 경우를
가리켰다. 이는 『에탈 코호몰로지』의 보조정리
\ref{etale-cohomology-lemma-compare-h-etale}에 기술된 상황에 관한 것이다.
\end{remark}

\begin{remark}
\label{obsolete-remark-V-implies-cohomology-etale-h-extra}
이 태그는 『사이트 위의 코호몰로지』의 보조정리
\ref{sites-cohomology-lemma-V-implies-cohomology-extra-general}의 특수한
경우를 가리켰다. 이는 『에탈 코호몰로지』의 보조정리
\ref{etale-cohomology-lemma-compare-h-etale}에 기술된 상황에 관한 것이다.
\end{remark}

\begin{remark}
\label{obsolete-remark-induction-step-V-C-etale-h}
이 태그는 『사이트 위의 코호몰로지』의 보조정리
\ref{sites-cohomology-lemma-induction-step-V-C-general}의 특수한 경우를
가리켰다. 이는 『에탈 코호몰로지』의 보조정리
\ref{etale-cohomology-lemma-compare-h-etale}에 기술된 상황에 관한 것이다.
\end{remark}

\begin{remark}
\label{obsolete-remark-how-used}
이 태그는 전에는 에탈 코호몰로지 장에 있었으나 재구성으로 인해 이제
그곳에 두기 적절하지 않다. 태그의 내용은 다음과 같았다.
『에탈 코호몰로지』의 보조정리 \ref{etale-cohomology-lemma-when-ctf}를
사용하면 $f : X \to Y$가 스킴의 분리 유한형 사상이고 $Y$가
뇌터이면 $Rf_!$가 함자
$D_{ctf}(X_\etale, \Lambda) \to D_{ctf}(Y_\etale, \Lambda)$를
유도함을 증명할 수 있다. $Y$가 체의 스펙트럼이고 $X$가 곡선인 경우의
이 논증의 예는 『대각합 공식』의 명제
\ref{trace-proposition-projective-curve-constructible-cohomology}에 있다.
\end{remark}

\begin{lemma}
\label{obsolete-lemma-glue-f-upper-shriek}
$f : X \to Y$를 스킴의 국소 준유한 사상이라 하자. 다음 성질을 갖는
유일한 함자 $f^! : \textit{Ab}(Y_\etale) \to
\textit{Ab}(X_\etale)$가 존재한다.
\begin{enumerate}
\item $f \circ j$가 분리인 임의의 열린 사상 $j : U \to X$에 대해
표준 동형 $j^! \circ f^! = (f \circ j)^!$이 있고,
\item $U \subset U' \subset X$에 대한 이 동형들은 『에탈 코호몰로지
심화』의 보조정리
\ref{more-etale-lemma-upper-shriek-restriction}의 동형들과 호환된다.
\end{enumerate}
\end{lemma}

\begin{proof}
『에탈 코호몰로지 심화』의 보조정리
\ref{more-etale-lemma-lqf-f-upper-shriek}와
\ref{more-etale-lemma-upper-shriek-restriction}에서 즉시 따른다.
\end{proof}

\begin{proposition}
\label{obsolete-proposition-lqf-shriek}
$f : X \to Y$를 국소 준유한 사상이라 하자. 다음 성질을 갖는 수반 함자
$f_! : \textit{Ab}(X_\etale) \to \textit{Ab}(Y_\etale)$와
$f^! : \textit{Ab}(Y_\etale) \to \textit{Ab}(X_\etale)$가 존재한다.
\begin{enumerate}
\item 함자 $f^!$는 『에탈 코호몰로지 심화』의 보조정리
\ref{more-etale-lemma-lqf-f-upper-shriek}에서 구성한 것이고,
\item $f \circ j$가 분리인 임의의 열린 사상 $j : U \to X$에 대해
표준 동형 $f_! \circ j_! = (f \circ j)_!$이 있으며,
\item $U \subset U' \subset X$에 대한 이 동형들은 『에탈 코호몰로지
심화』의 보조정리 \ref{more-etale-lemma-f-shriek-composition}의
동형들과 호환된다.
\end{enumerate}
\end{proposition}

\begin{proof}
『에탈 코호몰로지 심화』의 절
\ref{more-etale-section-finite-support}와
\ref{more-etale-section-duality-locally-quasi-finite}를 보라.
\end{proof}

\begin{lemma}
\label{obsolete-lemma-lqf-f-upper-shriek-stalk}
$f : X \to Y$를 국소 준유한인 스킴의 사상이라 하자.
아벨 군 $A$와 기하학적 점 $\overline{y} : \Spec(k) \to Y$에 대해
$f^!(\overline{y}_*A) = \prod\nolimits_{f(\overline{x}) = \overline{y}}
\overline{x}_*A$이다.
\end{lemma}

\begin{proof}
『에탈 코호몰로지 심화』의 보조정리
\ref{more-etale-lemma-lqf-f-upper-shriek}의 대응하는 명제에서 따른다.
\end{proof}

\begin{lemma}
\label{obsolete-lemma-lqf-f-shriek-composition}
$f : X \to Y$와 $g : Y \to Z$를 합성 가능한 스킴의 국소 준유한
사상이라 하자. 그러면 $g_! \circ f_! = (g \circ f)_!$이고
$f^! \circ g^! = (g \circ f)^!$이다.
\end{lemma}

\begin{proof}
『에탈 코호몰로지 심화』의 보조정리
\ref{more-etale-lemma-lqf-shriek-composition}과
\ref{more-etale-lemma-upper-shriek-restriction}을 결합하면 된다.
\end{proof}

\section{미분 등급 대수}
\label{obsolete-section-dga}


\begin{lemma}
\label{obsolete-lemma-P-not-preserved-base-change}
$(A, \text{d})$와 $(B, \text{d})$를 미분 등급 대수라 하자.
$N$을 성질 (P)를 갖는 미분 등급 $(A, B)$-쌍가군이라 하고,
$M$을 성질 (P)를 갖는 미분 등급 $A$-가군이라 하자.
그러면 $Q = M \otimes_A N$은 $D(B)$에서
$M \otimes_A^\mathbf{L} N$을 나타내는 미분 등급 $B$-가군이고,
미분 등급 부분가군들에 의한 여과
$$
0 = F_{-1}Q \subset F_0Q \subset F_1Q \subset \ldots \subset Q
$$
를 갖는다. 여기서 $Q = \bigcup F_pQ$이고, 포함사상
$F_iQ \to F_{i + 1}Q$는 허용 가능한 단사이며, 몫
$F_{i + 1}Q/F_iQ$는 미분 등급 $B$-가군으로서
$(A \otimes_R B)[k]$들의 직합과 동형이다.
\end{lemma}

\begin{proof}
$M$과 $N$ 위의 여과 $F_\bullet$을 택한다. 다음 식으로 주어지는
$Q = M \otimes_A N$ 위의 여과를 생각하자.
$$
F_n(Q) = \sum\nolimits_{i + j = n}
F_i(M) \otimes_A F_j(N)
$$
이는 명백히 미분 등급 $B$-부분가군이다. 예를 들어 $M$의 여과가
등급 $A$-가군의 범주에서 분할되므로
$$
F_n(Q)/F_{n - 1}(Q) =
\bigoplus\nolimits_{i + j = n}
F_i(M)/F_{i - 1}(M) \otimes_A F_j(N)/F_{j - 1}(N)
$$
임을 알 수 있다. 가정에 의해 오른쪽의 몫들은 각각 $A$ 및
$A \otimes_R B$의 옮김들의 직합과 동형이고,
$A \otimes_A (A \otimes_R B) = A \otimes_R B$이므로,
왼쪽은 미분 등급 $B$-가군으로서 $A \otimes_R B$의 옮김들의 직합이다.
(주의: $Q$에는 $(A, B)$-쌍가군 구조가 없다.)
이로써 보조정리의 첫 번째 명제가 증명된다. 두 번째 명제는
『미분 등급 대수』의 보조정리 \ref{dga-lemma-derived-bc}에서 함자를
정의한 방식으로부터 즉시 따른다.
\end{proof}

\section{단체적 방법}
\label{obsolete-section-simplicial}


\begin{lemma}
\label{obsolete-lemma-equiv}
『단체적 방법』의 보조정리 \ref{simplicial-lemma-section}과 같은 가정과
표기법을 사용하자. 사상 $f$의 단면 $g : U \to V$가 존재하고,
합성 $g \circ f$는 $V$ 위의 항등사상과 호모토피 동치이다.
특히 사상 $f$는 호모토피 동치이다.
\end{lemma}

\begin{proof}
『단체적 방법』의 보조정리 \ref{simplicial-lemma-section}과
\ref{simplicial-lemma-trivial-kan-homotopy}에서 즉시 따른다.
\end{proof}

\begin{lemma}
\label{obsolete-lemma-cosk-hom-deltak}
$\mathcal{C}$를 유한 쌍대곱과 유한 극한을 갖는 범주라 하자.
$X$를 $\mathcal{C}$의 대상이라 하고 $k \geq 0$이라 하자.
표준 사상
$$
\Hom(\Delta[k], X)
\longrightarrow
\text{cosk}_1 \text{sk}_1 \Hom(\Delta[k], X)
$$
은 동형이다.
\end{lemma}

\begin{proof}
임의의 단체적 대상 $V$에 대해
\begin{eqnarray*}
\Mor(V, \text{cosk}_1 \text{sk}_1 \Hom(\Delta[k], X))
& = &
\Mor(\text{sk}_1 V, \text{sk}_1 \Hom(\Delta[k], X)) \\
& = &
\Mor(i_{1!} \text{sk}_1 V, \Hom(\Delta[k], X)) \\
& = &
\Mor(i_{1!} \text{sk}_1 V \times \Delta[k], X)
\end{eqnarray*}
이다. 첫 번째 등식은 $\text{sk}$와 $\text{cosk}$의 수반성에서,
두 번째 등식은 $i_{1!}$와 $\text{sk}_1$의 수반성에서, 첫 번째 등식은
『단체적 방법』의 정의
\ref{simplicial-definition-hom-from-simplicial-set}에서 따른다. 마지막의
$X$는 값이 $X$인 상수 단체적 대상을 뜻한다.
『단체적 방법』의 보조정리 \ref{simplicial-lemma-augmentation-howto}에 의해
이 집합의 원소는 $i_{1!} \text{sk}_1 V \times \Delta[k]$의 차수 $0$과
$1$인 항들에만 의존한다. 이 항들은 $V \times \Delta[k]$의 차수 $0$과
$1$인 항들과 일치한다. 『단체적 방법』의 보조정리
\ref{simplicial-lemma-recovering-U-for-real}을 보라.
따라서 위 집합은
$\Mor(V \times \Delta[k], X) = \Mor(V, \Hom(\Delta[k], X))$와 같다.
\end{proof}

\begin{lemma}
\label{obsolete-lemma-cosk0-hom-deltak}
$\mathcal{C}$를 범주라 하자. $X$를 자기 곱
$X \times \ldots \times X$가 존재하는 $\mathcal{C}$의 대상이라 하자.
$k \geq 0$이라 하고, $C[k]$를 『단체적 방법』의 예
\ref{simplicial-example-simplex-cosimplicial-set}에 나오는 것으로 하자.
『단체적 방법』의 보조정리
\ref{simplicial-lemma-morphism-into-product}의 표기법을 사용하면,
표준 사상
$$
\Hom(C[k], X)_1
\longrightarrow
(\text{cosk}_0 \text{sk}_0 \Hom(C[k], X))_1
$$
은 사상
$$
\prod\nolimits_{\alpha : [k] \to [1]} X
\longrightarrow
X \times X
$$
과 동일시된다. 이 사상은 $\alpha$가 상수 사상인 인자들로의 사영이다.
\end{lemma}

\begin{proof}
이는 『초피복』의 보조정리 \ref{hypercovering-lemma-covering}의 증명에
나와 있다.
\end{proof}

\section{스킴에 관한 결과}
\label{obsolete-section-devissage}

\noindent
불필요해 보이는 보조정리들이다.

\begin{lemma}
\label{obsolete-lemma-stein-projective}
$(R, \mathfrak m, \kappa)$를 국소환이라 하자.
$X \subset \mathbf{P}^n_R$를 닫힌 부분스킴이라 하자.
$R = \Gamma(X, \mathcal{O}_X)$라고 가정하자. 그러면 특수 올
$X_k$는 기하학적으로 연결이다.
\end{lemma}

\begin{proof}
이는 『사상 심화』의 정리
\ref{more-morphisms-theorem-stein-factorization-general}의 특수한 경우이다.
\end{proof}

\begin{lemma}
\label{obsolete-lemma-property-irreducible-higher-rank}
$X$를 뇌터 스킴이라 하자.
$Z_0 \subset X$를 일반점 $\xi$를 갖는 기약 닫힌 부분집합이라 하자.
$\mathcal{P}$를 $X$ 위의 연접층의 성질이라 하고 다음을 가정하자.
\begin{enumerate}
\item 연접층의 임의의 짧은 완전열에서 세 층 가운데 두 층이 성질
$\mathcal{P}$를 가지면 나머지 층도 이 성질을 갖는다.
\item 연접층의 직합이 $\mathcal{P}$를 가지면 두 직합인자 모두
그 성질을 갖는다.
\item 모든 정수적 닫힌 부분스킴 $Z \subset Z_0 \subset X$,
$Z \not = Z_0$와 모든 준연접 아이디얼층
$\mathcal{I} \subset \mathcal{O}_Z$에 대해
$(Z \to X)_*\mathcal{I}$가 $\mathcal{P}$를 갖는다.
\item 다음을 만족하는 $X$ 위의 어떤 연접층 $\mathcal{G}$가 존재한다.
\begin{enumerate}
\item $\text{Supp}(\mathcal{G}) = Z_0$,
\item $\mathcal{G}_\xi$는 $\mathfrak m_\xi$에 의해 소멸하고,
\item $\mathcal{G}$는 성질 $\mathcal{P}$를 갖는다.
\end{enumerate}
\end{enumerate}
그러면 지지가 $Z_0$에 포함되는 $X$ 위의 모든 연접층 $\mathcal{F}$는
성질 $\mathcal{P}$를 갖는다.
\end{lemma}

\begin{proof}
증명은 『스킴의 코호몰로지』의 보조정리
\ref{coherent-lemma-property-irreducible}의 증명을 변형한 것이다.
그 증명과 정확히 같은 방식으로, 지지가 $Z_0$에 진부분집합으로 포함되는
모든 연접층이 성질 $\mathcal{P}$를 가짐을 알 수 있다.

\medskip\noindent
(3)에 나오는 연접층 $\mathcal{G}$를 생각하자.
『스킴의 코호몰로지』의 보조정리
\ref{coherent-lemma-prepare-filter-irreducible}에 의해 $Z_0$ 위의
아이디얼층 $\mathcal{I}$와 짧은 완전열
$$
0 \to
\left((Z_0 \to X)_*\mathcal{I}\right)^{\oplus r} \to
\mathcal{G} \to
\mathcal{Q} \to 0
$$
이 존재한다. 여기서 $\mathcal{Q}$의 지지는 $Z_0$에 진부분집합으로
포함된다. 특히 $\mathcal{G}$의 지지가 $Z$와 같으므로 $r > 0$이고
$\mathcal{I}$는 영이 아니다. $\mathcal{Q}$가 성질 $\mathcal{P}$를
가지므로 $\left((Z_0 \to X)_*\mathcal{I}\right)^{\oplus r}$도
성질 $\mathcal{P}$를 갖는다. (2)에 의해
$(Z_0 \to X)_*\mathcal{I}$가 성질 $\mathcal{P}$를 가짐을 얻는다.
이를 『스킴의 코호몰로지』의 보조정리
\ref{coherent-lemma-property-irreducible}의 증명 중 적절한 지점에
대입하면 보조정리가 따라 나온다. 몇몇 세부 사항은 생략한다.
\end{proof}

\begin{lemma}
\label{obsolete-lemma-property-higher-rank}
$X$를 뇌터 스킴이라 하자. $\mathcal{P}$를 $X$ 위의 연접층의
성질이라 하고 다음을 가정하자.
\begin{enumerate}
\item 연접층의 임의의 짧은 완전열에서 세 층 가운데 두 층이 성질
$\mathcal{P}$를 가지면 나머지 층도 이 성질을 갖는다.
\item 연접층의 직합이 $\mathcal{P}$를 가지면 두 직합인자 모두
그 성질을 갖는다.
\item 일반점 $\xi$를 갖는 모든 정수적 닫힌 부분스킴 $Z \subset X$에
대해 다음을 만족하는 어떤 연접층 $\mathcal{G}$가 존재한다.
\begin{enumerate}
\item $\text{Supp}(\mathcal{G}) = Z$,
\item $\mathcal{G}_\xi$는 $\mathfrak m_\xi$에 의해 소멸하고,
\item $\mathcal{G}$는 성질 $\mathcal{P}$를 갖는다.
\end{enumerate}
\end{enumerate}
그러면 $X$ 위의 모든 연접층은 성질 $\mathcal{P}$를 갖는다.
\end{lemma}

\begin{proof}
이 명제는 보조정리 \ref{obsolete-lemma-property-irreducible-higher-rank}에서
따른다. 그 방식은 『스킴의 코호몰로지』의 보조정리
\ref{coherent-lemma-property}가 『스킴의 코호몰로지』의 보조정리
\ref{coherent-lemma-property-irreducible}에서 따라 나오는 방식과 정확히
같다.
\end{proof}

\begin{lemma}
\label{obsolete-lemma-section-maps-back-into}
$X$를 스킴이라 하자. $\mathcal{L}$을 가역 $\mathcal{O}_X$-가군이라 하고
$s \in \Gamma(X, \mathcal{L})$를 절단이라 하자.
$\mathcal{F}' \subset \mathcal{F}$를 준연접 $\mathcal{O}_X$-가군이라
하자. 다음을 가정하자.
\begin{enumerate}
\item $X$는 준콤팩트이고,
\item $\mathcal{F}$는 유한형이며,
\item $\mathcal{F}'|_{X_s} = \mathcal{F}|_{X_s}$이다.
\end{enumerate}
그러면 $\mathcal{F}$ 위에서 $s^n$을 곱하는 사상이 $\mathcal{F}'$을
통해 분해되게 하는 $n \geq 0$이 존재한다.
\end{lemma}

\begin{proof}
달리 말하면, 어떤 $n \geq 0$에 대해
$s^n\mathcal{F} \subset
\mathcal{F}' \otimes_{\mathcal{O}_X} \mathcal{L}^{\otimes n}$이라고
주장하는 것이다. 다시 말해 몫사상
$\mathcal{F} \to \mathcal{F}/\mathcal{F}'$에 $s$의 어떤 거듭제곱을
곱하면 영이 된다고 주장한다. 이는 『스킴의 성질』의 보조정리
\ref{properties-lemma-section-maps-backwards}에서 따른다.
\end{proof}

\begin{lemma}
\label{obsolete-lemma-bound-degree-in-nbhd-generic-point}
$f : X \to Y$를 스킴의 사상이라 하자. 다음을 가정하자.
\begin{enumerate}
\item $X$와 $Y$는 정수적 스킴이고,
\item $f$는 국소 유한형이고 우세하며,
\item $f$는 준콤팩트이거나 분리이고,
\item $f$는 일반점에서 유한이다. 즉, 『스킴의 사상』의 보조정리
\ref{morphisms-lemma-finite-degree}의 (1)--(5) 가운데 하나가 성립한다.
\end{enumerate}
그러면 $f^{-1}(V) \to V$가 차수 $\deg(X/Y)$인 유한 국소 자유 사상이
되게 하는 공집합이 아닌 열린집합 $V \subset Y$가 존재한다.
특히 $f^{-1}(V) \to V$의 올들의 차수는 $\deg(X/Y)$로 유계이다.
\end{lemma}

\begin{proof}
$f^{-1}(V) \to V$가 유한이 되도록 $V$를 택할 수 있다.
그런 다음 일반 평탄성(『스킴의 사상』의 명제
\ref{morphisms-proposition-generic-flatness})에 의해 $V$를 줄여
$f^{-1}(V) \to V$가 평탄하고 유한 표시라고 가정할 수 있다.
그러면 『스킴의 사상』의 보조정리
\ref{morphisms-lemma-finite-flat}에 의해 이 사상은 유한 국소 자유이다.
$V$가 기약이므로 이 사상은 고정된 차수를 갖는다. 마지막 명제는 이 사실과
『스킴의 사상』의 보조정리
\ref{morphisms-lemma-finite-locally-free-universally-bounded}에서 따른다.
\end{proof}

\section{다양체의 유도 범주}
\label{obsolete-section-equiv}

\noindent
원래 다양체의 유도 범주 장에 속했으나 이제는 필요하지 않은
몇몇 보조정리이다.

\begin{lemma}
\label{obsolete-lemma-preserves-Coh}
$k$를 체라 하자. $X$를 $k$ 위의 정칙인 분리 유한형 스킴이라 하자.
$F : D_{perf}(\mathcal{O}_X) \to D_{perf}(\mathcal{O}_X)$를
$k$-선형 완전 함자라 하자. 모든 연접 $\mathcal{O}_X$-가군
$\mathcal{F}$로서 $\dim(\text{Supp}(\mathcal{F})) = 0$인 것에 대해,
$D_{perf}(\mathcal{O}_X)$의
$M$에 관해 함자적인 $k$-벡터 공간의 동형
$$
\Hom_X(\mathcal{F}, M) = \Hom_X(\mathcal{F}, F(M))
$$
이 있다고 가정하자. 그러면 밑바탕 위상 공간 위에서 항등사상을
유도하는 $k$ 위의 자기동형 $f : X \to X$
\footnote{이 조건은 흔히 $f$가 항등사상이 되도록 강제한다.
『다양체』의 보조정리 \ref{varieties-lemma-automorphism}를 보라.}와
가역 $\mathcal{O}_X$-가군 $\mathcal{L}$이 존재하여,
$F$와 $F'(M) = f^*M \otimes_{\mathcal{O}_X}^\mathbf{L} \mathcal{L}$은
서로 형제이다.
\end{lemma}

\begin{proof}
『다양체의 유도 범주』의 보조정리
\ref{equiv-lemma-duality-at-point}에 의해, 지지가 닫힌 점인 모든 연접
$\mathcal{O}_X$-가군 $\mathcal{F}$에 대해 $M$에 관해 함자적인 동형
$$
H^0(X, M \otimes^\mathbf{L}_{\mathcal{O}_X} \mathcal{F}) =
H^0(X, F(M) \otimes^\mathbf{L}_{\mathcal{O}_X} \mathcal{F})
$$
이 있음을 얻는다.

\medskip\noindent
$x \in X$를 닫힌 점이라 하고, $x$에서 값이 $\kappa(x)$인 마천루층
$\mathcal{F} = \mathcal{O}_x$에 위 사실을 적용하자. 모든
$p \in \mathbf{Z}$에 대해
$$
\dim_{\kappa(x)}
\text{Tor}^{\mathcal{O}_{X, x}}_p(M_x, \kappa(x)) =
\dim_{\kappa(x)}
\text{Tor}^{\mathcal{O}_{X, x}}_p(F(M)_x, \kappa(x))
$$
을 얻는다. 특히 $i > 0$에 대해 $H^i(M) = 0$이면,
『다양체의 유도 범주』의 보조정리
\ref{equiv-lemma-orthogonal-point-sheaf}에 의해 $i > 0$에 대해
$H^i(F(M)) = 0$이다.

\medskip\noindent
$\mathcal{E}$가 계수 $r$인 국소 자유 가군이면 $F(\mathcal{E})$도
계수 $r$인 국소 자유 가군이다. 실제로 다음을 만족하는
$\mathcal{O}_{X, x}$ 위의 완전 복합체 $K$는 차수 $0$에 놓인 계수
$r$인 자유 가군과 같다.
$$
\dim_{\kappa(x)} \text{Tor}^{\mathcal{O}_{X, x}}_i(K, \kappa(x)) =
\left\{
\begin{matrix}
r & \text{조건} & i = 0 \\
0 & \text{조건} & i \not = 0
\end{matrix}
\right.
$$
예를 들어 『대수학 심화』의 보조정리
\ref{more-algebra-lemma-lift-perfect-from-residue-field}를 보라.

\medskip\noindent
$M$이 닫힌 부분스킴 $Z \subset X$ 위에 지지되면 $F(M)$도 $Z$ 위에
지지된다. 실제로 $x \not \in Z$이면
$M \otimes_{\mathcal{O}_X}^\mathbf{L} \mathcal{O}_x = 0$이고, 따라서
$F(M)$에 대해서도 같으므로 『다양체의 유도 범주』의 보조정리
\ref{equiv-lemma-orthogonal-point-sheaf}에서 결론을 얻는다.

\medskip\noindent
특히 $F(\mathcal{O}_x)$는 $\{x\}$에 지지된다.
$H^i(\mathcal{O}_x) \not = 0$인 최소 정수를 $i \in \mathbf{Z}$라 하자.
$i \leq 0$임을 안다. $i < 0$이면 사상
$\mathcal{O}_x[-i] \to F(\mathcal{O}_x)$가 존재하는데, 이는 모든 사상
$\mathcal{O}_x[-i] \to \mathcal{O}_x$가 영이라는 사실과 모순이다.
따라서 $F(\mathcal{O}_x) = \mathcal{H}[0]$이고, 여기서
$\mathcal{H}$는 $x$의 마천루층이다.

\medskip\noindent
$\mathcal{G}$를 $\dim(\text{Supp}(\mathcal{G})) = 0$인 연접
$\mathcal{O}_X$-가군이라 하자. 여과
$$
0 = \mathcal{G}_0 \subset \mathcal{G}_1 \subset \ldots \subset
\mathcal{G}_n = \mathcal{G}
$$
가 존재하여, $n \geq i \geq 1$에 대해 몫
$\mathcal{G}_i/\mathcal{G}_{i - 1}$은 어떤 닫힌 점 $x_i \in X$에 대한
$\mathcal{O}_{x_i}$와 동형이다. 따라서 구별 삼각형
$$
F(\mathcal{G}_{i - 1}) \to F(\mathcal{G}_i) \to F(\mathcal{O}_{x_i})
$$
을 얻고, 귀납법으로 $F(\mathcal{G}_i)$가 차수 $0$에 놓인 연접층임을
알 수 있다.

\medskip\noindent
$\mathcal{G}$를 연접 $\mathcal{O}_X$-가군이라 하자.
$i > 0$에 대해 $H^i(F(\mathcal{G})) = 0$임을 안다. 모순을 얻기 위해
어떤 $i < 0$에 대해 $H^i(F(\mathcal{G}))$가 영이 아니라고 가정하자.
이 성질을 갖는 최소 $i$를 택하면 $D_{perf}(\mathcal{O}_X)$에서 사상
$H^i(F(\mathcal{G}))[-i] \to F(\mathcal{G})$를 얻는다.
$H^i(F(\mathcal{G}))$의 지지에 속하는 닫힌 점 $x \in X$를 택한다.
『대수학 심화』의 보조정리
\ref{more-algebra-lemma-kollar-kovacs-pseudo-coherent}에 의해 다음 사상이
영이 아니게 하는 $n > 0$이 존재한다.
$$
H^i(F(\mathcal{G}))_x \otimes_{\mathcal{O}_{X, x}}
\mathcal{O}_{X, x}/\mathfrak m_x^n
\longrightarrow
\text{Tor}^{\mathcal{O}_{X, x}}_{-i}(F(\mathcal{G})_x,
\mathcal{O}_{X, x}/\mathfrak m_x^n)
$$
다음으로 $m \geq 1$을 택하고 짧은 완전열
$$
0 \to \mathfrak m_x^m \mathcal{G} \to \mathcal{G} \to
\mathcal{G}/\mathfrak m_x^m\mathcal{G} \to 0
$$
을 생각한다. 위에서 $F(\mathcal{G}/\mathfrak m_x^m\mathcal{G})$가
차수 $0$에 놓인 층임을 보았다. 따라서
$H^i(F(\mathfrak m_x^m \mathcal{G})) \to H^i(F(\mathcal{G}))$는
동형이다. 다음 가환 도표를 생각하자.
$$
\xymatrix{
H^i(F(\mathfrak m_x^m\mathcal{G}))_x \otimes_{\mathcal{O}_{X, x}}
\mathcal{O}_{X, x}/\mathfrak m_x^n \ar[r] \ar[d] &
\text{Tor}^{\mathcal{O}_{X, x}}_{-i}(F(\mathfrak m_x^m\mathcal{G})_x,
\mathcal{O}_{X, x}/\mathfrak m_x^n) \ar[d] \\
H^i(F(\mathcal{G}))_x \otimes_{\mathcal{O}_{X, x}}
\mathcal{O}_{X, x}/\mathfrak m_x^n \ar[r] &
\text{Tor}^{\mathcal{O}_{X, x}}_{-i}(F(\mathcal{G})_x,
\mathcal{O}_{X, x}/\mathfrak m_x^n)
}
$$
왼쪽 수직 화살표는 동형이고 아래 화살표는 영이 아니므로, 모든
$m \geq 1$에 대해 오른쪽 수직 화살표는 영이 아니다. 한편 증명의 첫
문단에 의해 이 화살표는 다음 화살표와 동형임을 안다.
$$
\text{Tor}^{\mathcal{O}_{X, x}}_{-i}(\mathfrak m_x^m\mathcal{G}_x,
\mathcal{O}_{X, x}/\mathfrak m_x^n)
\longrightarrow
\text{Tor}^{\mathcal{O}_{X, x}}_{-i}(\mathcal{G}_x,
\mathcal{O}_{X, x}/\mathfrak m_x^n)
$$
그러나 『대수학 심화』의 보조정리
\ref{more-algebra-lemma-tor-annihilated}에 의해 $m \gg n$이면 이 화살표는
영이다. 이것이 원하던 모순이다.

\medskip\noindent
따라서 $F$가 연접 가군을 보존함을 안다. 『다양체의 유도 범주』의
보조정리 \ref{equiv-lemma-exact-functor-preserving-Coh}에 의해,
$F$는 다음과 같은 연접 $\mathcal{O}_{X \times X}$-가군
$\mathcal{K}$가 주는 Fourier--Mukai 함자 $F'$과 형제이다.
이 가군은 $\text{pr}_1$을 통해 $X$ 위에서 평탄하고
$\text{pr}_2$를 통해 $X$ 위에서 유한이다.
$F(\mathcal{O}_X)$가 차수 $0$에 놓인 가역 $\mathcal{O}_X$-가군
$\mathcal{L}$이므로
$$
\mathcal{L} \cong F(\mathcal{O}_X) \cong F'(\mathcal{O}_X) \cong
\text{pr}_{2, *}\mathcal{K}
$$
이다. 그러므로 『함자와 사상』의 보조정리
\ref{functors-lemma-pushforward-invertible-pre}에 의해
$\text{pr}_2 \circ s = \text{id}_X$이고
$\mathcal{K} = s_*\mathcal{L}$인 사상 $s : X \to X \times X$가
존재한다. $f = \text{pr}_1 \circ s$로 놓는다. 그러면
\begin{align*}
F'(M)
& =
R\text{pr}_{2, *}(L\text{pr}_1^*K \otimes \mathcal{K}) \\
& =
R\text{pr}_{2, *}(L\text{pr}_1^*M \otimes s_*\mathcal{L}) \\
& =
R\text{pr}_{2, *}(Rs_*(Lf^*M \otimes \mathcal{L})) \\
& =
Lf^*M \otimes \mathcal{L}
\end{align*}
이다. 세 번째 단계에서 『스킴의 유도 범주』의 보조정리
\ref{perfect-lemma-cohomology-base-change}를 사용했다.
모든 닫힌 점 $x \in X$에 대해 가군 $F(\mathcal{O}_x)$가 $x$에
지지되므로, $f$는 $X$의 밑바탕 위상 공간 위에서 항등사상을 유도한다.
이제 $f$가 동형임을 보여야 하며, 다음 문단에서 이를 하겠다.

\medskip\noindent
$x \in X$를 닫힌 점이라 하자. $n \geq 1$에 대해 $x$에서 값이
$\mathcal{O}_{X, x}/\mathfrak m_x^n$인 마천루층을
$\mathcal{O}_{x, n}$으로 나타내자. $\mathcal{O}_{x, n}$ 사이의
$\mathcal{O}_X$-가군 준동형에 관해 함자적인 동형
$$
\Hom_X(\mathcal{O}_{x, m}, \mathcal{O}_{x, n}) \cong
\Hom_X(\mathcal{O}_{x, m}, F(\mathcal{O}_{x, n})) \cong
\Hom_X(\mathcal{O}_{x, m}, f^*\mathcal{O}_{x, n} \otimes \mathcal{L})
$$
이 있다. (첫 번째 동형은 가정에 의해 존재하고, 두 번째 동형은 $F$와
$F'$이 형제이기 때문에 존재한다.) $m \geq n$이면
$\mathcal{O}_{x, n}$ 위의 작용을 통해
$\mathcal{O}_{X, x}/\mathfrak m^n =
\Hom_X(\mathcal{O}_{x, m}, \mathcal{O}_{x, n})$이므로,
모든 $n$에 대해
$f^\sharp : \mathcal{O}_{X, x}/\mathfrak m_x^n \to
\mathcal{O}_{X, x}/\mathfrak m_x^n$이 전단사임을 얻는다.
따라서 $f$는 닫힌 점에서의 완비 국소환들 사이의 동형을 유도하고,
그러므로 에탈이다(『에탈 사상』의 보조정리
\ref{etale-lemma-characterize-etale-completions}).
닫힌 점들을 살펴보면
$\Delta_f : X \to X \times_{f, X, f} X$는 ($f$가 에탈이므로 열린
몰입이고) 전단사이므로 동형이다. 따라서 $f$는 단사 사상이다.
마지막으로 『내림』의 보조정리
\ref{descent-lemma-flat-surjective-quasi-compact-monomorphism-isomorphism}가
이 사상이 열린 몰입임을 말해 주므로, $f$는 동형이라고 결론짓는다.
\end{proof}

\section{정칙 고유 경우의 표현가능성}
\label{obsolete-section-regular-proper}

\noindent
『다양체의 유도 범주』의 정리
\ref{equiv-theorem-bondal-van-den-bergh}를 개선하여 체 위의 모든 고유
스킴에 적용되게 했으므로 이 절은 폐기되었다(이전에는 체 위의 사영
스킴에 대해서만 증명했다).

\begin{lemma}
\label{obsolete-lemma-trace-map}
\begin{reference}
여기 제시한 증명은 \cite[주석 3.4]{MS}의 논증을 따른다.
\end{reference}
$f : X' \to X$를 정수적 뇌터 스킴의 고유 쌍유리 사상이라 하고
$X$가 정칙이라고 하자. 사상
$\mathcal{O}_X \to Rf_*\mathcal{O}_{X'}$은 $D(\mathcal{O}_X)$에서
표준적으로 분할된다.
\end{lemma}

\begin{proof}
$D(\mathcal{O}_X)$에서 $E = Rf_*\mathcal{O}_{X'}$로 놓는다.
『스킴의 유도 범주』의 보조정리
\ref{perfect-lemma-direct-image-coherent}에 의해 $E$는
$D^b_{\textit{Coh}}(\mathcal{O}_X)$에 속함에 유의하자.
『스킴의 유도 범주』의 보조정리
\ref{perfect-lemma-perfect-on-regular}에 의해 $E$는
$D(\mathcal{O}_X)$의 완전 대상이다.
$\mathcal{O}_{X'}$가 대수의 층이므로 『코호몰로지』의 주석
\ref{cohomology-remark-cup-product}에 의해 상대 컵곱
$\mu : E \otimes_{\mathcal{O}_X}^\mathbf{L} E \to E$가 있다.
$\sigma : E \otimes E^\vee \to E^\vee \otimes E$를 대칭 모노이드 범주
$D(\mathcal{O}_X)$ 위의 가환성 제약이라 하자
(『코호몰로지』의 보조정리
\ref{cohomology-lemma-symmetric-monoidal-derived}).
$\eta : \mathcal{O}_X \to E \otimes E^\vee$와
$\epsilon : E^\vee \otimes E \to \mathcal{O}_X$를 『코호몰로지』의 예
\ref{cohomology-example-dual-derived}에서 구성한 사상들이라 하자.
그러면 다음 사상을 생각할 수 있다.
$$
E \xrightarrow{\eta \otimes 1} E \otimes E^\vee \otimes E
\xrightarrow{\sigma \otimes 1} E^\vee \otimes E \otimes E
\xrightarrow{1 \otimes \mu} E^\vee \otimes E
\xrightarrow{\epsilon} \mathcal{O}_X
$$
이 사상이 보조정리의 사상의 한쪽 역이라고 주장한다. 이를 보이려면
합성 $\mathcal{O}_X \to \mathcal{O}_X$가 항등사상임을 보이면 충분하다.
이는 $X$의 일반점에서(또는 $f$가 동형인 $X$의 열린 부분스킴 위에서)
확인해도 된다. 이 경우 $E = \mathcal{O}_X$이고 $\mu$는 보통의 곱셈
사상이므로 결과가 명백하다.
\end{proof}

\begin{lemma}
\label{obsolete-lemma-characterize-dbcoh-proper-regular}
$X$를 체 $k$ 위의 정칙 고유 스킴이라 하고
$K \in \Ob(D_\QCoh(\mathcal{O}_X))$라 하자. 다음은 동치이다.
\begin{enumerate}
\item $K \in D^b_{\textit{Coh}}(\mathcal{O}_X) =
D_{perf}(\mathcal{O}_X)$이고,
\item $D(\mathcal{O}_X)$의 모든 완전 대상 $E$에 대해
$\sum_{i \in \mathbf{Z}} \dim_k \Ext^i_X(E, K) < \infty$이다.
\end{enumerate}
\end{lemma}

\begin{proof}
(1)의 등식은 『스킴의 유도 범주』의 보조정리
\ref{perfect-lemma-perfect-on-regular}에 의해 성립한다.
(1) $\Rightarrow$ (2)는 『다양체의 유도 범주』의 보조정리
\ref{equiv-lemma-finiteness}에서 따른다.
(2) $\Rightarrow$ (1)는 『사상 심화』의 보조정리
\ref{more-morphisms-lemma-characterize-relatively-perfect}에서 따른다.
\end{proof}

\begin{lemma}
\label{obsolete-lemma-bondal-van-den-bergh}
$X$를 체 $k$ 위의 정칙 고유 스킴이라 하자.
\begin{enumerate}
\item $F : D_{perf}(\mathcal{O}_X)^{opp} \to \text{Vect}_k$를 모든
$E \in D_{perf}(\mathcal{O}_X)$에 대해
$$
\sum\nolimits_{n \in \mathbf{Z}} \dim_k F(E[n]) < \infty
$$
를 만족하는 $k$-선형 코호몰로지적 함자라 하자. 그러면 $F$는 어떤
$K \in D_{perf}(\mathcal{O}_X)$에 대한 함자
$E \mapsto \Hom_X(E, K)$와 동형이다.
\item $G : D_{perf}(\mathcal{O}_X) \to \text{Vect}_k$를 모든
$E \in D_{perf}(\mathcal{O}_X)$에 대해
$$
\sum\nolimits_{n \in \mathbf{Z}} \dim_k G(E[n]) < \infty
$$
를 만족하는 $k$-선형 호몰로지적 함자라 하자. 그러면 $G$는 어떤
$K \in D_{perf}(\mathcal{O}_X)$에 대한 함자
$E \mapsto \Hom_X(K, E)$와 동형이다.
\end{enumerate}
\end{lemma}

\begin{proof}
이는 『다양체의 유도 범주』의 정리
\ref{equiv-theorem-bondal-van-den-bergh}와 보조정리
\ref{equiv-lemma-homological-representable}에서 따른다.
아래에 다른 증명도 제시한다.

\medskip\noindent
(1)의 증명. 유도 범주 $D_\QCoh(\mathcal{O}_X)$는 직합을 가지며,
콤팩트 생성되고, $D_{perf}(\mathcal{O}_X)$는 콤팩트 대상들의 충만한
부분범주이다. 『스킴의 유도 범주』의 보조정리
\ref{perfect-lemma-quasi-coherence-direct-sums}, 정리
\ref{perfect-theorem-bondal-van-den-Bergh}, 명제
\ref{perfect-proposition-compact-is-perfect}를 보라.
『다양체의 유도 범주』의 보조정리
\ref{equiv-lemma-van-den-bergh}에 의해 어떤
$K \in \Ob(D_\QCoh(\mathcal{O}_X))$에 대해
$F(E) = \Hom_X(E, K)$라고 가정할 수 있다. 그러면 보조정리
\ref{obsolete-lemma-characterize-dbcoh-proper-regular}에 의해 $K$는
$D^b_{\textit{Coh}}(\mathcal{O}_X)$에 속한다.

\medskip\noindent
(2)의 증명. $D_{perf}(\mathcal{O}_X)$ 위의 반변 함자
$E \mapsto E^\vee$를 생각하자. 『코호몰로지』의 보조정리
\ref{cohomology-lemma-dual-perfect-complex}를 보라.
이 함자는 $D_{perf}(\mathcal{O}_X)$의 완전 반자기동치이다.
따라서 함자 $F(E) = G(E^\vee)$에 (1)을 적용하면
$G(E^\vee) = \Hom_X(E, K)$인 $K \in D_{perf}(\mathcal{O}_X)$를 얻는다.
그러면 $G(E) = \Hom_X(E^\vee, K) = \Hom_X(K^\vee, E)$이므로
$K^\vee$를 택하면 된다.
\end{proof}

\section{몫의 함자}
\label{obsolete-section-quotients}

\begin{lemma}
\label{obsolete-lemma-factors-through-quotient}
$S = \Spec(R)$을 아핀 스킴이라 하고 $X$를 $S$ 위의 대수공간이라 하자.
$q_i : \mathcal{F} \to \mathcal{Q}_i$, $i = 1, 2$를 준연접
$\mathcal{O}_X$-가군의 전사라 하자. $\mathcal{Q}_1$이 $S$ 위에서
평탄하다고 가정하자. $T \to S$를 다음 분해가 존재하는 스킴의
준콤팩트 사상이라 하자.
$$
\xymatrix{
& \mathcal{F}_T \ar[rd]^{q_{2, T}} \ar[ld]_{q_{1, T}} \\
\mathcal{Q}_{1, T} & & \mathcal{Q}_{2, T} \ar@{..>}[ll]
}
$$
그러면 다음을 만족하는 닫힌 부분스킴 $Z \subset S$가 존재한다.
(a) $T \to S$는 $Z$를 통해 분해되고, (b) $q_{1, Z}$는
$q_{2, Z}$를 통해 분해된다. $\Ker(q_2)$가 유한형
$\mathcal{O}_X$-가군이고 $X$가 준콤팩트이면 $Z \to S$를 유한 표시로
택할 수 있다.
\end{lemma}

\begin{proof}
사상 $\Ker(q_2) \to \mathcal{Q}_1$에 『공간 위의 평탄성』의 보조정리
\ref{spaces-flat-lemma-F-zero-somewhat-closed}를 적용한다.
\end{proof}

\section{공간과 fpqc 피복}
\label{obsolete-section-fpqc}

\noindent
대수공간이 fpqc 피복에 관한 층 조건을 만족한다는 Gabber의 논증으로
인해 여기의 내용은 폐기되었다. 『공간의 성질』의 절
\ref{spaces-properties-section-fpqc}를 보라.

\begin{lemma}
\label{obsolete-lemma-separated-fpqc}
$S$를 스킴이라 하고 $X$를 $S$ 위의 대수공간이라 하자.
$\{f_i : T_i \to T\}_{i \in I}$를 $S$ 위의 스킴의 fpqc 피복이라 하자.
그러면 사상
$$
\Mor_S(T, X)
\longrightarrow
\prod\nolimits_{i \in I} \Mor_S(T_i, X)
$$
은 단사이다.
\end{lemma}

\begin{proof}
『공간의 성질』의 명제
\ref{spaces-properties-proposition-sheaf-fpqc}에서 즉시 따른다.
\end{proof}

\begin{lemma}
\label{obsolete-lemma-sheaf-fpqc-open-covering}
$S$를 스킴이라 하고 $X$를 $S$ 위의 대수공간이라 하자.
$X = \bigcup_{j \in J} X_j$를 Zariski 피복이라 하자. 『공간』의 정의
\ref{spaces-definition-Zariski-open-covering}을 보라.
각 $X_j$가 fpqc 위상에 대한 층 성질을 만족하면 $X$도 fpqc 위상에 대한
층 성질을 만족한다.
\end{lemma}

\begin{proof}
모든 대수공간이 fpqc 위상에 대한 층 성질을 만족하므로 참이다.
『공간의 성질』의 명제
\ref{spaces-properties-proposition-sheaf-fpqc}를 보라.
\end{proof}

\begin{lemma}
\label{obsolete-lemma-sheaf-fpqc-quasi-separated}
$S$를 스킴이라 하고 $X$를 $S$ 위의 대수공간이라 하자.
$X$가 $S$ 위에서 Zariski 국소적으로 준분리이면, $X$는 fpqc 위상에
대한 층 조건을 만족한다.
\end{lemma}

\begin{proof}
일반적인 『공간의 성질』의 명제
\ref{spaces-properties-proposition-sheaf-fpqc}에서 즉시 따른다.
\end{proof}

\begin{remark}
\label{obsolete-remark-proof-works-when}
이 주석은 전에는 보조정리 \ref{obsolete-lemma-sheaf-fpqc-quasi-separated}의
원래 증명(2009년 12월 18일)이 어느 정도까지 일반화되는지를 논의했다.
\end{remark}

\section{매우 합리적인 대수공간}
\label{obsolete-section-very-reasonable}

\noindent
다소 폐기된 내용이다.

\begin{lemma}
\label{obsolete-lemma-reasonable-kolmogorov}
$S$를 스킴이라 하고 $X$를 $S$ 위의 합리적인 대수공간이라 하자.
그러면 $|X|$는 Kolmogorov 공간이다(『위상 공간』의 정의
\ref{topology-definition-generic-point}을 보라).
\end{lemma}

\begin{proof}
정의들과 『괜찮은 공간』의 보조정리
\ref{decent-spaces-lemma-kolmogorov}에서 따른다.
\end{proof}

\noindent
이 절의 나머지에서는 매우 합리적인 대수공간에 관해 몇 가지 논평을 한다.
스킴 $U$와 전사인 에탈 준콤팩트 사상 $U \to X$가 존재하면 $X$는 매우
합리적이다. 『괜찮은 공간』의 보조정리
\ref{decent-spaces-lemma-characterize-very-reasonable}를 보라.

\begin{lemma}
\label{obsolete-lemma-scheme-very-reasonable}
모든 스킴은 매우 합리적이다.
\end{lemma}

\begin{proof}
항등사상이 준콤팩트인 전사 에탈 사상이므로 참이다.
\end{proof}

\begin{lemma}
\label{obsolete-lemma-very-reasonable-Zariski-local}
$S$를 스킴이라 하고 $X$를 $S$ 위의 대수공간이라 하자.
각 $X_i$가 매우 합리적인 Zariski 열린 피복 $X = \bigcup X_i$가
존재하면 $X$는 매우 합리적이다.
\end{lemma}

\begin{proof}
『괜찮은 공간』의 보조정리
\ref{decent-spaces-lemma-properties-local}의 경우 $(\epsilon)$이다.
\end{proof}

\begin{lemma}
\label{obsolete-lemma-quasi-separated-very-reasonable}
Zariski 국소적으로 준분리인 대수공간은 매우 합리적이다.
특히 모든 준분리 대수공간은 매우 합리적이다.
\end{lemma}

\begin{proof}
『괜찮은 공간』의 보조정리
\ref{decent-spaces-lemma-bounded-fibres}에 나오는 함의 가운데 하나이다.
\end{proof}

\begin{lemma}
\label{obsolete-lemma-representable-very-reasonable}
$S$를 스킴이라 하고 $X$, $Y$를 $S$ 위의 대수공간이라 하자.
$Y \to X$를 표현가능한 사상이라 하자. $X$가 매우 합리적이면
$Y$도 매우 합리적이다.
\end{lemma}

\begin{proof}
『괜찮은 공간』의 보조정리
\ref{decent-spaces-lemma-representable-properties}의 경우
$(\epsilon)$이다.
\end{proof}

\begin{remark}
\label{obsolete-remark-very-reasonable-Zariski-locally-quasi-separated}
매우 합리적인 대수공간들은 Zariski 국소적으로 준분리인 대수공간들보다
엄밀히 더 큰 모임을 이룬다. 고정점 없이 비준분리 스킴 $U$ 위에 작용하는
유한군 $G$에 대해 $X = [U/G]$ 꼴의 대수공간을 생각하자
(『공간』의 정의 \ref{spaces-definition-quotient}을 보라).
실제로 이 경우 $U \times_X U = U \times G$이고 $U$로 가는 두 사영은
명백히 준콤팩트이므로 $X$는 매우 합리적이다. 한편 대각사상
$U \times_X U \to U \times U$는 준콤팩트가 아니므로 이 대수공간은
준분리가 아니다. 이제 표수가 $\not = 2$인 체 $k$ 위의 무한 차원 아핀
공간에서 영점을 둘로 만든 공간을 $U$로 택하자. 『스킴』의 예
\ref{schemes-example-not-quasi-separated}를 보라.
$0_1, 0_2$를 $U$의 두 영점이라 하자. $G = \{+1, -1\}$로 놓고,
$-1$이 모든 좌표에는 $-1$로 작용하며 $0_1$과 $0_2$를 맞바꾸게 하자.
그러면 $[U/G]$는 매우 합리적이지만 Zariski 국소적으로 준분리이지는 않다
(세부 사항은 생략한다).
\end{remark}

\noindent
경고: 다음 보조정리는 주의해서 사용해야 한다. 여기에 나오는 스킴
$U_i$들은 반드시 분리이거나 심지어 준분리일 필요도 없다.

\begin{lemma}
\label{obsolete-lemma-very-reasonable-quasi-compact-pieces}
$S$를 스킴이라 하고 $X$를 $S$ 위의 매우 합리적인 대수공간이라 하자.
다음을 만족하는 스킴들의 집합 $U_i$와 사상 $U_i \to X$가 존재한다.
\begin{enumerate}
\item 각 $U_i$는 준콤팩트 스킴이고,
\item 각 $U_i \to X$는 에탈이며,
\item 두 사영 $U_i \times_X U_i \to U_i$는 모두 준콤팩트이고,
\item 사상 $\coprod U_i \to X$는 전사이다(그리고 에탈이다).
\end{enumerate}
\end{lemma}

\begin{proof}
『괜찮은 공간』의 정의
\ref{decent-spaces-definition-very-reasonable}는 (2), (3), (4)를 만족하는
$U_i \to X$가 존재한다고 말한다. $i$를 고정하고
$R_i = U_i \times_X U_i$로 놓으며 사영들을
$s, t : R_i \to U_i$로 나타내자. 임의의 아핀 열린집합
$W \subset U_i$에 대해 열린집합 $W' = t(s^{-1}(W)) \subset U_i$는
준콤팩트이고 $R_i$-불변이다(『군체』의 보조정리
\ref{groupoids-lemma-constructing-invariant-opens}를 보라).
따라서 $W'$은 준콤팩트 스킴이고 $W' \to X$는 에탈이며,
$W' \times_X W' = s^{-1}(W') = t^{-1}(W')$이므로 두 사영
$W' \times_X W' \to W'$은 준콤팩트이다. 즉, $W \subset U_i$가
$U_i$의 아핀 열린 피복들의 원소들을 달릴 때 얻는 $W' \to X$의 족이
원하는 것을 준다.
\end{proof}

\section{대수공간에 관한 폐기된 보조정리}
\label{obsolete-section-obsolete-on-spaces}

\noindent
불필요해 보이거나 본문에서 더 이상 쓰이지 않는 보조정리들이다.

\begin{lemma}
\label{obsolete-lemma-vanishing-surjective}
『공간의 코호몰로지』의 상황
\ref{spaces-cohomology-situation-vanishing}에서 사상
$p : X \to \Spec(A)$는 전사이다.
\end{lemma}

\begin{proof}
이 보조정리는 원래 『공간의 코호몰로지』의 명제
\ref{spaces-cohomology-proposition-vanishing-affine}의 증명에 사용되었으나,
이제는 그 명제의 결과이다.
\end{proof}

\begin{lemma}
\label{obsolete-lemma-vanishing-universally-closed}
『공간의 코호몰로지』의 상황
\ref{spaces-cohomology-situation-vanishing}에서 사상
$p : X \to \Spec(A)$는 보편적으로 닫혀 있다.
\end{lemma}

\begin{proof}
이 보조정리는 원래 『공간의 코호몰로지』의 명제
\ref{spaces-cohomology-proposition-vanishing-affine}의 증명에 사용되었으나,
이제는 그 명제의 결과이다.
\end{proof}

\begin{remark}
\label{obsolete-remark-equation-first}
이 태그는 전에는 『형식 공간』의 보조정리
\ref{formal-spaces-lemma-iff-adic-star}의 증명에 있는 한 등식을 가리켰다.
\end{remark}

\begin{remark}
\label{obsolete-remark-equation-second}
이 태그는 전에는 『형식 공간』의 보조정리
\ref{formal-spaces-lemma-iff-adic-star}의 증명에 있는 한 등식을 가리켰다.
\end{remark}




\section{대수적 스택에 관한 폐기된 보조정리}
\label{obsolete-section-obsolete-on-stacks}

\noindent
불필요해 보이거나 본문에서 더 이상 쓰이지 않는 보조정리들이다.

\begin{lemma}
\label{obsolete-lemma-infinite-sequence}
$S$를 국소 뇌터 스킴이라 하자. $\mathcal{X}$를 (RS*)를 갖는
$(\Sch/S)_{fppf}$ 위의 군체로 올화된 범주라 하자.
$x$를 $S$ 위에서 유한형인 아핀 스킴 $U$ 위의 $\mathcal{X}$의
대상이라 하자. $u_n \in U$, $n \geq 1$을 서로 다르고 유한형인 점들이라
하고, 모든 $n$에 대해 $x$가 $u_n$에서 버설이 아니라고 하자.
$u_n$을 부분열로 바꾼 뒤, $S$ 위에서
$$
x \to x_1 \to x_2 \to \ldots
\quad\text{범주 }\mathcal{X}\text{ 안에서 다음 열 위에 놓인}\quad
U \to U_1 \to U_2 \to \ldots
$$
인 사상들이 존재하여 다음을 만족한다.
\begin{enumerate}
\item 각 $n$에 대해 사상 $U \to U_n$은 일차 비후이고,
\item 각 $n$에 대해 짧은 완전열
$$
0 \to \kappa(u_n) \to \mathcal{O}_{U_n} \to \mathcal{O}_{U_{n - 1}} \to 0
$$
이 있으며 $n = 1$일 때 $U_0 = U$이고,
\item 각 $n$에 대해 $u_n$의 열린 근방 $W \subset U_n$과 사상
$\alpha : x_n|_W \to x$로 이루어진 쌍 $(W, \alpha)$ 가운데,
합성
$$
x|_{U \cap W} \xrightarrow{\text{사상의 제한 }x \to x_n}
x_n|_W \xrightarrow{\alpha} x
$$
이 표준 사상 $x|_{U \cap W} \to x$가 되는 것은 {\bf 존재하지
않는다}.
\end{enumerate}
\end{lemma}

\begin{proof}
이 보조정리는 원래 버설성의 개방성 판정법
(『Artin의 공리』의 보조정리
\ref{artin-lemma-SGE-implies-openness-versality})의 증명에 사용되었으나,
『Artin의 공리』의 보조정리
\ref{artin-lemma-infinite-sequence-pre}로 대체되었고 여기서 쉽게 따른다.
실제로 $u_n$, $n \geq 1$을 부분열로 바꾼 뒤 이 점들 사이에는 특수화가
없다고 가정할 수 있고 그렇게 한다. 『스킴의 성질』의 보조정리
\ref{properties-lemma-thin-infinite-sequence}를 보라. 이제
『Artin의 공리』의 보조정리
\ref{artin-lemma-infinite-sequence-pre}를 적용하면 증명이 끝난다.
\end{proof}

\section{스킴의 여접 복합체의 변형}
\label{obsolete-section-cotangent-schemes-variant}

\noindent
이 절에서는 스킴의 사상의 여접 복합체를 구성하는 다른 방법을 제시한다.
응용에는 『여접 복합체』의 절
\ref{cotangent-section-cotangent-schemes-variant}에서 논의한 더 쉬운 판으로
충분하므로, 이 절은 현재 폐기된 결과 장에 들어 있다.

\medskip\noindent
$f : X \to Y$를 스킴의 사상이라 하자. $\mathcal{C}_{X/Y}$를 다음과 같은
스킴의 가환 도표를 대상으로 갖는 범주라 하자.
\begin{equation}
\label{obsolete-equation-object}
\vcenter{
\xymatrix{
X \ar[d] & U \ar[l] \ar[d] \ar[r]_i & A \ar[ld] \\
Y & V \ar[l]
}
}
\end{equation}
여기서
\begin{enumerate}
\item $U$는 $X$의 열린 부분스킴이고,
\item $V$는 $Y$의 열린 부분스킴이며,
\item $P$가 $\mathbf{Z}$ 위의 (어떤 변수 집합에 대한) 다항식 대수일 때
$V$ 위의 동형 $A = V \times \Spec(P)$가 존재한다.
\end{enumerate}
즉, $A$는 $V$ 위의 (무한 차원) 아핀 공간이다. 사상은 가환 도표로
주어진다.

\medskip\noindent
{\bf 표기법.} $\mathcal{C}_{X/Y}$의 대상, 즉 도표
(\ref{obsolete-equation-object})는 흔히 $U \to A$로 나타낸다. 이때 다음을
암묵적으로 이해한다. (a) $U$는 $X$의 열린 부분스킴이고, (b)
$U \to A$는 $Y$ 위의 사상이며, (c) 구조 사상 $A \to Y$의 상은 열린집합
$V \subset Y$이고, (d) $A \to V$는 아핀 공간이다.
$V \subset Y$가 $A \to Y$의 상임을 나타내기 위해 $U \to A/V$라고
쓰겠다. $X_{Zar}$가 $X$의 작은 Zariski 사이트를 나타냄을 상기하자.
망각 함자
$$
\mathcal{C}_{X/Y} \to X_{Zar},\ (U \to A) \mapsto U
\quad\text{및}\quad
\mathcal{C}_{X/Y} \mapsto Y_{Zar},\ (U \to A/V) \mapsto V.
$$
가 있다.

\begin{lemma}
\label{obsolete-lemma-category-fibred}
$X \to Y$를 스킴의 사상이라 하자.
\begin{enumerate}
\item 범주 $\mathcal{C}_{X/Y}$는 $X_{Zar}$ 위에 올화되어 있다.
\item 범주 $\mathcal{C}_{X/Y}$는 $Y_{Zar}$ 위에 올화되어 있다.
\item 범주 $\mathcal{C}_{X/Y}$는 $U \subset X$, $V \subset Y$가
열린집합이고 $f(U) \subset V$인 쌍 $(U, V)$들의 범주 위에 올화되어 있다.
\end{enumerate}
\end{lemma}

\begin{proof}
(1)에 대하여. $\mathcal{C}_{X/Y}$의 대상 $U \to A$와
$X_{Zar}$의 사상 $U' \to U$가 주어졌다고 하자. $i'$을 $i$와
$U' \to U$의 합성으로 두어 $\mathcal{C}_{X/Y}$의 대상
$i' : U' \to A$를 생각한다. $\mathcal{C}_{X/Y}$의 사상
$(U' \to A) \to (U \to A)$는 $X_{Zar}$ 위에서 강한 Cartesian이다.

\medskip\noindent
(2)에 대하여. 대상 $U \to A/V$와 $V' \to V$가 주어졌을 때
$U' = U \cap f^{-1}(V')$, $A' = V' \times_V A$로 놓으면,
$V' \to V$ 위의 강한 Cartesian 사상
$(U' \to A') \to (U \to A)$를 얻는다.

\medskip\noindent
(3)에 대하여. (3)의 범주를 $(X/Y)_{Zar}$로 나타내자.
$U \to A/V$와 $(X/Y)_{Zar}$의 사상
$(U', V') \to (U, V)$가 주어졌을 때 $A' = V' \times_V A$를
생각할 수 있다. 그러면 사상
$(U' \to A'/V') \to (U \to A/V)$는 $(X/Y)_{Zar}$ 위에서
$\mathcal{C}_{X/Y}$의 강한 Cartesian 사상이다.
\end{proof}

\noindent
$X_{Zar}$에서 물려받은 위상을 사용하여 $\mathcal{C}_{X/Y}$ 위의 위상
$\tau_X$를 얻는다(『스택』의 절 \ref{stacks-section-topology}를 보라).
달리 명시하지 않는 한 이것이 우리가 $\mathcal{C}_{X/Y}$ 위에서
고려하는 위상이다. 정확히 말해, 사상들의 족
$\{(U_i \to A_i) \to (U \to A)\}$가 $\mathcal{C}_{X/Y}$의 피복일
필요충분조건은 다음과 같다.
\begin{enumerate}
\item $U = \bigcup U_i$이고,
\item 모든 $i$에 대해 $A_i = A$이다.
\end{enumerate}
(2)에서 $A_i \cong A$를 허용해도 같은 층들의 모임을 얻는다.
함자 $u$는 토포스의 사상
$\pi : \Sh(\mathcal{C}_{X/Y}) \to \Sh(X_{Zar})$를 정의한다.

\medskip\noindent
사이트 $\mathcal{C}_{X/Y}$에는 다음과 같은 여러 환의 층이 있다.
\begin{enumerate}
\item 규칙 $(U \to A) \mapsto \mathcal{O}(A)$로 주어지는 층
$\mathcal{O}$.
\item 규칙 $(U \to A) \mapsto \mathcal{O}(U)$로 주어지는 층
$\underline{\mathcal{O}}_X = \pi^{-1}\mathcal{O}_X$.
\item 규칙 $(U \to A/V) \mapsto \mathcal{O}(V)$로 주어지는 층
$\underline{\mathcal{O}}_Y$.
\end{enumerate}
환 달린 토포스의 사상
\begin{equation}
\label{obsolete-equation-pi-schemes}
\vcenter{
\xymatrix{
(\Sh(\mathcal{C}_{X/Y}), \underline{\mathcal{O}}_X) \ar[r]_i \ar[d]_\pi &
(\Sh(\mathcal{C}_{X/Y}), \mathcal{O}) \\
(\Sh(X_{Zar}), \mathcal{O}_X)
}
}
\end{equation}
을 얻는다. 사상 $i$는 밑바탕 토포스 위에서 항등이고,
$i^\sharp : \mathcal{O} \to \underline{\mathcal{O}}_X$는 명백한 사상이다.
사상 $\pi$는 『사이트 위의 코호몰로지』의 상황
\ref{sites-cohomology-situation-fibred-category}의 특수한 경우이다.
다음에서는 유도 함자
$
Li^* : D(\mathcal{O}) \longrightarrow D(\underline{\mathcal{O}}_X)
$,
$Ri_* = i_* : D(\underline{\mathcal{O}}_X) \to D(\mathcal{O})$의 왼쪽
수반과,
$
L\pi_! : D(\underline{\mathcal{O}}_X) \longrightarrow D(\mathcal{O}_X)
$,
$\pi^* = \pi^{-1} : D(\mathcal{O}_X) \to
D(\underline{\mathcal{O}}_X)$의 왼쪽 수반이 중요한 역할을 한다.

\begin{remark}
\label{obsolete-remark-different-topologies}
$Y_{Zar}$에서 물려받은 위상을 택하여 $\mathcal{C}_{X/Y}$ 위의 두 번째
위상 $\tau_Y$를 얻는다. 사상들의 족
$\{(U_i \to A_i) \to (U \to A)\}$가 피복일 필요충분조건이
$U = \bigcup U_i$, $V = \bigcup V_i$, 그리고
$A_i \cong V_i \times_V A$인 세 번째 위상 $\tau_{X \to Y}$도 있다.
이는 보조정리 \ref{obsolete-lemma-category-fibred}의 (3)에 나오는 쌍 $(U, V)$들의
범주를 밑바탕 범주로 하는 사이트 $(X/Y)_{Zar}$의 위상에서 물려받은
위상이다. $(X/Y)_{Zar}$의 피복들은
$U = \bigcup U_i$와 $V = \bigcup V_i$를 만족하는 족
$\{(U_i, V_i) \to (U, V)\}$이다. 토포스의 사상
$$
\xymatrix{
\Sh(\mathcal{C}_{X/Y})
= \Sh(\mathcal{C}_{X/Y}, \tau_X) &
\Sh(\mathcal{C}_{X/Y}, \tau_{X \to Y}) \ar[l] \ar[r] &
\Sh(\mathcal{C}_{X/Y}, \tau_Y)
}
$$
이 있다($\tau_X$가 우리의 ``기본'' 위상임을 상기하자).
이 화살표들의 당김뒤 함자는 층화이고, 밂앞은 밑바탕 준층 위에서
항등이다. 토포스의 도표
$$
\xymatrix{
\Sh(X_{Zar}) \ar[d]^f & \Sh(\mathcal{C}_{X/Y}) \ar[l]^\pi &
\Sh(\mathcal{C}_{X/Y}, \tau_{X \to Y}) \ar[l] \ar[d] \\
\Sh(Y_{Zar}) & & \Sh(\mathcal{C}_{X/Y}, \tau_Y) \ar[ll]
}
$$
는 {\bf 가환하지 않는다}. 실제로 $Y$ 위의 영이 아닌 아벨 층을
$(\mathcal{C}_{X/Y}, \tau_{X \to Y})$로 당기면 영이 아닌 아벨 층이
되지만, 그런 층을 $X$로 당기면 영이 되는 예를 분명히 찾을 수 있다.
$Y_{Zar}$ 위의 임의의 준층 $\mathcal{F}$는 객체 $(U \to A/V)$에
집합 $\mathcal{F}(V)$를 대응시키는 규칙에 의해 $\mathcal{C}_{Y/X}$ 위의
층 $\underline{\mathcal{F}}$를 준다는 점에 유의하자.
$\mathcal{F}$가 우연히 층이더라도 일반적으로
$\underline{\mathcal{F}} = \pi^{-1}f^{-1}\mathcal{F}$인 것은 아니다.
이는 위 도표의 비가환성과 관련된다. 실제로 아래 수평 화살표와 오른쪽
수직 화살표를 통해 $\mathcal{F}$를
$\Sh(\mathcal{C}_{X/Y}, \tau_{X \to Y})$로 당긴 뒤 밂앞한 것으로
$\underline{\mathcal{F}}$를 기술할 수 있다. 한 예는 층
$\underline{\mathcal{O}}_Y$이다. 그러나 참인 사실은 사상
$\underline{\mathcal{F}} \to \pi^{-1}f^{-1}\mathcal{F}$가 있고,
(나중에 보겠지만) $\pi_!$를 적용하면 이 사상이 동형으로 바뀐다는 것이다.
\end{remark}

\section{평탄 가군의 변형과 장애물}
\label{obsolete-section-modules}

\noindent
이 절에서는 환 $\Lambda$ 위의 임의의 대수공간 $X$에 대한 연접층의
스택의 변형 이론 구성을 개략적으로 설명한다. 『Quot』의 절
\ref{quot-section-not-flat}에서 논의를 개선했으므로 이 내용은 폐기되었다.

\medskip\noindent
설정은 다음과 같다. 다음이 주어졌다고 가정한다.
\begin{enumerate}
\item 환 $\Lambda$,
\item $\Lambda$ 위의 대수공간 $X$,
\item $\Lambda$-대수 $A$와
$X_A = X \times_{\Spec(\Lambda)} \Spec(A)$,
\item $A$ 위에서 평탄한 유한 표시 $\mathcal{O}_{X_A}$-가군
$\mathcal{F}$.
\end{enumerate}
이 상황에서
$$
0 \to I \to A' \to A \to 0
$$
꼴의 가능한 모든 전사를 생각한다. 여기서 $A'$은 핵 $I$가 $A'$에서
제곱 영 아이디얼인 $\Lambda$-대수이다. $A'$이 주어지면
$\Spec(\Lambda)$ 위의 대수공간의 일차 비후 $X_A \to X_{A'}$를 얻는다.
각 비후에 대해 $\mathcal{F}$를 $A'$ 위에서 평탄한 $X_{A'}$ 위의
유한 표시 가군 $\mathcal{F}'$으로 올리는 문제를 생각한다.
이 설정에서 『변형 이론』의 보조정리
\ref{defos-lemma-flat-ringed-topoi}의 결과들을 재현하고자 한다.

\medskip\noindent
더 정확히 말해, $\textit{Lift}(\mathcal{F}, A')$를 다음 쌍
$(\mathcal{F}', \alpha)$들의 범주라 하자. $\mathcal{F}'$은 $A'$ 위에서
평탄한 $X_{A'}$ 위의 유한 표시 가군이고,
$\alpha : \mathcal{F}'|_{X_A} \to \mathcal{F}$는 동형이다.
사상 $(\mathcal{F}'_1, \alpha_1) \to (\mathcal{F}'_2, \alpha_2)$는
$\alpha_1$, $\alpha_2$와 호환되는 동형
$\mathcal{F}'_1 \to \mathcal{F}'_2$이다.
$\textit{Lift}(\mathcal{F}, A')$의 동형류들의 집합을
$\text{Lift}(\mathcal{F}, A')$로 나타낸다.

\medskip\noindent
$\mathcal{G}$를 $A$ 위에서 평탄한 $X_\etale$ 위의
$\mathcal{O}_X \otimes_\Lambda A$-가군층이라 하자.
$\mathcal{G}'$가 $A'$ 위에서 평탄한
$\mathcal{O}_X \otimes_\Lambda A'$-가군층이고 $\beta$가 동형
$\mathcal{G}' \otimes_{A'} A \to \mathcal{G}$인 쌍
$(\mathcal{G}', \beta)$들의 범주
$\textit{Lift}(\mathcal{G}, A')$를 도입한다.

\begin{lemma}
\label{obsolete-lemma-equivalence}
위와 같은 표기법과 가정을 사용하자. $p : X_A \to X$를 사영이라 하자.
$A'$이 주어졌을 때 사영을 $p' : X_{A'} \to X$로 나타내자.
함자 $p'_*$는 다음 두 범주 사이의 동치를 유도한다.
\begin{enumerate}
\item 범주 $\textit{Lift}(\mathcal{F}, A')$,
\item 범주 $\textit{Lift}(p_*\mathcal{F}, A')$.
\end{enumerate}
\end{lemma}

\begin{proof}
원문에는 증명이 제시되어 있지 않다.
\end{proof}

\noindent
$\mathcal{H}$를 $A$ 위에서 평탄한 $\mathcal{C}_{X/\Lambda}$ 위의
$\mathcal{O} \otimes_\Lambda A$-가군층이라 하자.
$\textit{Lift}_\mathcal{O}(\mathcal{H}, A')$를 다음 쌍
$(\mathcal{H}', \gamma)$들을 대상으로 하는 범주로 도입한다.
$\mathcal{H}'$은 $A'$ 위에서 평탄한
$\mathcal{O} \otimes_\Lambda A'$-가군층이고,
$\gamma : \mathcal{H}' \otimes_A A' \to \mathcal{H}$는
$\mathcal{O} \otimes_\Lambda A$-가군의 동형이다.

\medskip\noindent
$\mathcal{G}$를 $A$ 위에서 평탄한 $X_\etale$ 위의
$\mathcal{O}_X \otimes_\Lambda A$-가군층이라 하자.
『여접 복합체』의 등식 (\ref{cotangent-equation-pi-spaces})에 나오는
사상 $i$와 $\pi$를 생각하자.
$\underline{\mathcal{G}} = \pi^{-1}(\mathcal{G})$로 나타내자.
이는 단순히 규칙 $(U \to \mathbf{A}) \mapsto \mathcal{G}(U)$로
주어지므로 $\underline{\mathcal{O}}_X \otimes_\Lambda A$-가군층이다.
같은 층을 $\mathcal{O} \otimes_\Lambda A$-가군층으로 보았을 때
$i_*\underline{\mathcal{G}}$로 나타낸다.

\begin{lemma}
\label{obsolete-lemma-second-equivalence}
위와 같은 표기법과 가정을 사용하자.
함자 $\pi_!$는 다음 두 범주 사이의 동치를 유도한다.
\begin{enumerate}
\item 범주
$\textit{Lift}_\mathcal{O}(i_*\underline{\mathcal{G}}, A')$,
\item 범주 $\textit{Lift}(\mathcal{G}, A')$.
\end{enumerate}
\end{lemma}

\begin{proof}
원문에는 증명이 제시되어 있지 않다.
\end{proof}

\begin{lemma}
\label{obsolete-lemma-second-equivalence-obs}
보조정리 \ref{obsolete-lemma-second-equivalence}와 같은 표기법과 가정을 사용하자.
$D(\mathcal{O}_X \otimes_\Lambda A)$의 대상
$$
L = L(\Lambda, X, A, \mathcal{G}) = L\pi_!(Li^*(i_*(\underline{\mathcal{G}})))
$$
을 생각하자. 제곱 영 핵 $I$를 갖는 $\Lambda$-대수의 전사
$A' \to A$가 주어지면 다음이 성립한다.
\begin{enumerate}
\item 범주 $\textit{Lift}(\mathcal{G}, A')$가 공집합이 아닐
필요충분조건은 어떤 류
$\xi \in \Ext^2_{\mathcal{O}_X \otimes A}(L,
\mathcal{G} \otimes_A I)$가 영인 것이다.
\item $\textit{Lift}(\mathcal{G}, A')$가 공집합이 아니면,
$\text{Lift}(\mathcal{G}, A')$는
$\Ext^1_{\mathcal{O}_X \otimes A}(L, \mathcal{G} \otimes_A I)$ 아래의
주동차 공간이다.
\item 올림 $\mathcal{G}'$이 주어졌을 때, $\text{id}_\mathcal{G}$로
당겨지는 $\mathcal{G}'$의 자기동형들의 집합은 표준적으로
$\Ext^0_{\mathcal{O}_X \otimes A}(L, \mathcal{G} \otimes_A I)$와
동형이다.
\end{enumerate}
\end{lemma}

\begin{proof}
원문에는 증명이 제시되어 있지 않다.
\end{proof}

\noindent
마지막으로 모든 것을 다음과 같이 합친다.

\begin{proposition}
\label{obsolete-proposition-conclusion}
$\Lambda$, $X$, $A$, $\mathcal{F}$가 위와 같다고 하자.
$D(X_A)$의 표준 대상 $L = L(\Lambda, X, A, \mathcal{F})$가 존재하여,
제곱 영 핵 $I$를 갖는 $\Lambda$-대수의 전사 $A' \to A$가 주어지면
다음이 성립한다.
\begin{enumerate}
\item 범주 $\textit{Lift}(\mathcal{F}, A')$가 공집합이 아닐
필요충분조건은 어떤 류 $\xi \in \Ext^2_{X_A}(L,
\mathcal{F} \otimes_A I)$가 영인 것이다.
\item $\textit{Lift}(\mathcal{F}, A')$가 공집합이 아니면,
$\text{Lift}(\mathcal{F}, A')$는
$\Ext^1_{X_A}(L, \mathcal{F} \otimes_A I)$ 아래의 주동차 공간이다.
\item 올림 $\mathcal{F}'$이 주어졌을 때, $\text{id}_\mathcal{F}$로
당겨지는 $\mathcal{F}'$의 자기동형들의 집합은 표준적으로
$\Ext^0_{X_A}(L, \mathcal{F} \otimes_A I)$와 동형이다.
\end{enumerate}
\end{proposition}

\begin{proof}
원문에는 증명이 제시되어 있지 않다.
\end{proof}

\begin{lemma}
\label{obsolete-lemma-pseudo-coherent}
명제 \ref{obsolete-proposition-conclusion}의 상황에서
$X \to \Spec(\Lambda)$가 국소 유한형이고 $\Lambda$가 뇌터이면
$L$은 유사 연접이다.
\end{lemma}

\begin{proof}
원문에는 증명이 제시되어 있지 않다.
\end{proof}

\section{비평탄 경우의 연접층 스택}
\label{obsolete-section-not-flat}

\noindent
『Quot』의 정리 \ref{quot-theorem-coherent-algebraic}에서
$f : X \to B$가 평탄하다는 가정은 필요하지 않다.
이 절에서는 유도 대수기하학의 아이디어를 바탕으로 증명 방법을 수정하여
평탄성 가정을 피한다. 『Quot』의 절 \ref{quot-section-not-flat}에서는
완전히 다른 방법으로 정확히 같은 결과를 얻는다. 이것이 이 절의 방법이
폐기된 이유이다.

\medskip\noindent
『Quot』의 정리 \ref{quot-theorem-coherent-algebraic}의 증명에서 평탄성을
사용하는 유일한 단계는 『Quot』의 보조정리
\ref{quot-lemma-coherent-defo-thy}를 적용하는 부분이다. 이 보조정리는
『Artin의 공리』의 절 \ref{artin-section-dual}에서와 같은 장애물 이론을
구성하는 데 사용된다. 그 증명은 『변형 이론』의 보조정리
\ref{defos-lemma-flat-ringed-topoi}와
\ref{defos-lemma-verify-iv-ringed-topoi}에 의존한다. 이 보조정리들은
『변형 이론』의 절 \ref{defos-section-flat-ringed-topoi}에 있다. 이 때문에 $f$가
평탄하다는 가정이 생긴다. 계속하기 전에 『변형 이론』의 보조정리
\ref{defos-lemma-flat-ringed-topoi}의 결과 (2), (3)은 $f$가 평탄하다는
가정 없이도 성립함에 유의하자. 실제로 이 결과들은 평탄성 가정이 없는
『변형 이론』의 보조정리 \ref{defos-lemma-inf-ext-rel-ringed-topoi}와
\ref{defos-lemma-inf-map-rel-ringed-topoi}에 의존한다.

\medskip\noindent
세부 사항에 들어가기 전에 유도 대수기하학에서 온 이 구성의 동기를
설명한다. 이 설명이 뒤따르는 내용을 분명하게 해 줄 것이라 생각한다.
$A$를 국소 뇌터 밑 $S$ 위의 유한형 대수라 하자.
$X$를 $A$로 ``유도 밑변환''한 것을 $X \otimes^\mathbf{L} A$라 쓰고,
표준 포함 사상을 $i : X_A \to X \otimes^\mathbf{L} A$로 나타내자.
대상 $X \otimes^\mathbf{L} A$는 Stacks Project에 (아직) 정의되어 있지
않다. 이를 호몰로지 층이
$$
H_i(\mathcal{O}_{X \otimes^\mathbf{L} A}) =
\text{Tor}^{\mathcal{O}_S}_i(\mathcal{O}_X, A).
$$
인 단체적 환의 층 $\mathcal{O}_{X \otimes^\mathbf{L} A}$가 주어진
대수공간 $X_A$라고 생각할 수 있다. 사상
$X \otimes^\mathbf{L} A \to \Spec(A)$는 평탄하다(단체적 환의 층의
각 항이 $A$-평탄하다). 따라서 평탄 가군의 변형에 관한 보통의 내용을
적용할 수 있다. 이로부터 장애물 이론을 얻으며, 사용되는 군은
$$
\Ext^i_{X \otimes^\mathbf{L} A}(i_*\mathcal{F},
i_*\mathcal{F} \otimes_A M)
$$
이다. 여기서 $i = 0, 1, 2$는 각각 무한소 자기동형, 무한소 변형,
장애물에 대응한다. $i_*\mathcal{F}$를
$X \otimes^\mathbf{L} A'$로 평탄하게 변형한 것은 자동으로
$i'_*\mathcal{F}'$ 꼴이며, 여기서 $\mathcal{F}'$은 $\mathcal{F}$의
평탄 변형이다. 함자 $Li^*$와 $i_* = Ri_*$의 수반성에 의해 이 Ext 군들은
$$
\Ext^i_{X_A}(Li^*(i_*\mathcal{F}), \mathcal{F} \otimes_A M)
$$
과 같다. 따라서 『Quot』의 보조정리
\ref{quot-lemma-coherent-defo-thy}의 증명과 정확히 같은 꼴의 장애물 군을
얻으며, 유일한 차이는 $\mathcal{F}$가 처음 나타나는 곳을 복합체
$Li^*(i_*\mathcal{F})$로 바꾸는 것이다.

\medskip\noindent
아래에서는 $E(\mathcal{F}) = Li^*(i_*\mathcal{F})$를 ``직접'' 구성하고
$\mathcal{F}$의 변형 이론과 맺는 관계를 직접 증명하여 보조정리의
비평탄 판을 증명한다. 실제로 우리가 관심을 갖는 것은 $i = 0, 1, 2$에
대한 Ext 군
$\Ext^i_{X_A}(Li^*(i_*\mathcal{F}), \mathcal{F} \otimes_A M)$뿐이므로
$\tau_{\geq -2}E(\mathcal{F})$만 구성하면 충분하다. 코호몰로지 층을
다음과 같이 식별할 수도 있다.
$$
H^i(E(\mathcal{F})) =
\left\{
\begin{matrix}
0 & \text{조건 }i > 0 \\
\mathcal{F} & \text{조건 } i = 0 \\
0 & \text{조건 } i = -1 \\
\text{Tor}_1^{\mathcal{O}_S}(\mathcal{O}_X, A)
\otimes_{\mathcal{O}_X} \mathcal{F} &
\text{조건 } i = -2
\end{matrix}
\right.
$$
이 관찰은 아래 주석들에서 $E(\mathcal{F})$를 구성하는 길잡이가 된다.

\begin{remark}[직접 구성]
\label{obsolete-remark-construction-E}
$S$를 스킴이라 하자. $f : X \to B$를 $S$ 위의 대수공간의 사상이라
하고 $U$를 $B$ 위의 또 다른 대수공간이라 하자.
두 번째 사영을 $q : X \times_B U \to U$로 나타내자.
『여접 복합체』의 절 \ref{cotangent-section-fibre-product}의 구별 삼각형
$$
Lq^*L_{U/B} \to L_{X \times_B U/B} \to E \to Lq^*L_{U/B}[1]
$$
을 생각하자. 임의의 $\mathcal{O}_{X \times_B U}$-가군층
$\mathcal{F}$에 대해 Atiyah 류
$$
\mathcal{F} \to
L_{X \times_B U/B}
\otimes_{\mathcal{O}_{X \times_B U}}^\mathbf{L} \mathcal{F}[1]
$$
가 있다. 『여접 복합체』의 절
\ref{cotangent-section-atiyah-general}을 보라.
이를 $E$로 가는 사상과 합성하고 $D(\mathcal{O}_{X \times_B U})$의
구별 삼각형
$$
E(\mathcal{F}) \to \mathcal{F} \to
\mathcal{F} \otimes_{\mathcal{O}_{X \times_B U}}^\mathbf{L} E[1] \to
E(\mathcal{F})[1]
$$
을 택할 수 있다. 구성에 의해 Atiyah 류는 다음 사상으로 올라간다.
$$
e_\mathcal{F} :
E(\mathcal{F})
\longrightarrow
Lq^*L_{U/B} \otimes_{\mathcal{O}_{X \times_B U}}^\mathbf{L} \mathcal{F}[1]
$$
이 사상은 다음 구별 삼각형들의 사상에 들어간다.
$$
\xymatrix{
\mathcal{F} \otimes^\mathbf{L} Lq^*L_{U/B}[1] \ar[r] &
\mathcal{F} \otimes^\mathbf{L} L_{X \times_B U/B}[1] \ar[r] &
\mathcal{F} \otimes^\mathbf{L} E[1] \\
E(\mathcal{F}) \ar[r] \ar[u]^{e_\mathcal{F}} &
\mathcal{F} \ar[r] \ar[u]^{Atiyah} &
\mathcal{F} \otimes^\mathbf{L} E[1] \ar[u]^{=}
}
$$
$S, B, X, f, U, \mathcal{F}$가 주어졌을 때 $E(\mathcal{F})$와
$e_\mathcal{F}$의 선택 하나를 고정한다.
\end{remark}

\begin{remark}[장애물 류의 구성]
\label{obsolete-remark-construction-ob}
주석 \ref{obsolete-remark-construction-E}의 표기법을 사용하고,
$i : U \to U'$을 $B$ 위에서 $U$의 일차 비후라 하자.
$\mathcal{I} \subset \mathcal{O}_{U'}$를 $B'$에서 $B$를 잘라 내는
준연접 아이디얼층이라 하자. 기본 삼각형
$$
Li^*L_{U'/B} \to L_{U/B} \to L_{U/U'} \to Li^*L_{U'/B}[1]
$$
과 사상 $L_{U/U'} \to \mathcal{I}[1]$은 사상
$e_{U'} : L_{U/B} \to \mathcal{I}[1]$을 결정한다. 이를 앞 주석의
사상 $e_\mathcal{F}$와 결합하면
$$
(\text{id}_\mathcal{F} \otimes Lq^*e_{U'}) \cup e_\mathcal{F} :
E(\mathcal{F})
\longrightarrow
\mathcal{F} \otimes_{\mathcal{O}_{X \times_B U}} q^*\mathcal{I}[2]
$$
를 얻는다(유도 텐서곱에서 보통의 텐서곱으로 가는 사상과도 합성했다).
다시 말해 다음 원소를 얻는다.
$$
\xi_{U'} \in
\Ext^2_{\mathcal{O}_{X \times_B U}}(
E(\mathcal{F}),
\mathcal{F} \otimes_{\mathcal{O}_{X \times_B U}} q^*\mathcal{I})
$$
\end{remark}

\begin{lemma}
\label{obsolete-lemma-ob-is-obstruction}
주석 \ref{obsolete-remark-construction-ob}의 상황에서 $\mathcal{F}$가 $U$ 위에서
평탄하다고 가정하자. 류 $\xi_{U'}$가 영일 필요충분조건은
$i^*\mathcal{F}' \cong \mathcal{F}$인, $U'$ 위에서 평탄한
$\mathcal{O}_{X \times_B U'}$-가군 $\mathcal{F}'$이 존재하는 것이다.
\end{lemma}

\begin{proof}[증명 (개요)]
『변형 이론』의 보조정리
\ref{defos-lemma-inf-obs-ext-rel-ringed-topoi}의 판정법을 사용하겠다.
$\mathcal{O} = \mathcal{O}_{X \times_B U}$,
$\mathcal{O}' = \mathcal{O}_{X \times_B U'}$로 줄여 쓰겠다.
짧은 완전열
$$
0 \to \mathcal{I} \to \mathcal{O}_{U'} \to \mathcal{O}_U \to 0.
$$
을 생각하자. $\mathcal{J} \subset \mathcal{O}'$를
$X \times_B U$를 잘라 내는 준연접 아이디얼층이라 하자.
위 완전열로부터 완전열
$$
\text{Tor}_1^{\mathcal{O}_B}(\mathcal{O}_X, \mathcal{O}_U) \to
q^*\mathcal{I} \to \mathcal{J} \to 0
$$
을 얻는다. 여기서
$\text{Tor}_1^{\mathcal{O}_B}(\mathcal{O}_X, \mathcal{O}_U)$는 다음을
줄여 쓴 것이다.
$$
\text{Tor}_1^{h^{-1}\mathcal{O}_B}(p^{-1}\mathcal{O}_X, q^{-1}\mathcal{O}_U)
\otimes_{(p^{-1}\mathcal{O}_X\otimes_{h^{-1}\mathcal{O}_B}q^{-1}\mathcal{O}_U)}
\mathcal{O}.
$$
$\mathcal{F}$와 텐서곱하면 완전열
$$
\mathcal{F} \otimes_\mathcal{O}
\text{Tor}_1^{\mathcal{O}_B}(\mathcal{O}_X, \mathcal{O}_U) \to
\mathcal{F} \otimes_\mathcal{O}
q^*\mathcal{I} \to
\mathcal{F} \otimes_\mathcal{O} \mathcal{J} \to 0
$$
을 얻는다. (『변형 이론』의 보조정리
\ref{defos-lemma-inf-obs-ext-rel-ringed-topoi}의 표기법과 비교하면 문자
$\mathcal{I}$와 $\mathcal{J}$의 역할이 뒤바뀌었음에 유의하자.)
그 보조정리의 조건 (1)은 위의 마지막 사상이 동형이라는 것, 즉 첫 번째
사상이 영이라는 것이다. 이 사상의 소멸은 기하학적 점
$\overline{z} = (\overline{x}, \overline{u}) : \Spec(k) \to
X \times_B U$에서 줄기별로 확인할 수 있다.
$R = \mathcal{O}_{B, \overline{b}}$,
$A = \mathcal{O}_{X, \overline{x}}$,
$B = \mathcal{O}_{U, \overline{u}}$,
$C = \mathcal{O}_{\overline{z}}$로 놓는다.
『여접 복합체』의 보조정리 \ref{cotangent-lemma-fibre-product}와
$E(\mathcal{F})$의 정의 삼각형에 의해
$$
H^{-2}(E(\mathcal{F}))_{\overline{z}} =
\mathcal{F}_{\overline{z}} \otimes \text{Tor}_1^R(A, B)
$$
이다. 따라서 사상 $\xi_{U'}$는 사상
$$
\mathcal{F}_{\overline{z}} \otimes \text{Tor}_1^R(A, B)
\longrightarrow
\mathcal{F}_{\overline{z}} \otimes_B \mathcal{I}_{\overline{u}}
$$
을 유도한다. 이 사상이 위에서 기술한 사상의 줄기와 같다고 주장한다
(증명은 생략한다. 이는 순수하게 환론적인 명제이다).
따라서 『변형 이론』의 보조정리
\ref{defos-lemma-inf-obs-ext-rel-ringed-topoi}의 조건 (1)은
$H^{-2}(\xi_{U'}) :
H^{-2}(E(\mathcal{F})) \to \mathcal{F} \otimes \mathcal{I}$가 영인 것과
동치이다.

\medskip\noindent
증명을 끝내기 위해 조건 (1)이 성립한다고 가정할 때 조건 (2)가
$\xi_{U'}$의 소멸과 동치임을 보인다. 증명의 나머지에서는
$\mathcal{F} \otimes \mathcal{I}$로
$\mathcal{F} \otimes_\mathcal{O} q^*\mathcal{I} =
\mathcal{F} \otimes_\mathcal{O} \mathcal{J}$를 나타낸다.
$H^0(E(\mathcal{F})) = \mathcal{F}$와
$H^{-1}(E(\mathcal{F})) = 0$을 사용하여 스펙트럼 열
$$
\Ext^i(H^{-j}(E(\mathcal{F})), \mathcal{F} \otimes \mathcal{I})
\Rightarrow
\Ext^{i + j}(E(\mathcal{F}), \mathcal{F} \otimes \mathcal{I})
$$
을 살펴보면 완전열
$$
0 \to
\Ext^2(\mathcal{F}, \mathcal{F} \otimes \mathcal{I}) \to
\Ext^2(E(\mathcal{F}), \mathcal{F} \otimes \mathcal{I}) \to
\Hom(H^{-2}(E(\mathcal{F})), \mathcal{F} \otimes \mathcal{I})
$$
이 있음을 알 수 있다. 따라서 원소 $\xi_{U'}$는
$\Ext^2(\mathcal{F}, \mathcal{F} \otimes \mathcal{I})$의 원소이다.
이 원소가 『변형 이론』의 보조정리
\ref{defos-lemma-inf-obs-ext-rel-ringed-topoi}의 원소와 일치함을 보이면
증명이 끝난다. 이 확인은 생략한다.
\end{proof}

\begin{lemma}
\label{obsolete-lemma-coherent-defo-thy-general}
『Quot』의 상황 \ref{quot-situation-coherent}에서 $S$가 국소 뇌터
스킴이고 $S = B$라고 가정하자.
$\mathcal{X} = \textit{Coh}_{X/B}$로 놓는다.
그러면 $\mathcal{X}$에 대해 버설성의 개방성이 성립한다
(『Artin의 공리』의 정의
\ref{artin-definition-openness-versality}를 보라).
\end{lemma}

\begin{proof}[증명 (개요)]
$U \to S$를 스킴의 유한형 사상이라 하고, $x$를 $U$ 위의
$\mathcal{X}$의 대상, $u_0 \in U$를 유한형 점이라 하며 $x$가
$u_0$에서 버설이라고 하자.
$U$를 줄인 뒤 $u_0$가 닫힌 점이라고 가정할 수 있고
(『스킴의 사상』의 보조정리
\ref{morphisms-lemma-point-finite-type}), $U = \Spec(A)$이며
$U \to S$가 $S$의 아핀 열린집합 $\Spec(\Lambda)$ 안으로 사상된다고
가정할 수 있다. 보조정리를 증명하기 위해 『Artin의 공리』의 보조정리
\ref{artin-lemma-dual-openness}를 사용하겠다.
$\mathcal{F}$를 주어진 대상 $x$에 대응하며 $A$ 위에서 평탄한
$X_A = \Spec(A) \times_S X$ 위의 연접 가군이라 하자.

\medskip\noindent
주석 \ref{obsolete-remark-construction-E}에서와 같이 $E(\mathcal{F})$와
$e_\mathcal{F}$를 택한다. $E(\mathcal{F})$의 코호몰로지 층에 관한
기술로부터 임의의 $A$-가군 $M$에 대해
$$
\Ext^1(E(\mathcal{F}), \mathcal{F} \otimes_A M) =
\Ext^1(\mathcal{F}, \mathcal{F} \otimes_A M)
$$
임을 알 수 있다. 이것과 『변형 이론』의 보조정리
\ref{defos-lemma-inf-ext-rel-ringed-topoi}를 사용하면 함자들의 동형
$$
T_x(M) = \Ext^1_{X_A}(E(\mathcal{F}), \mathcal{F} \otimes_A M)
$$
을 얻는다. 보조정리 \ref{obsolete-lemma-ob-is-obstruction}에 의해 제곱 영 핵
$I$를 갖는 $\Lambda$-대수의 임의의 전사 $A' \to A$에 대해 장애물 류
$$
\xi_{A'} \in \Ext^2_{X_A}(E(\mathcal{F}), \mathcal{F} \otimes_A I)
$$
가 있다. $m = 2$로 놓고 $i \leq m$에 대한 Ext 군
$\Ext^i_{X_A}(E(\mathcal{F}), \mathcal{F} \otimes_A M)$의 계산에
『공간의 유도 범주』의 보조정리
\ref{spaces-perfect-lemma-compute-ext}를 적용한다.
$E(\mathcal{F})$가 $D^-_{\textit{Coh}}$에 속함의 확인은 생략한다.
힌트: 『여접 복합체』의 보조정리
\ref{cotangent-lemma-cotangent-finite}를 사용하라.
완전 대상 $K \in D(A)$와 경계 사상들과 호환되는 함자적 동형
$$
H^i(K \otimes_A^\mathbf{L} M)
\longrightarrow
\Ext^i_{X_A}(E(\mathcal{F}), \mathcal{F} \otimes_A M)
$$
을 $i \leq m$에 대해 얻는다. 이 대상 $K$와 위에 표시한 동일시들은
『Artin의 공리』의 상황 \ref{artin-situation-dual}과 같은 자료를 준다.
마지막으로 『Artin의 공리』의 보조정리
\ref{artin-lemma-dual-obstruction}의 조건 (iv)는 『변형 이론』의 보조정리
\ref{defos-lemma-verify-iv-ringed-topoi}의 한 변형에 의해 성립한다.
그 변형의 명제와 증명은 생략한다. 따라서 『Artin의 공리』의 보조정리
\ref{artin-lemma-dual-openness}를 적용할 수 있고 보조정리가 증명된다.
\end{proof}

\begin{theorem}
\label{obsolete-theorem-coherent-algebraic-general}
$S$를 스킴이라 하고 $f : X \to B$를 $S$ 위의 대수공간의 사상이라 하자.
$f$가 유한 표시이고 분리라고 가정하자. 그러면
$\textit{Coh}_{X/B}$는 $S$ 위의 대수적 스택이다.
\end{theorem}

\begin{proof}
이 정리는 『Quot』의 정리
\ref{quot-theorem-coherent-algebraic-general}을 복사한 것이다.
여기에 사본을 둔 이유는 이 절의 내용으로 두 번째 증명을 얻기 때문이다
(이 절의 시작에서 논의했다). 실제로 『Quot』의 정리
\ref{quot-theorem-coherent-algebraic}의 증명과 정확히 같은 방식으로
논증하되, 보조정리 \ref{obsolete-lemma-coherent-defo-thy-general}을
『Quot』의 보조정리 \ref{quot-lemma-coherent-defo-thy} 대신 대입한다.
\end{proof}

\section{수정}
\label{obsolete-section-modifications}

\noindent
다음은 『형식 공간의 대수화』의 등식
(\ref{restricted-equation-modification})에 나오는 범주에 관한 폐기된
결과이다. 현재의 내용은 『형식 공간의 대수화』의 절
\ref{restricted-section-modifications}을 보라.

\begin{lemma}
\label{obsolete-lemma-henselian}
$(A, \mathfrak m, \kappa)$를 뇌터 국소환이라 하자.
$A$에 대한 『형식 공간의 대수화』의 등식
(\ref{restricted-equation-modification})의 범주는 $A$의 Hensel화
$A^h$에 대한 『형식 공간의 대수화』의 등식
(\ref{restricted-equation-modification})의 범주와 동치이다.
\end{lemma}

\begin{proof}
이는 『형식 공간의 대수화』의 보조정리
\ref{restricted-lemma-equivalence-to-completion}의 특수한 경우이다.
\end{proof}

\noindent
유리 특이점에 관한 다음 보조정리는 표면 특이점의 해소 장에서 더 이상
필요하지 않다.

\begin{lemma}
\label{obsolete-lemma-double-dual-rational}
『표면의 해소』의 상황 \ref{resolve-situation-rational}에서,
$M$을 유한 반사 $A$-가군이라 하자. 딸린 $\mathcal{O}_S$-가군의
당김뒤를 $M \otimes_A \mathcal{O}_X$로 나타내자. 그러면
$M \otimes_A \mathcal{O}_X$는 그 이중 쌍대 위로 사상된다.
\end{lemma}

\begin{proof}
이중 쌍대를 $\mathcal{F} = (M \otimes_A \mathcal{O}_X)^{**}$로 놓고,
$\mathcal{F}' \subset \mathcal{F}$를 평가 사상
$M \otimes_A \mathcal{O}_X \to \mathcal{F}$의 상이라 하자.
그러면 짧은 완전열
$$
0 \to \mathcal{F}' \to \mathcal{F} \to \mathcal{Q} \to 0
$$
을 얻는다. $X$가 정규이므로 여차원 $1$인 점에서 국소환
$\mathcal{O}_{X, x}$는 이산 값매김환이다(『스킴의 성질』의 보조정리
\ref{properties-lemma-criterion-normal}을 보라).
따라서 『대수학 심화』의 보조정리
\ref{more-algebra-lemma-cokernel-map-double-dual-dvr}에 의해 그런 점들에서
$\mathcal{Q}_x = 0$이다. 그러므로 $\mathcal{Q}$는 유한 개의 닫힌 점에
지지되고, 『스킴의 코호몰로지』의 보조정리
\ref{coherent-lemma-coherent-support-dimension-0}에 의해 대역적으로
생성된다. $\mathcal{F}'$도 대역 절단들로 생성되므로
(『표면의 해소』의 보조정리
\ref{resolve-lemma-globally-generated}) 완전열
$$
0 \to H^0(X, \mathcal{F}') \to H^0(X, \mathcal{F}) \to
H^0(X, \mathcal{Q}) \to 0
$$
을 얻는다. $X \to \Spec(A)$는 닫힌 점의 여집합 위에서 동형이고
$M$은 반사이므로, 사상
$$
M \to H^0(X, \mathcal{F}') \to H^0(X, \mathcal{F})
$$
은 $A$의 임의의 비극대 소 아이디얼에서 국소화하면 동형을 유도한다.
따라서 『대수학 심화』의 보조정리
\ref{more-algebra-lemma-check-isomorphism-via-depth-and-ass}와 정규환 위의
반사 가군이 성질 $(S_2)$를 갖는다는 사실(『대수학 심화』의 보조정리
\ref{more-algebra-lemma-reflexive-over-normal})에 의해 이 사상들은
동형이다. 따라서 원하는 대로 $\mathcal{Q} = 0$이다.
\end{proof}

\section{교차 이론}
\label{obsolete-section-intersection-theory}

\begin{lemma}
\label{obsolete-lemma-good-blowing-up}
$b : X' \to X$를 체 $k$ 위의 매끄러운 사영 스킴을 매끄러운 닫힌 부분스킴
$Z \subset X$에서 부풀리기한 것이라 하자. 다음 도표를 생각하자.
$$
\xymatrix{
E \ar[r]_j \ar[d]_\pi & X' \ar[d]^b \\
Z \ar[r]^i & X
}
$$
$K_0(X)$의 어떤 원소를 $Z$로 제한하면 $K_0(Z)$에서
$\mathcal{C}_{Z/X}$의 류와 같다고 가정하자. 그러면 어떤
$\alpha'' \in K_0(X')$에 대해
$K_0(X')$에서 $[Lb^*\mathcal{O}_Z] = [\mathcal{O}_E] \cdot \alpha''$이다.
\end{lemma}

\begin{proof}
스킴 $X$, $X'$, $E$, $Z$는 $k$ 위에서 매끄럽고 사영이므로
$K'_0(X) = K_0(X) = K_0(\textit{Vect}(X)) =
K_0(D^b_{\textit{Coh}}(X)))$이고, 나머지 $3$개의 스킴에 대해서도 마찬가지이다.
『스킴의 유도 범주』의 보조정리
\ref{perfect-lemma-Noetherian-Kprime},
\ref{perfect-lemma-Kprime-K}, \ref{perfect-lemma-K-is-old-K}를 보라.
이 증명에서는 이 판들 사이를 자유롭게 오가겠다.
$\mathcal{F}$를 정의하는 유한 국소 자유 $\mathcal{O}_E$-가군의
짧은 완전열
$$
0 \to \mathcal{F} \to \pi^*\mathcal{C}_{Z/X} \to
\mathcal{C}_{E/X'} \to 0
$$
을 생각하자. $\mathcal{C}_{E/X'} = \mathcal{O}_{X'}(-E)|_E$는 가역
$\mathcal{O}_X$-가군 $\mathcal{O}_{X'}(-E)$의 제한임에 유의하자.
$K_0(Z)$에서 $i^*\alpha = [\mathcal{C}_{Z/X}]$인
$\alpha \in K_0(X)$를 택한다.
$\alpha' = b^*\alpha - [\mathcal{O}_{X'}(-E)]$로 놓자.
그러면 $j^*\alpha' = [\mathcal{F}]$이다. 『Weil 코호몰로지 이론』의
보조정리 \ref{weil-lemma-lambda-operations}에 의해
$j^*\lambda^i(\alpha') = [\wedge^i(\mathcal{F})]$임을 얻는다.
이는 $K_0(X)$에서
$[\mathcal{O}_E] \cdot \alpha' = [\wedge^i\mathcal{F}]$이라는 뜻이다.
『스킴의 유도 범주』의 보조정리
\ref{perfect-lemma-projection-formula}를 보라.
$r$을 $X$ 안에서 $Z$의 기약 성분들이 갖는 여차원의 최댓값이라 하자.
생략하는 계산에 의하면 $i \geq 0, 1, \ldots, r - 1$에 대해
$H^{-i}(Lb^*\mathcal{O}_Z) = \wedge^i\mathcal{F}$이고 다른 차수에서는
영이다. 따라서 $K_0(X)$에서
\begin{align*}
[Lb^*\mathcal{O}_Z] & =
\sum\nolimits_{i = 0, \ldots, r - 1} (-1)^i[\wedge^i\mathcal{F}] \\
& =
\sum\nolimits_{i = 0, \ldots, r - 1} (-1)^i[\mathcal{O}_E] \lambda^i(\alpha') \\
& =
[\mathcal{O}_E] \left(\sum\nolimits_{i = 0, \ldots, r - 1}
(-1)^i \lambda^i(\alpha')\right)
\end{align*}
이다. 이로써
$\alpha'' = \sum_{i = 0, \ldots, r - 1} (-1)^i \lambda^i(\alpha')$에 대해
보조정리가 증명된다.
\end{proof}

\begin{lemma}
\label{obsolete-lemma-gysin-factors-principal}
$(S, \delta)$가 『Chow 호몰로지』의 상황
\ref{chow-situation-setup}에서와 같다고 하자.
$X$를 $S$ 위에서 국소 유한형이라 하고, $X$는 정수적이며
$n = \dim_\delta(X)$라고 하자.
$a \in \Gamma(X, \mathcal{O}_X)$를 영이 아닌 함수라 하고,
$i : D = Z(a) \to X$를 $a$의 영점 스킴의 닫힌 몰입이라 하자.
$f \in R(X)^*$라 하자. 이 경우 $A_{n - 2}(D)$에서
$i^*\text{div}_X(f) = 0$이다.
\end{lemma}

\begin{proof}
『Chow 호몰로지』의 보조정리
\ref{chow-lemma-gysin-factors-general}의 특수한 경우이다.
\end{proof}

\begin{remark}
\label{obsolete-remark-not-true-not-quasi-compact}
이 주석은 전에는 『Chow 호몰로지』의 보조정리
\ref{chow-lemma-cycles-k-group}의 화살표들이 일반적으로 동형인지
분명하지 않다고 말했다. 그러나 이제 이 사실의 증명을 찾았다.
\end{remark}






\subsection{부풀리기 보조정리}
\label{obsolete-section-blowing-up-lemmas}

\noindent
이 절에서는 부풀리기 위에서 적절한 유효 Cartier 인자로 Cartier 인자를
표현하는 것에 관한 몇 가지 보조정리를 증명한다. 이 보조정리들은
\cite[제2.4절]{F}에서 찾을 수 있다. 유한형이 아닌 설정에서도
작동하도록 명제를 조정했다. 보조정리
\ref{obsolete-lemma-blowing-up-intersections}의 사상 $b$는 무한히 많은 부풀리기의
합성일 수 있지만, 주어진 임의의 준콤팩트 열린집합 $W \subset X$ 위에서는
유한 번의 부풀리기만 필요하다(그리고 이것이 앞서 인용한 곳의 결과이다).

\begin{lemma}
\label{obsolete-lemma-push-pull-effective-Cartier}
$(S, \delta)$가 『Chow 호몰로지』의 상황
\ref{chow-situation-setup}에서와 같다고 하자.
$X$, $Y$를 $S$ 위에서 국소 유한형이라 하고 $f : X \to Y$를 고유
사상이라 하자. $D \subset Y$를 유효 Cartier 인자라 하자.
$X$, $Y$가 정수적이고
$n = \dim_\delta(X) = \dim_\delta(Y)$이며 $f$가 우세라고 가정하자.
그러면
$$
f_*[f^{-1}(D)]_{n - 1} = [R(X) : R(Y)] [D]_{n - 1}.
$$
특히 $f$가 쌍유리이면
$f_*[f^{-1}(D)]_{n - 1} = [D]_{n - 1}$이다.
\end{lemma}

\begin{proof}
『Chow 호몰로지』의 보조정리
\ref{chow-lemma-equal-c1-as-cycles}와, $D$가
$\mathcal{O}_X(D)$의 표준 절단 $1_D$의 영점 스킴이라는 사실에서
즉시 따른다.
\end{proof}

\begin{lemma}
\label{obsolete-lemma-blowing-up-denominators}
$(S, \delta)$가 『Chow 호몰로지』의 상황
\ref{chow-situation-setup}에서와 같다고 하자.
$X$를 $S$ 위에서 국소 유한형이라 하고, $X$가 정수적이며
$\dim_\delta(X) = n$이라고 가정하자.
$\mathcal{L}$을 가역 $\mathcal{O}_X$-가군이라 하고 $s$를
$\mathcal{L}$의 영이 아닌 유리형 절단이라 하자.
$s$가 $U$ 위의 $\mathcal{L}$의 절단에 대응하도록 하는 가장 큰 열린
부분스킴을 $U \subset X$라 하자. 다음을 만족하는 사영 사상
$$
\pi : X' \longrightarrow X
$$
가 존재한다.
\begin{enumerate}
\item $X'$은 정수적이고,
\item $\pi|_{\pi^{-1}(U)} : \pi^{-1}(U) \to U$는 동형이며,
\item 다음을 만족하는 유효 Cartier 인자 $D, E \subset X'$이 존재하고,
$$
\pi^*\mathcal{L} = \mathcal{O}_{X'}(D - E),
$$
\item 유리형 절단 $s$는 위 동형을 통해 유리형 절단
$1_D \otimes (1_E)^{-1}$에 대응하며(『인자』의 정의
\ref{divisors-definition-invertible-sheaf-effective-Cartier-divisor}),
\item $Z_{n - 1}(X)$에서
$$
\pi_*([D]_{n - 1} - [E]_{n - 1}) = \text{div}_\mathcal{L}(s)
$$
이다.
\end{enumerate}
\end{lemma}

\begin{proof}
$\mathcal{I} \subset \mathcal{O}_X$를 $s$의 분모들의 준연접 아이디얼층이라
하자. 『인자』의 정의
\ref{divisors-definition-regular-meromorphic-ideal-denominators}를 보라.
『인자』의 보조정리 \ref{divisors-lemma-blowing-up-denominators}에서
(2), (3), (4)를 얻는다. 『인자』의 보조정리
\ref{divisors-lemma-blow-up-integral-scheme}에서 (1)을 얻는다.
『인자』의 보조정리 \ref{divisors-lemma-blowing-up-projective}에 의해
사상 $\pi$는 사영이다. 이제 (5)를 증명하면 된다.
『Chow 호몰로지』의 보조정리
\ref{chow-lemma-equal-c1-as-cycles}에 의해
$$
\pi_*(\text{div}_{\mathcal{L}'}(s')) = \text{div}_\mathcal{L}(s).
$$
따라서 $\text{div}_{\mathcal{L}'}(s') =
[D]_{n - 1} - [E]_{n - 1}$임을 보이면 충분하다.
이는 등식 $s' = 1_D \otimes 1_E^{-1}$과 가법성에서 따른다.
『인자』의 보조정리 \ref{divisors-lemma-c1-additive}를 보라.
\end{proof}

\begin{definition}
\label{obsolete-definition-epsilon}
$(S, \delta)$가 『Chow 호몰로지』의 상황
\ref{chow-situation-setup}에서와 같다고 하자.
$X$를 $S$ 위에서 국소 유한형이라 하고, $X$가 정수적이며
$\dim_\delta(X) = n$이라고 가정하자.
$D_1, D_2$를 $X$의 두 유효 Cartier 인자라 하자.
$Z \subset X$를 $\dim_\delta(Z) = n - 1$인 정수적 닫힌 부분스킴이라
하자. 이 상황의 {\it $\epsilon$-불변량}은
$$
\epsilon_Z(D_1, D_2) = n_Z \cdot m_Z
$$
이다. 여기서 $n_Z$와 $m_Z$는 각각 $(n - 1)$-순환
$[D_1]_{n - 1}$과 $[D_2]_{n - 1}$에서 $Z$의 계수이다.
\end{definition}

\begin{lemma}
\label{obsolete-lemma-two-divisors}
$(S, \delta)$가 『Chow 호몰로지』의 상황
\ref{chow-situation-setup}에서와 같다고 하자.
$X$를 $S$ 위에서 국소 유한형이라 하고, $X$가 정수적이며
$\dim_\delta(X) = n$이라고 가정하자.
$D_1, D_2$를 $X$의 두 유효 Cartier 인자라 하자.
$Z$를 스킴 $D_1 \cap D_2$의 열린닫힌 부분스킴이라 하자.
$\dim_\delta(D_1 \cap D_2 \setminus Z) \leq n - 2$라고 가정하자.
그러면 사상 $b : X' \to X$와 $X'$ 위의 Cartier 인자
$D_1', D_2', E$가 존재하여 다음을 만족한다.
\begin{enumerate}
\item $X'$은 정수적이고,
\item $b$는 사영이며,
\item $b$는 닫힌 부분스킴 $Z$에서 $X$를 부풀리기한 것이고,
\item $E = b^{-1}(Z)$이며,
\item $b^{-1}(D_1) = D'_1 + E$이고 $b^{-1}D_2 = D_2' + E$이며,
\item $\dim_\delta(D'_1 \cap D'_2) \leq n - 2$이고,
$Z = D_1 \cap D_2$이면 $D'_1 \cap D'_2 = \emptyset$이며,
\item $\dim_\delta(W') = n - 1$인 모든 정수적 닫힌 부분스킴 $W'$에
대해 다음이 성립한다.
\begin{enumerate}
\item $\epsilon_{W'}(D'_1, E) > 0$이면 $W = b(W')$로 놓았을 때
$\dim_\delta(W) = n - 1$이고
$$
\epsilon_{W'}(D'_1, E) < \epsilon_W(D_1, D_2),
$$
\item $\epsilon_{W'}(D'_2, E) > 0$이면 $W = b(W')$로 놓았을 때
$\dim_\delta(W) = n - 1$이고
$$
\epsilon_{W'}(D'_2, E) < \epsilon_W(D_1, D_2),
$$
\end{enumerate}
\end{enumerate}
\end{lemma}

\begin{proof}
준연접 아이디얼층
$\mathcal{I} = \mathcal{I}_{D_1} + \mathcal{I}_{D_2}$가 스킴론적 교차
$D_1 \cap D_2 \subset X$를 정의함에 유의하자.
$Z$가 $D_1 \cap D_2$의 연결 성분들의 합집합이므로 모든 $z \in Z$에
대해 $\mathcal{O}_{X, z} \to \mathcal{O}_{Z, z}$의 핵은
$\mathcal{I}_z$와 같다. $b : X' \to X$를 $Z$에서 $X$의 부풀리기라
하자. (따라서 $Z$ 주위에서 Zariski 국소적으로 이는 $\mathcal{I}$에서
$X$를 부풀리기한 것이다.) 대응하는 유효 Cartier 인자를
$E = b^{-1}(Z)$로 나타내자. 『인자』의 보조정리
\ref{divisors-lemma-blowing-up-gives-effective-Cartier-divisor}를 보라.
$Z \subset D_1$이므로 $E \subset f^{-1}(D_1)$이고, 따라서 어떤 유효
Cartier 인자 $D'_1 \subset X'$에 대해 $D_1 = D_1' + E$이다.
『인자』의 보조정리
\ref{divisors-lemma-difference-effective-Cartier-divisors}를 보라.
마찬가지로 $D_2 = D_2' + E$이다. 이로써 주장 (1)--(5)가 처리되었다.

\medskip\noindent
$W'$이 (7)(a) 또는 (7)(b)와 같으면 $W'$의 상 $W$는
$D_1 \cap D_2$에 포함된다. $W$가 $Z$에 포함되지 않으면 $b$는 $W$의
일반점에서 동형이고, $\dim_\delta(W) = \dim_\delta(W') = n - 1$임을
얻는다. 이는
$\dim_\delta(D_1 \cap D_2 \setminus Z) \leq n - 2$라는 가정과
모순이다. 따라서 $W \subset Z$이다. 즉, (6)과 (7)을 증명하기 위해
$X$ 위에서 $Z$ 주위로 국소화해도 된다.

\medskip\noindent
따라서 $A$가 뇌터 정역이고 $X = \Spec(A)$이며,
$D_1 = \Spec(A/a)$, $D_2 = \Spec(A/b)$,
$Z = D_1 \cap D_2$라고 가정할 수 있다. $I = (a, b)$로 놓는다.
$A$가 정역이고 $a, b \not = 0$이므로 부풀리기를 두 조각
$U = \Spec(A[s]/(as - b))$와 $V = \Spec(A[t]/(bt -a))$로 덮을 수 있다.
이 조각들은 $s$를 $t^{-1}$로 보내는 동형
$A[s, s^{-1}]/(as - b) \cong A[t, t^{-1}]/(bt - a)$을 사용하여
붙인다. 유효 Cartier 인자 $E$는
$\Spec(A[s]/(as - b, a)) \subset U$와
$\Spec(A[t]/(bt - a, b)) \subset V$로 기술된다.
닫힌 부분스킴 $D'_1$은
$\Spec(A[t]/(bt - a, t)) \subset U$에 대응한다.
닫힌 부분스킴 $D'_2$는
$\Spec(A[s]/(as -b, s)) \subset V$에 대응한다.
``$ts = 1$''이므로 $D'_1 \cap D'_2 = \emptyset$이다.

\medskip\noindent
높이가 일이고 $s, a \in \mathfrak q$인 소 아이디얼
$\mathfrak q \subset A[s]/(as - b)$가 있다고 하자.
대응하는 $A$의 소 아이디얼을 $\mathfrak p \subset A$라 하자.
$a, b \in \mathfrak p$임에 유의하자. 차원 공식에 의해
$\dim(A_{\mathfrak p}) = 1$이기도 하다. 이제 보여야 할 마지막 주장은
$$
\text{ord}_{A_{\mathfrak p}}(a)
\text{ord}_{A_{\mathfrak p}}(b)
>
\text{ord}_{B_{\mathfrak q}}(a)
\text{ord}_{B_{\mathfrak q}}(s)
$$
이다. 여기서 $B = A[s]/(as - b)$이다. 『대수학』의 보조정리
\ref{algebra-lemma-quasi-finite-extension-dim-1}에 의해
$x = a, b$에 대해
$\text{ord}_{A_{\mathfrak p}}(x) \geq
\text{ord}_{B_{\mathfrak q}}(x)$이다.
$\text{ord}_{B_{\mathfrak q}}(s) > 0$이므로 $\text{ord}$ 함수의
가법성과 $as = b$라는 사실을 사용하면 결론이 나온다.
\end{proof}

\begin{definition}
\label{obsolete-definition-locally-finite-sum-effective-Cartier-divisors}
$X$를 스킴이라 하고 $\{D_i\}_{i \in I}$를 $X$ 위의 유효 Cartier
인자들의 국소 유한 모임이라 하자. 함수
$I \to \mathbf{Z}_{\geq 0}$, $i \mapsto n_i$가 주어졌다고 하자.
{\it 유효 Cartier 인자들의 합}
$D = \sum n_i D_i$는 다음 성질을 갖는 유일한 유효 Cartier 인자
$D \subset X$이다. 임의의 준콤팩트 열린집합 $U \subset X$ 위에서
$D|_U = \sum_{D_i \cap U \not = \emptyset} n_iD_i|_U$이고,
이는 『인자』의 정의
\ref{divisors-definition-sum-effective-Cartier-divisors}에서와 같은 합이다.
\end{definition}

\begin{lemma}
\label{obsolete-lemma-sum-divisors-associated-Weil}
$(S, \delta)$가 『Chow 호몰로지』의 상황
\ref{chow-situation-setup}에서와 같다고 하자.
$X$를 $S$ 위에서 국소 유한형이라 하고, $X$가 정수적이며
$\dim_\delta(X) = n$이라고 가정하자.
$\{D_i\}_{i \in I}$를 $X$ 위의 유효 Cartier 인자들의 국소 유한
모임이라 하자. $i \in I$에 대해 $n_i \geq 0$이 주어졌다고 하자.
그러면 $Z_{n - 1}(X)$에서
$$
[D]_{n - 1} = \sum\nolimits_i n_i[D_i]_{n - 1}
$$
이다.
\end{lemma}

\begin{proof}
순환들의 등식을 증명하므로 $X$ 위에서 국소적으로 작업해도 된다.
따라서 유한합의 경우로, 다시 귀납법을 사용하여 두 유효 Cartier 인자의
합 $D = D_1 + D_2$의 경우로 환원된다.
『Chow 호몰로지』의 보조정리 \ref{chow-lemma-compute-c1}에 의해
$D_1 = \text{div}_{\mathcal{O}_X(D_1)}(1_{D_1})$이다. 여기서
$1_{D_1}$은 $\mathcal{O}_X(D_1)$의 표준 절단이다. 물론 $D_2$와
$D$에 대해서도 같은 명제가 성립한다. 동일시
$\mathcal{O}_X(D) = \mathcal{O}_X(D_1) \otimes \mathcal{O}_X(D_2)$를
통해 $1_D = 1_{D_1} \otimes 1_{D_2}$이므로 『인자』의 보조정리
\ref{divisors-lemma-c1-additive}를 적용하면 된다.
\end{proof}

\begin{lemma}
\label{obsolete-lemma-blowing-up-intersections}
$(S, \delta)$가 『Chow 호몰로지』의 상황
\ref{chow-situation-setup}에서와 같다고 하자.
$X$를 $S$ 위에서 국소 유한형이라 하고, $X$가 정수적이며
$\dim_\delta(X) = d$라고 가정하자.
$\{D_i\}_{i \in I}$를 $X$ 위의 유효 Cartier 인자들의 국소 유한
모임이라 하자. 모든 $\{i, j, k\} \subset I$,
$\#\{i, j, k\} = 3$에 대해
$D_i \cap D_j \cap D_k = \emptyset$이라고 가정하자.
그러면 다음이 존재한다.
\begin{enumerate}
\item $\dim_\delta(X \setminus U) \leq d - 3$인 열린 부분스킴
$U \subset X$,
\item 사상 $b : U' \to U$,
\item $U'$ 위의 유효 Cartier 인자들 $\{D'_j\}_{j \in J}$.
\end{enumerate}
이들은 다음 성질을 갖는다.
\begin{enumerate}
\item $b$는 고유 사상 $b : U' \to U$이고,
\item $U'$은 정수적이며,
\item $b$는 $D_i|_U$들의 두 개씩의 교차들의 합집합의 여집합 위에서
동형이고,
\item $\{D'_j\}_{j \in J}$는 $U'$ 위의 유효 Cartier 인자들의
국소 유한 모임이며,
\item $j \not = j'$이면
$\dim_\delta(D'_j \cap D'_{j'}) \leq d - 2$이고,
\item 어떤 $n_{ij} \geq 0$에 대해
$b^{-1}(D_i|_U) = \sum n_{ij} D'_j$이다.
\end{enumerate}
또한 $X$가 준콤팩트이면 위에서 $U = X$라고 가정할 수 있다.
\end{lemma}

\begin{proof}
먼저 아마도 가장 흥미로운 준콤팩트 경우를 증명하자. 이 경우 부풀리기의
열
$$
X = X_0 \xleftarrow{b_0} X_1 \xleftarrow{b_1} X_2 \leftarrow \ldots
$$
과 유효 Cartier 인자들의 유한 집합
$\{D_{n, i}\}_{i \in I_n}$을 귀납적으로 만든다.
각 단계에서 모든 삼중 교차
$D_{n, i} \cap D_{n, j} \cap D_{n, k}$가 공집합이라는 성질을 갖게 한다.
또한 각 $n \geq 0$에 대해
$I_{n + 1} = I_n \amalg P(I_n)$이 되게 한다. 여기서 $P(I_n)$은
$I_n$의 원소들의 쌍들의 집합이다. 마지막으로
$$
b_n^{-1}(D_{n, i}) = D_{n + 1, i} +
\sum\nolimits_{i' \in I_n, i' \not = i} D_{n + 1, \{i, i'\}}
$$
가 되게 한다. 이로부터 각 $n \geq 0$에 대해
$(b_0 \circ \ldots \circ b_n)^{-1}(D_i)$가
$j \in I_{n + 1}$에 대한 인자 $D_{n + 1, j}$들의 음이 아닌 정수
선형결합임을 얻는다.

\medskip\noindent
귀납법을 시작하기 위해 $X_0 = X$, $I_0 = I$,
$D_{0, i} = D_i$로 놓는다.

\medskip\noindent
$(X_n, \{D_{n, i}\}_{i \in I_n})$이 주어졌다고 하자.
$X_{n + 1}$을 닫힌 부분스킴
$Z_n = \bigcup_{\{i, i'\} \in P(I_n)} D_{n, i} \cap D_{n, i'}$에서
$X_n$을 부풀리기한 것으로 정의한다.
삼중 교차에 관한 가정에 의해 닫힌 부분스킴
$D_{n, i} \cap D_{n, i'}$들이 서로소임에 유의하자. 즉,
$Z_n = \coprod_{\{i, i'\} \in P(I_n)} D_{n, i} \cap D_{n, i'}$라고
쓸 수 있다. 또한 $D_{n, i} \cap D_{n, i'}$의 Zariski 근방에서
사상 $b_n$은 닫힌 부분스킴 $D_{n, i} \cap D_{n, i'}$에서 스킴
$X_n$을 부풀리기한 것과 같고, 보조정리 \ref{obsolete-lemma-two-divisors}의
결과들을 적용할 수 있다. 따라서
$D_{n + 1, \{i, i'\}} = b_n^{-1}(D_i \cap D_{i'})$로 놓으면 유효
Cartier 인자를 얻는다. Cartier 인자
$D_{n + 1, \{i, i'\}}$들은 서로소이다.
모든 $i' \in I_n$, $i' \not = i$에 대해
$b_n^{-1}(D_{n, i}) \supset D_{n + 1, \{i, i'\}}$임은 명백하다.
그러므로 『인자』의 보조정리
\ref{divisors-lemma-difference-effective-Cartier-divisors}를 적용하면
$X_{n + 1}$ 위의 어떤 유효 Cartier 인자 $D_{n + 1, i}$에 대해 실제로
$b^{-1}(D_{n, i}) = D_{n + 1, i} +
\sum\nolimits_{i' \in I_n, i' \not = i} D_{n + 1, \{i, i'\}}$이다.
$D_{n + 1, \{i, i'\}}$의 근방에서 이 인자들
$D_{n + 1, i}$는 보조정리 \ref{obsolete-lemma-two-divisors}의 프라임이 붙은
인자들의 역할을 한다. 특히 그 보조정리의 (6)에 의해
$i \not = i'$, $i, i' \in I_n$이면
$D_{n + 1, i} \cap D_{n + 1, i'} = \emptyset$이다.
보조정리 \ref{obsolete-lemma-two-divisors}의 (6)에 의해 이는 이미 인자
$D_{n + 1, i}$들의 삼중 교차가 영임을 뜻한다.

\medskip\noindent
좋다. 이제 $X$의 준콤팩트성을 사용하여 불변량
\begin{equation}
\label{obsolete-equation-invariant}
\epsilon(X, \{D_i\}_{i \in I}) =
\max\{\epsilon_Z(D_i, D_{i'}) \mid
Z \subset X,
\dim_\delta(Z) = d - 1,
\{i, i'\} \in P(I)\}
\end{equation}
이 유한임을 알 수 있다. 실제로 각 $D_i$는 기약 성분을 유한 개 이하만
갖는다. 어떤 $n$에 대해 불변량
$\epsilon(X_n, \{D_{n, i}\}_{i \in I_n})$이 영이라고 주장한다.
그렇지 않다면 보조정리 \ref{obsolete-lemma-two-divisors}에 의해 양의 정수들의
엄밀히 감소하는 열
$$
\epsilon(X, \{D_i\}_{i \in I})
=
\epsilon(X_0, \{D_{0, i}\}_{i \in I_0})
>
\epsilon(X_1, \{D_{1, i}\}_{i \in I_1})
>
\ldots
$$
을 얻어 모순이다. 불변량
$\epsilon(X_n, \{D_{n, i}\}_{i \in I_n})$이 영인 $n$을 택한다.
이는 $\epsilon_Z(D_{n, i}, D_{n, i'}) > 0$인 정수적 닫힌 부분스킴
$Z \subset X_n$과 지표 쌍 $i, i' \in I_n$이 없다는 뜻이다.
즉, 원하는 대로 모든 $\{i, i'\} \in P(I_n)$에 대해
$\dim_\delta(D_{n, i}, D_{n, i'}) \leq d - 2$이다.

\medskip\noindent
다음으로 스킴 $X$가 더 이상 준콤팩트라고 가정하지 않는 일반적인 경우를
다룬다. 증명의 첫 부분의 아이디어에는 $X$의 어떤 고정된 점을 우세하게
덮는 중심을 갖는 부풀리기의 무한 열이 생길 수 있다는 문제가 있다.
이를 피하기 위해 각 단계에서 여차원이 $\geq 3$인 적절한 닫힌 부분집합을
잘라 낸다. 구체적으로 다음 꼴의 사상들의 열을 귀납적으로 구성한다.
$$
\xymatrix{
X = X_0 \\
U_0 \ar[u]^{j_0} & X_1 \ar[l]_{b_0} \\
 & U_1 \ar[u]^{j_1} & X_2 \ar[l]_{b_1} \\
 & & U_2 \ar[u]^{j_2} & X_3 \ar[l]_{b_2}
}
$$
각 사상 $j_n : U_n \to X_n$은 열린 몰입이다.
각 사상 $b_n : X_{n + 1} \to U_n$은 정수적 스킴의 고유 쌍유리
사상이다. 준콤팩트 경우와 마찬가지로 $X_n$ 위의 유효 Cartier 인자
$\{D_{n, i}\}_{i \in I_n}$를 갖게 한다. 각 단계에서 모든 삼중 교차
$D_{n, i} \cap D_{n, j} \cap D_{n, k}$가 공집합이라는 성질을 갖게 한다.
또한 각 $n \geq 0$에 대해
$I_{n + 1} = I_n \amalg P(I_n)$이 되게 한다. 여기서 $P(I_n)$은
$I_n$의 원소들의 쌍들의 집합이다. 마지막으로
$$
b_n^{-1}(D_{n, i}|_{U_n}) = D_{n + 1, i} +
\sum\nolimits_{i' \in I_n, i' \not = i} D_{n + 1, \{i, i'\}}
$$
가 되게 한다.

\medskip\noindent
$X_0 = X$, $I_0 = I$, $D_{0, i} = D_i$로 놓아 귀납법을 시작한다.

\medskip\noindent
$(X_n, \{D_{n, i}\})$이 주어졌을 때 열린 부분스킴 $U_n$을 다음과 같이
구성한다. 각 쌍 $\{i, i'\} \in P(I_n)$에 대해 닫힌 부분스킴
$D_{n, i} \cap D_{n, i'}$을 생각하자. 이 부분스킴에는 $\delta$-차원이
$d - 2$인 ``좋은'' 기약 성분들과 $\delta$-차원이 $d - 1$인 ``나쁜''
기약 성분들이 있다. 다음과 같이 놓는다.
$$
\text{Bad}(i, i')
=
\bigcup\nolimits_{W \subset D_{n, i} \cap D_{n, i'}
\text{ 다음 조건의 기약 성분 }\dim_\delta(W) = d - 1} W
$$
마찬가지로
$$
\text{Good}(i, i')
=
\bigcup\nolimits_{W \subset D_{n, i} \cap D_{n, i'}
\text{ 다음 조건의 기약 성분 }\dim_\delta(W) = d - 2} W.
$$
그러면
$D_{n, i} \cap D_{n, i'} = \text{Bad}(i, i') \cup \text{Good}(i, i')$이고
$\dim_\delta(\text{Bad}(i, i') \cap \text{Good}(i, i')) \leq d - 3$이다.
$U_n$을 다음과 같이 택한다.
$$
U_n
=
X_n
\setminus
\bigcup\nolimits_{\{i, i'\} \in P(I_n)}
\text{Bad}(i, i') \cap \text{Good}(i, i').
$$
인자 $D_{n, i}$들의 삼중 교차에 관한 조건에 의해 이 합집합은 실제로
서로소 합집합이다. 또한 스킴으로서
$$
D_{n, i}|_{U_n} \cap D_{n, i'}|_{U_n}
=
Z_{n, i, i'} \amalg G_{n, i, i'}
$$
임을 알 수 있다. 여기서 $Z_{n, i, i'}$는 $\delta$-차원 $d - 1$의
등차원이고, $G_{n, i, i'}$는 $\delta$-차원 $d - 2$의 등차원이다.
(따라서 위상적으로 $Z_{n, i, i'}$는 나쁜 성분들의 합집합에서 좋은
성분들과의 교차를 버린 것이다.) 마지막으로
$$
Z_n =
\bigcup\nolimits_{\{i, i'\} \in P(I_n)} Z_{n, i, i'} =
\coprod\nolimits_{\{i, i'\} \in P(I_n)} Z_{n, i, i'},
$$
로 놓고 $b_n : X_{n + 1} \to X_n$을 $Z_n$에서의 부풀리기로 정의한다.
보조정리 \ref{obsolete-lemma-two-divisors}는 각 자취
$D_{n, i}|_{U_n} \cap D_{n, i'}|_{U_n}$의 근방에서 사상
$b_n : X_{n + 1} \to X_n$에 적용된다는 점에 유의하자.
따라서 증명의 첫 부분과 정확히 같은 방식으로
$\{i, i'\} \in P(I_n)$에 대한 유효 Cartier 인자
$D_{n + 1, \{i, i'\}}$와 $i \in I_n$에 대한 유효 Cartier 인자
$D_{n + 1, i}$를 얻으며, 이들은
$b_n^{-1}(D_{n, i}|_{U_n}) = D_{n + 1, i} +
\sum\nolimits_{i' \in I_n, i' \not = i} D_{n + 1, \{i, i'\}}$를
만족한다. 각 $n$에 대해 합성
$j_0 \circ \ldots \circ j_{n - 1} \circ b_{n - 1}$로 얻는 사상을
$\pi_n : X_n \to X$로 나타내자.

\medskip\noindent
{\bf 주장:} 임의의 준콤팩트 열린집합 $V \subset X$가 주어지면 충분히
큰 모든 $n$에 대해 사상
$$
\pi_n^{-1}(V) \leftarrow \pi_{n + 1}^{-1}(V) \leftarrow \ldots
$$
은 모두 동형이다. 실제로 사상
$\pi_n^{-1}(V) \leftarrow \pi_{n + 1}^{-1}(V)$이 동형이 아니면 어떤
$\{i, i'\} \in P(I_n)$에 대해
$Z_{n, i, i'} \cap \pi_n^{-1}(V) \not = \emptyset$이다.
따라서 $\dim_\delta(W) = d - 1$인 기약 성분
$W \subset D_{n, i} \cap D_{n, i'}$이 존재한다. 특히
$\epsilon_W(D_{n, i}, D_{n, i'}) > 0$이다.
보조정리 \ref{obsolete-lemma-two-divisors}를 반복해서 적용하면
$$
\epsilon_W(D_{n, i}, D_{n, i'})
<
\epsilon(V, \{D_i|_V\}) - n
$$
을 얻는다. 여기서 $\epsilon(V, \{D_i|_V\})$는
(\ref{obsolete-equation-invariant})에서와 같다. $V$가 준콤팩트이므로
$\epsilon(V, \{D_i|_V\}) < \infty$이고,
$n > \epsilon(V, \{D_i|_V\})$을 택하면 결과를 얻는다.

\medskip\noindent
구성에 의해 차집합 $X_n \setminus U_n$은
$\dim_\delta(X_n \setminus U_n) \leq d - 3$임에 유의하자.
$T_n = \pi_n(X_n \setminus U_n)$을 그 $X$ 안의 상이라 하자.
위 사상 도표를 따라가면서 각 $b_n$이 닫혀 있음을 사용하면, 모든
$n$에 대해 $T_0 \cup \ldots \cup T_n$이 $X$의 닫힌 부분집합임이
따른다. 구성에 의해 모든 $t \in T_n$은
$\delta(t) \leq d - 3$을 만족한다. 따라서
$\overline{T_n} \subset X$는
$\dim_\delta(T_n) \leq d - 3$인 닫힌 부분집합이다.
위 주장에 의해 임의의 준콤팩트 열린집합 $V \subset X$에 대해
$T_n \cap V \not = \emptyset$인 $n$은 유한 개뿐이다.
따라서 $\{\overline{T_n}\}_{n \geq 0}$은 닫힌 부분집합들의 국소 유한
모임이고, $U = X \setminus \bigcup \overline{T_n}$으로 놓을 수 있다.
이것이 보조정리에 나오는 $U$가 된다.

\medskip\noindent
$U$의 구성에 의해
$U_n \cap \pi_n^{-1}(U) = \pi_n^{-1}(U)$임에 유의하자.
따라서 모든 사상
$$
b_n : \pi_{n + 1}^{-1}(U) \longrightarrow \pi_n^{-1}(U)
$$
은 고유이다. 또한 위 주장에 의해 이 사상들은 $X$의 각 준콤팩트
열린집합 위에서 결국 동형이 된다. 따라서
$$
U' = \lim_n \pi_n^{-1}(U).
$$
를 정의할 수 있다. 이는 $U$ 위에서 국소적인 성질이고 각 콤팩트
열린집합 위에서 극한이 안정되므로, 유도된 사상 $b : U' \to U$는
고유이다. 마찬가지로 구성에서의 포함
$I_n \to I_{n + 1}$을 사용하여 $J = \bigcup_{n \geq 0} I_n$으로 놓는다.
$j \in J$에 대해 $j$가 $i \in I_{n_0}$에 대응하는 $n_0$를 택하고
$D'_j = \lim_{n \geq n_0} D_{n, i}$로 정의한다. $X$ 위에서 국소적으로
사상들이 안정되므로 이 역시 잘 정의된다. 보조정리의 나머지 주장들은
준콤팩트 $X$의 경우와 같이 확인한다.
\end{proof}

\section{교차하는 인자들의 가환성}
\label{obsolete-section-commutativity}

\noindent
이 절의 결과들은 원래 『Chow 호몰로지』의 절
\ref{chow-section-commutativity}의 보조정리들에 대한 다른 증명과
『Chow 호몰로지』의 보조정리
\ref{chow-lemma-gysin-commutes-gysin}의 약한 판을 제시하는 데 사용되었다.

\begin{lemma}
\label{obsolete-lemma-improved-additivity}
$(S, \delta)$가 『Chow 호몰로지』의 상황
\ref{chow-situation-setup}에서와 같다고 하자.
$X$를 $S$ 위에서 국소 유한형이라 하자.
$\{i_j : D_j \to X \}_{j \in J}$를 $X$ 위의 유효 Cartier 인자들의
국소 유한 모임이라 하자. $j\in J$에 대해 $n_j > 0$이라 하자.
$D = \sum_{j \in J} n_j D_j$로 놓고 포함사상을 $i : D \to X$로
나타내자. $\alpha \in Z_{k + 1}(X)$라 하자. 그러면
$$
p : \coprod\nolimits_{j \in J} D_j \longrightarrow D
$$
는 고유이고 $\CH_k(D)$에서
$$
i^*\alpha = p_*\left(\sum n_j i_j^*\alpha\right)
$$
이다.
\end{lemma}

\begin{proof}
유리 동치들의 무한합에 관한 미묘한 점 때문에 이 보조정리의 증명은
예상보다 조금 길어진다. 준콤팩트 경우에는 족 $D_j$가 유한이고,
결과는 매우 쉬우며 『Chow 호몰로지』의 보조정리
\ref{chow-lemma-compute-c1}, 『인자』의 보조정리
\ref{divisors-lemma-c1-additive}, 그리고 정의들에서 곧바로 따른다.

\medskip\noindent
족 $\{D_j\}_{j \in J}$가 국소 유한이므로 사상 $p$는 고유이다.
$W_a \subset X$가 $\delta$-차원 $k + 1$의 정수적 닫힌 부분스킴이 되도록
$\alpha = \sum_{a \in A} m_a [W_a]$라고 쓰자.
닫힌 몰입을 $i_a : W_a \to X$로 나타내자.
$\{W_a\}_{a \in A}$가 $X$ 위에서 국소 유한이고 모든 $a \in A$에 대해
$m_a \not = 0$이라고 가정한다.

\medskip\noindent
『Chow 호몰로지』의 정의
\ref{chow-definition-gysin-homomorphism}에 의해 류 $i^*\alpha$는 어떤
$\beta_a \in Z_k(W_a \cap D)$에 대한 순환
$\sum m_a\beta_a$의 류임에 유의하자. 구체적으로
$W_a \not \subset D$이면 $\beta_a = [D \cap W_a]_k$이고,
$W_a \subset D$이면 $\beta_a$는
$c_1(\mathcal{O}_X(D)) \cap [W_a]$를 나타내는 순환이다.

\medskip\noindent
각 $a \in A$에 대해 $J = J_{a, 1} \amalg J_{a, 2} \amalg J_{a, 3}$로
쓰자. 여기서
\begin{enumerate}
\item $j \in J_{a, 1}$일 필요충분조건은
$W_a \cap D_j = \emptyset$이고,
\item $j \in J_{a, 2}$일 필요충분조건은
$W_a \not = W_a \cap D_1 \not = \emptyset$이며,
\item $j \in J_{a, 3}$일 필요충분조건은 $W_a \subset D_j$이다.
\end{enumerate}
족 $\{D_j\}$가 국소 유한이므로 $J_{a, 3}$은 유한 집합이다.
각 $a \in A$와 $j \in J$에 대해 순환
$\beta_{a, j} \in Z_k(W_a \cap D_j)$를 다음과 같이 택한다.
\begin{enumerate}
\item $j \in J_{a, 1}$이면 $\beta_{a, j} = 0$으로 놓고,
\item $j \in J_{a, 2}$이면
$\beta_{a, j} = [D_j \cap W_a]_k$로 놓으며,
\item $j \in J_{a, 3}$이면
$c_1(i_a^*\mathcal{O}_X(D_j)) \cap [W_j]$를 나타내는
$\beta_{a, j} \in Z_k(W_a)$를 택한다.
\end{enumerate}
$\CH_k(W_a \cap D)$에서
$$
\beta_a \sim_{rat}
\sum\nolimits_{j \in J} n_j \beta_{a, j}
$$
라고 주장한다.

\medskip\noindent
경우 I: $W_a \not \subset D$. 이 경우 $J_{a, 3} = \emptyset$이다.
따라서 순환들의 등식
$[D \cap W_a]_k = \sum n_j [D_j \cap W_a]_k$을 보이면 충분하다.
이는 보조정리 \ref{obsolete-lemma-sum-divisors-associated-Weil}이다.

\medskip\noindent
경우 II: $W_a \subset D$. 이 경우 $\beta_a$는
$c_1(i_a^*\mathcal{O}_X(D)) \cap [W_a]$를 나타내는 순환이다.
$D = D_{a, 1} + D_{a, 2} + D_{a, 3}$으로 쓰자. 여기서
$D_{a, s} = \sum_{j \in J_{a, s}} n_jD_j$이다.
『인자』의 보조정리 \ref{divisors-lemma-c1-additive}에 의해
\begin{eqnarray*}
c_1(i_a^*\mathcal{O}_X(D)) \cap [W_a] & = &
c_1(i_a^*\mathcal{O}_X(D_{a, 1})) \cap [W_a] +
c_1(i_a^*\mathcal{O}_X(D_{a, 2})) \cap [W_a] \\
& &
 + c_1(i_a^*\mathcal{O}_X(D_{a, 3})) \cap [W_a].
\end{eqnarray*}
합의 첫 번째 항은 영임이 명백하다. $J_{a, 3}$이 유한이므로 마지막 항은
$\sum\nolimits_{j \in J_{a, 3}} n_jc_1(i_a^*\mathcal{L}_j) \cap [W_a]$
와 일치한다. 『인자』의 보조정리
\ref{divisors-lemma-c1-additive}를 보라. 이는
$\sum_{j \in J_{a, 3}} n_j \beta_{a, j}$로 나타내어진다.
마지막으로 경우 I에 의해 가운데 항은 순환
$\sum\nolimits_{j \in J_{a, 2}} n_j[D_j \cap W_a]_k =
\sum_{j \in J_{a, 2}} n_j\beta_{a, j}$로 나타내어진다.
따라서 이 경우에도 주장이 성립한다.

\medskip\noindent
이제 보조정리의 증명을 끝낼 준비가 되었다. $\beta_a$를 택한 방식에 의해
$i^*D \sim_{rat} \sum m_a\beta_a$이다. 각 $a$에 대해
$\beta_a \sim_{rat} \sum_j \beta_{a, j}$이고, 이 유리 동치는
$D \cap W_a$ 위에서 이루어진다. 닫힌 부분스킴 $D \cap W_a$들의 모임이
$D$ 위에서 국소 유한이므로 $D$ 위에서도
$\sum m_a \beta_a \sim_{rat} \sum_{a, j} m_a\beta_{a, j}$이다!
(『Chow 호몰로지』의 주석
\ref{chow-remark-infinite-sums-rational-equivalences}를 보라.)
좋다. 이제 $D_j$ 위의 순환으로 본 $\sum_a m_a\beta_{a, j}$는
$i_j^*\alpha$를 나타내므로 $\sum_{a, j} m_a\beta_{a, j}$는
$p_* \sum_j i_j^*\alpha$를 나타낸다. 이것으로 끝이다.
\end{proof}

\begin{lemma}
\label{obsolete-lemma-commutativity-effective-Cartier-proper-intersection}
$(S, \delta)$가 『Chow 호몰로지』의 상황
\ref{chow-situation-setup}에서와 같다고 하자.
$X$를 $S$ 위에서 국소 유한형이라 하고, $X$가 정수적이며
$\dim_\delta(X) = n$이라고 가정하자.
$D$, $D'$를 $X$ 위의 유효 Cartier 인자라 하자.
$\dim_\delta(D \cap D') = n - 2$라고 가정하자.
대응하는 닫힌 몰입을 각각 $i : D \to X$와 $i' : D' \to X$라 하자.
그러면
\begin{enumerate}
\item 순환 $\alpha \in Z_{n - 2}(D \cap D')$가 존재하여,
$D$로의 밂앞은 $i^*[D']_{n - 1} \in \CH_{n - 2}(D)$를 나타내고
$D'$으로의 밂앞은
$(i')^*[D]_{n - 1} \in \CH_{n - 2}(D')$를 나타내며,
\item $\CH_{n - 2}(X)$에서
$$
D \cdot [D']_{n - 1}
=
D' \cdot [D]_{n - 1}
$$
이다.
\end{enumerate}
\end{lemma}

\begin{proof}
(2)는 (1)에서 자명하게 따른다.
$Z_a$를 $D$의 기약 성분, $[Z_b]$를 $D'$의 기약 성분으로 하여
$[D]_{n - 1} = \sum n_a[Z_a]$와
$[D']_{n - 1} = \sum m_b[Z_b]$라고 쓰자.
『Chow 호몰로지』의 정의
\ref{chow-definition-gysin-homomorphism}에 의해
$i^*D' = \sum m_b i^*[Z_b]$이고
$(i')^*D = \sum n_a(i')^*[Z_a]$이다.
가정에 의해 어떤 기약 성분 $Z_b$도 $D$에 포함되지 않으므로 정의에 따라
$i^*[Z_b] = [Z_b\cap D]_{n - 2}$이다. 마찬가지로
$(i')^*[Z_a] = [Z_a \cap D']_{n - 2}$이다. 따라서 증명하려는 것은
$D \cap D'$ 위에 실제로 지지되는 순환들의 등식
$$
\sum n_a[Z_a \cap D']_{n - 2} = \sum m_b[Z_b \cap D]_{n - 2}
$$
이다. $W \subset X$를 $\dim_\delta(W) = n - 2$인 정수적 닫힌
부분스킴이라 하고 $\xi \in W$를 일반점이라 하자.
$R = \mathcal{O}_{X, \xi}$로 놓는다. 이는 뇌터 국소 정역이고
$\dim(R) = 2$이다. $D$의 아이디얼을 정의하는 원소를 $f \in R$,
$D'$의 아이디얼을 정의하는 원소를 $f' \in R$라 하자.
가정에 의해 $\dim(R/(f, f')) = 0$이다.
$(f')$ 위의 극소 소 아이디얼들을
$\mathfrak q'_1, \ldots, \mathfrak q'_t \subset R$라 하고,
$(f)$ 위의 극소 소 아이디얼들을
$\mathfrak q_1, \ldots, \mathfrak q_s \subset R$라 하자.
위 등식은 다음 등식으로 귀착된다.
$$
\sum_{i = 1, \ldots, s}
\text{length}_{R_{\mathfrak q_i}}(R_{\mathfrak q_i}/(f))
\text{ord}_{R/\mathfrak q_i}(f')
=
\sum_{j = 1, \ldots, t}
\text{length}_{R_{\mathfrak q'_j}}(R_{\mathfrak q'_j}/(f'))
\text{ord}_{R/\mathfrak q'_j}(f).
$$
『Chow 호몰로지』의 보조정리
\ref{chow-lemma-length-multiplication}를 $M = R/(f)$에 적용하면
이 등식의 왼쪽은
$$
\text{length}_R(R/(f, f'))
-
\text{length}_R(\Ker(f' : R/(f) \to R/(f)))
$$
과 같다. 좋다. 이제 사상
$x \bmod (f) \mapsto f'x \bmod (ff')$을 통해
$\Ker(f' : R/(f) \to R/(f))$가
$((f) \cap (f'))/(ff')$와 표준적으로 동형임에 유의하자.
따라서 왼쪽은
$$
\text{length}_R(R/(f, f'))
-
\text{length}_R((f) \cap (f')/(ff'))
$$
이다. 이는 $f$와 $f'$에 관해 대칭이므로 결론이 나온다.
\end{proof}

\begin{lemma}
\label{obsolete-lemma-commutativity-effective-Cartier-proper-intersection-infinite}
$(S, \delta)$가 『Chow 호몰로지』의 상황
\ref{chow-situation-setup}에서와 같다고 하자.
$X$를 $S$ 위에서 국소 유한형이라 하고, $X$가 정수적이며
$\dim_\delta(X) = n$이라고 가정하자.
$\{D_j\}_{j \in J}$를 $X$ 위의 유효 Cartier 인자들의 국소 유한
모임이라 하자. $n_j, m_j \geq 0$을 음이 아닌 정수들의 모임이라 하자.
$D = \sum n_j D_j$, $D' = \sum m_j D_j$로 놓는다.
모든 $j \not = j'$에 대해
$\dim_\delta(D_j \cap D_{j'}) = n - 2$라고 가정하자.
그러면 $\CH_{n - 2}(X)$에서
$D \cdot [D']_{n - 1} = D' \cdot [D]_{n - 1}$이다.
\end{lemma}

\begin{proof}
합들이 유한인 경우, 예를 들어 $X$가 준콤팩트인 경우에는 이 보조정리는
보조정리 \ref{obsolete-lemma-sum-divisors-associated-Weil}과
\ref{obsolete-lemma-commutativity-effective-Cartier-proper-intersection}에서
자명하게 따른다. 따라서 독자에게 증명을 건너뛰기를 권한다.

\medskip\noindent
일반적인 경우의 증명은 다음과 같다.
$i_j : D_j \to X$를 닫힌 몰입이라 하자.
$p : \coprod D_j \to X$를 사상 $i_j$들의 쌍대곱이라 하자.
$\{Z_a\}_{a \in A}$를 $\bigcup D_j$의 기약 성분들의 모임이라 하자.
각 $j$에 대해
$$
[D_j]_{n - 1} = \sum d_{j, a}[Z_a].
$$
라고 쓰자. 보조정리 \ref{obsolete-lemma-sum-divisors-associated-Weil}에 의해
$$
[D]_{n - 1} = \sum n_j d_{j, a} [Z_a],
\quad
[D']_{n - 1} = \sum m_j d_{j, a} [Z_a].
$$
보조정리 \ref{obsolete-lemma-improved-additivity}에 의해
$$
D \cdot [D']_{n - 1} = p_*\left(\sum n_j i_j^*[D']_{n - 1} \right),
\quad
D' \cdot [D]_{n - 1} = p_*\left(\sum m_{j'} i_{j'}^*[D]_{n - 1} \right).
$$
Gysin 준동형의 정의에서와 같이(『Chow 호몰로지』의 정의
\ref{chow-definition-gysin-homomorphism}를 보라), $D_j \cap Z_a$ 위에서
$i_j^*[Z_a]$를 나타내는 순환 $\beta_{a, j}$를 택한다.
(실제로 $Z_a$가 $D_j$에 포함되지 않으면
$\beta_{a, j} = [D_j \cap Z_a]_{n - 2}$이고, 따라서 이 경우에는 선택의
여지가 없음에 유의하자.) 각 $D_j$로 제한했을 때 $p$가 닫힌 몰입이므로
$\beta_{a, j}$를 $X$ 위의 순환으로 볼 수 있고, 그렇게 하겠다.
위에서 얻은 $[D]_{n - 1}$과 $[D']_{n - 1}$의 공식들을 대입하면
$$
D \cdot [D']_{n - 1} =
\sum\nolimits_{j, j', a} n_j m_{j'} d_{j', a} \beta_{a, j},
\quad
D' \cdot [D]_{n - 1} =
\sum\nolimits_{j, j', a} m_{j'} n_j d_{j, a} \beta_{a, j'}.
$$
또한 같은 규약 아래
$$
D_j \cdot [D_{j'}]_{n - 1} = \sum d_{j', a} \beta_{a, j}.
$$
이 용어로 보조정리
\ref{obsolete-lemma-commutativity-effective-Cartier-proper-intersection}는
(그 증명도 보라) $j \not = j'$일 때 순환
$\sum d_{j', a} \beta_{a, j}$와
$\sum d_{j, a} \beta_{a, j'}$가 순환으로서 같다고 말한다!
따라서
\begin{eqnarray*}
D \cdot [D']_{n - 1}
& = &
\sum\nolimits_{j, j', a} n_j m_{j'} d_{j', a} \beta_{a, j} \\
& = &
\sum\nolimits_{j \not = j'} n_j m_{j'}
\left(\sum\nolimits_a d_{j', a} \beta_{a, j}\right) +
\sum\nolimits_{j, a} n_j m_j d_{j, a} \beta_{a, j} \\
& = &
\sum\nolimits_{j \not = j'} n_j m_{j'}
\left(\sum\nolimits_a d_{j, a} \beta_{a, j'}\right) +
\sum\nolimits_{j, a} n_j m_j d_{j, a} \beta_{a, j} \\
& = &
\sum\nolimits_{j, j', a} m_{j'} n_j d_{j, a} \beta_{a, j'} \\
& = &
D' \cdot [D]_{n - 1}
\end{eqnarray*}
이고, 결론이 나온다.
\end{proof}

\begin{lemma}
\label{obsolete-lemma-commutativity-effective-Cartier}
$(S, \delta)$가 『Chow 호몰로지』의 상황
\ref{chow-situation-setup}에서와 같다고 하자.
$X$를 $S$ 위에서 국소 유한형이라 하고, $X$가 정수적이며
$\dim_\delta(X) = n$이라고 가정하자.
$D$, $D'$를 $X$ 위의 유효 Cartier 인자라 하자. 그러면
$$
D \cdot [D']_{n - 1} = D' \cdot [D]_{n - 1}
$$
이 $\CH_{n - 2}(X)$에서 성립한다.
\end{lemma}

\begin{proof}[보조정리 \ref{obsolete-lemma-commutativity-effective-Cartier}의 첫 번째 증명]
먼저 $X$가 준콤팩트인 경우를 증명하자. $X$와 두 원소로 된 유효
Cartier 인자들의 집합 $\{D, D'\}$에 보조정리
\ref{obsolete-lemma-blowing-up-intersections}를 적용한다. 그러면 고유 사상
$b : X' \to X$와 서로 두 개씩 여차원 $\geq 2$에서 교차하는 유효
Cartier 인자들의 유한 모임 $D'_j \subset X'$을 얻으며,
$b^{-1}(D) = \sum n_j D'_j$이고
$b^{-1}(D') = \sum m_j D'_j$이다.
보조정리 \ref{obsolete-lemma-push-pull-effective-Cartier}에 의해
$Z_{n - 1}(X)$에서
$b_*[b^{-1}(D)]_{n - 1} = [D]_{n - 1}$이고 $D'$에 대해서도 마찬가지이다.
따라서 『Chow 호몰로지』의 보조정리
\ref{chow-lemma-pushforward-cap-c1}에 의해 $\CH_{n - 2}(X)$에서
$$
D \cdot [D']_{n - 1} = b_*\left(b^{-1}(D) \cdot
[b^{-1}(D')]_{n - 1}\right)
$$
이고 다른 항도 마찬가지이다. 그러므로 보조정리
\ref{obsolete-lemma-commutativity-effective-Cartier-proper-intersection-infinite}의
$\CH_{n - 2}(X')$에서의 등식
$b^{-1}(D) \cdot [b^{-1}(D')]_{n - 1} =
b^{-1}(D') \cdot [b^{-1}(D)]_{n - 1}$에서 보조정리가 따른다.

\medskip\noindent
위 증명에서 참조한 각 보조정리는 일반적인 경우($X$가 준콤팩트라고
가정하지 않는 경우)에도 성립함에 유의하자. 일반적인 경우의 유일한 작은
변화는 보조정리 \ref{obsolete-lemma-blowing-up-intersections}를 적용하여 얻는
사상 $b : U' \to U$의 치역이 여집합의 여차원이 $\geq 3$인 열린집합
$U \subset X$라는 것이다. 따라서 『Chow 호몰로지』의 보조정리
\ref{chow-lemma-restrict-to-open}에 의해
$\CH_{n - 2}(U) = \CH_{n - 2}(X)$이고, $X$를 $U$로 바꾸면 증명의
나머지는 변함없이 진행된다.
\end{proof}

\begin{proof}[보조정리 \ref{obsolete-lemma-commutativity-effective-Cartier}의 두 번째 증명]
$\mathcal{I} = \mathcal{O}_X(-D)$와
$\mathcal{I}' = \mathcal{O}_X(-D')$를 각각 $D$와 $D'$의 가역 아이디얼층이라
하자. 다음과 같이 나타낸다.
$\mathcal{I}_{D'} = \mathcal{I} \otimes_{\mathcal{O}_X} \mathcal{O}_{D'}$이고
$\mathcal{I}'_D = \mathcal{I}' \otimes_{\mathcal{O}_X} \mathcal{O}_D$이다.
포함사상 $\mathcal{I} \to \mathcal{O}_X$을 $D'$으로 제한하여 사상
$$
\varphi : \mathcal{I}_{D'} \longrightarrow \mathcal{O}_{D'}
$$
을 얻을 수 있고, 마찬가지로
$$
\psi : \mathcal{I}'_D \longrightarrow \mathcal{O}_D
$$
를 얻는다. 다음은 명백하다.
$$
\Coker(\varphi)
\cong
\mathcal{O}_{D \cap D'}
\cong
\Coker(\psi)
$$
그리고
$$
\Ker(\varphi)
\cong
\frac{\mathcal{I} \cap \mathcal{I}'}{\mathcal{I}\mathcal{I}'}
\cong
\Ker(\psi).
$$
따라서 $K_0(\textit{Coh}_{\leq n - 1}(X))$에서
$$
\gamma =
[\mathcal{I}_{D'}] - [\mathcal{O}_{D'}]
=
[\mathcal{I}'_D] - [\mathcal{O}_D]
$$
임을 알 수 있다. 한편
$$
[\mathcal{I}'_D]_{n - 1} = [D]_{n - 1}, \quad
[\mathcal{I}_{D'}]_{n - 1} = [D']_{n - 1}.
$$
이고
$$
\mathcal{O}_X(D') \otimes \mathcal{I}'_D = \mathcal{O}_D, \quad
\mathcal{O}_X(D) \otimes \mathcal{I}_{D'} = \mathcal{O}_{D'}.
$$
이다. 『Chow 호몰로지』의 보조정리
\ref{chow-lemma-coherent-sheaf-cap-c1}를 두 번 적용하면, 이는 원소
$\gamma$가 $B_{n - 2}(X)$의 원소이고
$c_1(\mathcal{O}_X(D')) \cap [D]_{n - 1}$과
$c_1(\mathcal{O}_X(D)) \cap [D']_{n - 1}$ 둘 다로 사상된다는 뜻이다.
따라서 결론이 나온다(사상
$B_{n - 2}(X) \to \CH_{n - 2}(X)$가 잘 정의된다는 것이 이 증명의
핵심이다).
\end{proof}

\section{정칙 고유 모형 위의 쌍대화 가군}
\label{obsolete-section-dualizing}

\noindent
『반안정 축소』의 상황 \ref{models-situation-regular-model}에서
$\omega_{X/R}^\bullet = f^!\mathcal{O}_{\Spec(R)}$를
$f : X \to \Spec(R)$의 상대 쌍대화 복합체라 하자. 이는
『스킴의 쌍대성』의 주석
\ref{duality-remark-relative-dualizing-complex}에서 도입되었다.
『반안정 축소』의 보조정리 \ref{models-lemma-gorenstein}에 의해
$f$는 상대 차원 $1$인 Gorenstein 사상이므로 『스킴의 쌍대성』의
보조정리
\ref{duality-lemma-affine-flat-Noetherian-gorenstein},
\ref{duality-lemma-CM-shriek},
\ref{duality-lemma-gorenstein-CM-morphism}을 사용하면 어떤 가역
$\mathcal{O}_X$-가군 $\omega_X$에 대해
$$
\omega_{X/R}^\bullet = \omega_X[1]
$$
임을 알 수 있다. 이 가역 가군은 흔히
{\it $R$ 위의 $X$의 상대 쌍대화 가군}이라고 불린다.
$R$은 차원 $1$인 정칙(따라서 Gorenstein) 환이므로
$\omega_R^\bullet = R[1]$은 $R$의 정규화된 쌍대화 복합체이다.
따라서 $\omega_X = H^{-2}(f^!\omega_R^\bullet)$이고,
$\omega_X$가 상대 쌍대화 가군일 뿐 아니라 쌍대화 가군이기도 함을
알 수 있다. 『스킴의 쌍대성』의 예
\ref{duality-example-proper-over-local}를 보라. 그러므로 『스킴의
쌍대성』의 보조정리
\ref{duality-lemma-dualizing-module-proper-over-A}에 의해 $\omega_X$는
함자
$$
\textit{Coh}(\mathcal{O}_X) \to \textit{Sets},\quad
\mathcal{F} \mapsto \Hom_R(H^1(X, \mathcal{F}), R)
$$
를 표현한다. 이는 『반안정 축소』의 상황
\ref{models-situation-regular-model}에서 상대 쌍대화 가군을 정의하는
다른 방법을 준다. $\omega_X$의 형성은 임의의 밑변환과 가환한다
(주어진 상대 차원의 임의의 고유 Gorenstein 사상에 대해).
이는 『스킴의 쌍대성』의 주석
\ref{duality-remark-relative-dualizing-complex}에서 논의한 상대 쌍대화
복합체의 대응하는 사실에서 따르고, 그 사실은 『스킴의 쌍대성』의
보조정리 \ref{duality-lemma-proper-flat-base-change}로 거슬러 올라간다.
따라서 $\omega_X$는 『대수 곡선』의 보조정리
\ref{curves-lemma-duality-dim-1-CM}에서 논의한 $K$ 위의 $C$의 쌍대화
가군 $\omega_C$로 당겨진다. 『대수 곡선』의 보조정리
\ref{curves-lemma-duality-dim-1}에 의해
$\omega_C$는 $\Omega_{C/K}$와 동형임에 유의하자.
마찬가지로 $\omega_X|_{X_k}$는 $k$ 위의 $X_k$의 쌍대화 가군
$\omega_{X_k}$이다.

\begin{lemma}
\label{obsolete-lemma-dualizing-components}
『반안정 축소』의 상황 \ref{models-situation-regular-model}에서
$k$ 위의 $C_i$의 쌍대화 가군은
$$
\omega_{C_i} = \omega_X(C_i)|_{C_i}
$$
이다. 여기서 $\omega_X$는 위와 같다.
\end{lemma}

\begin{proof}
$t : C_i \to X$를 닫힌 몰입이라 하자. $t$는 유효 Cartier 인자의
포함이므로 『스킴의 쌍대성』의 보조정리
\ref{duality-lemma-twisted-inverse-image-closed}와
\ref{duality-lemma-sheaf-with-exact-support-effective-Cartier}에서,
모든 가역 $\mathcal{O}_X$-가군 $\mathcal{L}$에 대해
$t^!(\mathcal{L}) = \mathcal{L}(C_i)|_{C_i}$임을 얻는다.
다음 가환 도표를 생각하자.
$$
\xymatrix{
C_i \ar[r]_t \ar[d]_g & X \ar[d]^f \\
\Spec(k) \ar[r]^s & \Spec(R)
}
$$
$C_i$는 가역 쌍대화 가군 $\omega_{C_i}$를 갖는 Gorenstein 곡선이며
(『반안정 축소』의 보조정리 \ref{models-lemma-gorenstein}),
이 가군은 성질 $\omega_{C_i}[0] = g^!\mathcal{O}_{\Spec(k)}$로
특징지어진다. 『대수 곡선』의 보조정리
\ref{curves-lemma-duality-dim-1}과 그 증명, 그리고 보조정리
\ref{curves-lemma-duality-dim-1-CM}, \ref{curves-lemma-rr}를 보라.
한편 $s^!(R[1]) = k$이고, 따라서
$$
\omega_{C_i}[0] =
g^! s^!(R[1]) = t^!f^!(R[1]) = t^!\omega_X
$$
이다. 위 사실들을 결합하면 보조정리의 명제를 얻는다.
\end{proof}

\section{중복되거나 나뉜 참조}
\label{obsolete-section-duplicates}

\noindent
이 절은 중복된 결과들과, 전에는 여러 사실을 한꺼번에 말했지만 이제는
구성 부분들로 나뉜 참조들을 모아 두는 곳이다.

\begin{lemma}
\label{obsolete-lemma-characterize-epi-mono}
$X$를 위상 공간이라 하자.
\begin{enumerate}
\item 준층의 범주에서 범주론적 전사상(각각 단사상)은 정확히 성분별로
전사(각각 단사)인 준층 사상이다.
\item 층의 범주에서 범주론적 전사상(각각 단사상)은 정확히 전사
(각각 단사)인 층 사상이고, 정확히 모든 줄기에서 전사(각각 단사)인
사상들이다.
\item 집합의 준층의 전사(각각 단사) 사상을 층화하면 전사(각각 단사)가
된다.
\end{enumerate}
\end{lemma}

\begin{proof}
『공간 위의 층』의 보조정리
\ref{sheaves-lemma-characterize-epi-mono}와
\ref{sheaves-lemma-sheafify-epi-mono}를 결합하면 된다.
\end{proof}

\begin{lemma}
\label{obsolete-lemma-directed-colimit-finite-type}
$X$를 스킴이라 하고, $X$가 준콤팩트이고 준분리라고 하자.
$\mathcal{F}$를 준연접 $\mathcal{O}_X$-가군이라 하자.
그러면 $\mathcal{F}$는 그 유한형 준연접 부분가군들의 유향 쌍극한이다.
\end{lemma}

\begin{proof}
이는 『스킴의 성질』의 보조정리
\ref{properties-lemma-quasi-coherent-colimit-finite-type}와 중복된다.
\end{proof}

\begin{lemma}
\label{obsolete-lemma-points-monomorphism}
$S$를 스킴이라 하고 $X$를 $S$ 위의 대수공간이라 하자.
사상 $\{\Spec(k) \to X \text{ 단사 사상}\} \to |X|$은 단사이다.
\end{lemma}

\begin{proof}
이는 『공간의 성질』의 보조정리
\ref{spaces-properties-lemma-points-monomorphism}과 중복된다.
\end{proof}

\begin{theorem}
\label{obsolete-theorem-equivalence-sheaves-point}
$K$가 체이고 $S = \Spec(K)$라 하자.
$\overline{s}$를 $S$의 기하학적 점이라 하자.
$G = \text{Gal}_{\kappa(s)}$를 절대 Galois 군이라 하자.
그러면 범주의 동치
$\Sh(S_\etale) \to G\textit{-Sets}$,
$\mathcal{F} \mapsto \mathcal{F}_{\overline{s}}$가 있다.
\end{theorem}

\begin{proof}
이는 『에탈 코호몰로지』의 정리
\ref{etale-cohomology-theorem-equivalence-sheaves-point}와 중복된다.
\end{proof}

\begin{remark}
\label{obsolete-remark-tangent-spaces}
중복된 태그 때문에 이곳에 도착했다. 실제 내용은 『형식 변형 이론』의
절 \ref{formal-defos-section-tangent-spaces}을 보라.
\end{remark}

\begin{lemma}
\label{obsolete-lemma-locally-ringed-space-direct-summand-free}
$X$를 국소 환 달린 공간이라 하자. 유한 자유 $\mathcal{O}_X$-가군의
직합인자는 유한 국소 자유이다.
\end{lemma}

\begin{proof}
이는 『가군층』의 보조정리
\ref{modules-lemma-direct-summand-of-locally-free-is-locally-free}와
중복된다.
\end{proof}

\begin{lemma}
\label{obsolete-lemma-characterize-injective}
$R$을 환이라 하고 $E$를 $R$-가군이라 하자. 다음은 동치이다.
\begin{enumerate}
\item $E$는 단사 $R$-가군이고,
\item 아이디얼 $I \subset R$와 가군 준동형 $\varphi : I \to E$가
주어지면 $\varphi$를 확장하는 $R$-가군 준동형 $R \to E$가 존재한다.
\end{enumerate}
\end{lemma}

\begin{proof}
이는 Baer 판정법이다. 『단사 대상』의 보조정리
\ref{injectives-lemma-criterion-baer}를 보라.
\end{proof}

\begin{lemma}
\label{obsolete-lemma-periodic-length}
$R$을 국소환이라 하자.
\begin{enumerate}
\item $(M, N, \varphi, \psi)$가 $2$-주기 복합체이고 $M$, $N$의 길이가
유한이면
$e_R(M, N, \varphi, \psi) = \text{length}_R(M) - \text{length}_R(N)$이다.
\item $(M, \varphi, \psi)$가 $(2, 1)$-주기 복합체이고 $M$의 길이가
유한이면 $e_R(M, \varphi, \psi) = 0$이다.
\item $2$-주기 복합체의 짧은 완전열
$$
0 \to (M_1, N_1, \varphi_1, \psi_1)
\to (M_2, N_2, \varphi_2, \psi_2)
\to (M_3, N_3, \varphi_3, \psi_3)
\to 0
$$
이 주어졌다고 하자. 셋 가운데 둘의 코호몰로지 가군이 유한 길이를
가지면 셋째도 그러하며
$$
e_R(M_2, N_2, \varphi_2, \psi_2) =
e_R(M_1, N_1, \varphi_1, \psi_1) +
e_R(M_3, N_3, \varphi_3, \psi_3).
$$
이다.
\end{enumerate}
\end{lemma}

\begin{proof}
이는 『Chow 호몰로지』의 보조정리
\ref{chow-lemma-additivity-periodic-length}와
\ref{chow-lemma-finite-periodic-length}에서 따른다.
\end{proof}


\begin{lemma}
\label{obsolete-lemma-extensions-of-rings}
$A$를 환이라 하고 $I$를 $A$-가군이라 하자.
\begin{enumerate}
\item $I$가 제곱 영 아이디얼인 환의 확대
$0 \to I \to A' \to A \to 0$들의 집합은 표준적으로
$\Ext^1_A(\NL_{A/\mathbf{Z}}, I)$와 전단사이다.
\item 환 준동형 $A \to B$, $B$-가군 $N$, $A$-가군 준동형
$c : I \to N$과 다음 제곱 영 핵을 갖는 환의 확대들이 주어졌다고 하자.
\begin{enumerate}
\item[(a)] $\alpha \in \Ext^1_A(\NL_{A/\mathbf{Z}}, I)$에 대응하는
$0 \to I \to A' \to A \to 0$,
\item[(b)] $\beta \in \Ext^1_B(\NL_{B/\mathbf{Z}}, N)$에 대응하는
$0 \to N \to B' \to B \to 0$.
\end{enumerate}
그러면 『변형 이론』의 등식 (\ref{defos-equation-to-solve})에 들어가는
사상 $A' \to B'$가 존재할 필요충분조건은 $\beta$와 $\alpha$가
$\Ext^1_A(\NL_{A/\mathbf{Z}}, N)$의 같은 원소로 사상되는 것이다.
\end{enumerate}
\end{lemma}

\begin{proof}
이는 『변형 이론』의 보조정리
\ref{defos-lemma-extensions-of-algebras}와
\ref{defos-lemma-extensions-of-algebras-functorial}에서 따른다.
\end{proof}

\begin{lemma}
\label{obsolete-lemma-extensions-of-ringed-spaces}
$(S, \mathcal{O}_S)$를 환 달린 공간이라 하고 $\mathcal{J}$를
$\mathcal{O}_S$-가군이라 하자.
\begin{enumerate}
\item $\mathcal{J}$가 제곱 영 아이디얼인 환의 층의 확대
$0 \to \mathcal{J} \to \mathcal{O}_{S'} \to \mathcal{O}_S \to 0$들의
집합은 표준적으로
$\Ext^1_{\mathcal{O}_S}(\NL_{S/\mathbf{Z}}, \mathcal{J})$와 전단사이다.
\item 환 달린 공간의 사상
$f : (X, \mathcal{O}_X) \to (S, \mathcal{O}_S)$,
$\mathcal{O}_X$-가군 $\mathcal{G}$, $f$-사상
$c : \mathcal{J} \to \mathcal{G}$와 다음 제곱 영 핵을 갖는 환의 층의
확대들이 주어졌다고 하자.
\begin{enumerate}
\item[(a)]
$\alpha \in \Ext^1_{\mathcal{O}_S}(\NL_{S/\mathbf{Z}}, \mathcal{J})$에
대응하는
$0 \to \mathcal{J} \to \mathcal{O}_{S'} \to \mathcal{O}_S \to 0$,
\item[(b)]
$\beta \in \Ext^1_{\mathcal{O}_X}(\NL_{X/\mathbf{Z}}, \mathcal{G})$에
대응하는
$0 \to \mathcal{G} \to \mathcal{O}_{X'} \to \mathcal{O}_X \to 0$.
\end{enumerate}
그러면 『변형 이론』의 등식
(\ref{defos-equation-to-solve-ringed-spaces})에 들어가는 사상
$X' \to S'$이 존재할 필요충분조건은 $\beta$와 $\alpha$가
$\Ext^1_{\mathcal{O}_X}(Lf^*\NL_{S/\mathbf{Z}}, \mathcal{G})$의 같은
원소로 사상되는 것이다.
\end{enumerate}
\end{lemma}

\begin{proof}
이는 『변형 이론』의 보조정리
\ref{defos-lemma-extensions-of-relative-ringed-spaces}와
\ref{defos-lemma-extensions-of-relative-ringed-spaces-functorial}에서 따른다.
\end{proof}

\begin{lemma}
\label{obsolete-lemma-extensions-of-ringed-topoi}
$(\Sh(\mathcal{B}), \mathcal{O}_\mathcal{B})$를 환 달린 토포스라 하고
$\mathcal{J}$를 $\mathcal{O}_\mathcal{B}$-가군이라 하자.
\begin{enumerate}
\item $\mathcal{J}$가 제곱 영 아이디얼인 환의 층의 확대
$0 \to \mathcal{J} \to \mathcal{O}_{\mathcal{B}'} \to
\mathcal{O}_\mathcal{B} \to 0$들의 집합은 표준적으로
$\Ext^1_{\mathcal{O}_\mathcal{B}}(
\NL_{\mathcal{O}_\mathcal{B}/\mathbf{Z}}, \mathcal{J})$와 전단사이다.
\item 환 달린 토포스의 사상
$f : (\Sh(\mathcal{C}), \mathcal{O}) \to
(\Sh(\mathcal{B}), \mathcal{O}_\mathcal{B})$,
$\mathcal{O}$-가군 $\mathcal{G}$,
$f^{-1}\mathcal{O}_\mathcal{B}$-가군 준동형
$c : f^{-1}\mathcal{J} \to \mathcal{G}$와 다음 제곱 영 핵을 갖는 환의
층의 확대들이 주어졌다고 하자.
\begin{enumerate}
\item[(a)]
$\alpha \in \Ext^1_{\mathcal{O}_\mathcal{B}}(
\NL_{\mathcal{O}_\mathcal{B}/\mathbf{Z}}, \mathcal{J})$에 대응하는
$0 \to \mathcal{J} \to \mathcal{O}_{\mathcal{B}'} \to
\mathcal{O}_\mathcal{B} \to 0$,
\item[(b)]
$\beta \in \Ext^1_\mathcal{O}(\NL_{\mathcal{O}/\mathbf{Z}}, \mathcal{G})$에
대응하는
$0 \to \mathcal{G} \to \mathcal{O}' \to \mathcal{O} \to 0$.
\end{enumerate}
그러면 『변형 이론』의 등식
(\ref{defos-equation-to-solve-ringed-topoi})에 들어가는 사상
$(\Sh(\mathcal{C}), \mathcal{O}') \to
(\Sh(\mathcal{B}, \mathcal{O}_{\mathcal{B}'})$가 존재할 필요충분조건은
$\beta$와 $\alpha$가
$\Ext^1_\mathcal{O}(
Lf^*\NL_{\mathcal{O}_\mathcal{B}/\mathbf{Z}}, \mathcal{G})$의 같은 원소로
사상되는 것이다.
\end{enumerate}
\end{lemma}

\begin{proof}
이는 『변형 이론』의 보조정리
\ref{defos-lemma-extensions-of-relative-ringed-topoi}와
\ref{defos-lemma-extensions-of-relative-ringed-topoi-functorial}에서 따른다.
\end{proof}

\begin{remark}
\label{obsolete-remark-examples-formal-defos}
이 태그는 전에는 변형 문제의 여러 예를 기술하는 한 절을 가리켰다.
이제는 각 예마다 독립된 절이 있다. 『변형 문제』의 절
\ref{examples-defos-section-finite-projective-modules},
\ref{examples-defos-section-representations},
\ref{examples-defos-section-continuous-representations},
\ref{examples-defos-section-graded-algebras}를 보라.
\end{remark}

\begin{lemma}
\label{obsolete-lemma-examples-have-RS}
『변형 문제』의 예
\ref{examples-defos-example-finite-projective-modules},
\ref{examples-defos-example-representations},
\ref{examples-defos-example-continuous-representations},
\ref{examples-defos-example-graded-algebras}는 Rim--Schlessinger 조건
(RS)를 만족한다.
\end{lemma}

\begin{proof}
이는 『변형 문제』의 보조정리
\ref{examples-defos-lemma-finite-projective-modules-RS},
\ref{examples-defos-lemma-representations-RS},
\ref{examples-defos-lemma-continuous-representations-RS},
\ref{examples-defos-lemma-graded-algebras-RS}에서 따른다.
\end{proof}

\begin{lemma}
\label{obsolete-lemma-tangent-and-inf}
다음과 같은 표준 $k$-벡터 공간의 동일시들이 있다.
\begin{enumerate}
\item 『변형 문제』의 예
\ref{examples-defos-example-finite-projective-modules}에서
$x_0 = (k, V)$이면 $T_{x_0}\mathcal{F} = (0)$이고
$\text{Inf}_{x_0}(\mathcal{F}) = \text{End}_k(V)$이며, 이들은 유한
차원이다.
\item 『변형 문제』의 예
\ref{examples-defos-example-representations}에서
$x_0 = (k, V, \rho_0)$이면
$T_{x_0}\mathcal{F} = \Ext^1_{k[\Gamma]}(V, V) =
H^1(\Gamma, \text{End}_k(V))$이고
$\text{Inf}_{x_0}(\mathcal{F}) = H^0(\Gamma, \text{End}_k(V))$이다.
$\Gamma$가 유한 생성이면 이들은 유한 차원이다.
\item 『변형 문제』의 예
\ref{examples-defos-example-continuous-representations}에서
$x_0 = (k, V, \rho_0)$이면
$T_{x_0}\mathcal{F} = H^1_{cont}(\Gamma, \text{End}_k(V))$이고
$\text{Inf}_{x_0}(\mathcal{F}) = H^0_{cont}(\Gamma, \text{End}_k(V))$이다.
$\Gamma$가 위상적으로 유한 생성이면 이들은 유한 차원이다.
\item 『변형 문제』의 예
\ref{examples-defos-example-graded-algebras}에서 $x_0 = (k, P)$이면
$T_{x_0}\mathcal{F}$와
$\text{Inf}_{x_0}(\mathcal{F}) = \text{Der}_k(P, P)$는 $P$가 $k$ 위에서
유한 생성일 때 유한 차원이다.
\end{enumerate}
\end{lemma}

\begin{proof}
이는 『변형 문제』의 보조정리
\ref{examples-defos-lemma-finite-projective-modules-TI},
\ref{examples-defos-lemma-representations-TI},
\ref{examples-defos-lemma-continuous-representations-TI},
\ref{examples-defos-lemma-graded-algebras-TI}에서 따른다.
\end{proof}
